\documentclass[12pt, letterpaper]{amsart}

% ======================================================================
% == META AND PACKAGES                                              ==
% ======================================================================
\usepackage[margin=1in]{geometry}
\usepackage{amssymb}
\usepackage{amsfonts}
\usepackage{amsmath}
\usepackage{lmodern}
\usepackage{microtype}
\usepackage{mathtools}
\usepackage[dvipsnames]{xcolor}
\usepackage[
    colorlinks,
    linkcolor=blue,
    citecolor=blue,
    urlcolor=blue,
    implicit=true,
    breaklinks=true
]{hyperref}
\usepackage[noabbrev,capitalize]{cleveref}
\usepackage[T1]{fontenc}
\usepackage[all]{hypcap} % Fixes hyperref pointing to captions instead of figures
\usepackage{enumitem}
\usepackage{appendix}

% ======================================================================
% == CITATION AND REFERENCING                                       ==
% ======================================================================
\bibliographystyle{amsplain}

% ======================================================================
% == CUSTOM DEFINITIONS                                   ==
% ======================================================================
% --- Caligraphic letters ---
\newcommand{\Ecal}{\mathcal{E}}
\newcommand{\Efrak}{\mathfrak{E}}
\newcommand{\Efrakgen}{\Efrak^{F, \Ucal}}
\newcommand{\Thetacal}{\mathcal{A}}
\newcommand{\Ucal}{\mathcal{U}}
\newcommand{\Ucaltheta}{\mathcal{U}_\theta}
\newcommand{\Pcal}{\mathcal{P}}
\newcommand{\Bcal}{\mathcal{B}}
\newcommand{\Hcal}{\mathcal{H}} 
\DeclareMathOperator{\Image}{Im} 

% --- Script letters for function spaces (with arguments) ---
\newcommand{\Scalp}[1]{\mathcal{S}^{#1}}
\newcommand{\Hcalp}[1]{\mathcal{H}^{#1}}
\newcommand{\Lcalp}[1]{\mathcal{L}^{#1}}

% --- Spaces of matrices ---
\newcommand{\Sdpn}{\mathbb{S}_d^{\ge 0}}
\newcommand{\Sdp}{\mathbb{S}_d^{>0}}
\newcommand{\Sd}{\mathbb{S}_d}
\newcommand{\R}{\mathbb{R}}

% --- Mathematical operators and symbols ---
\newcommand{\dd}{\mathrm{d}}
\newcommand{\norm}[1]{\left\lVert#1\right\rVert}
\newcommand{\abs}[1]{\left\lvert#1\right\rvert}
\newcommand{\scpr}[2]{\left\langle #1, #2 \right\rangle}
\newcommand{\Tr}{\mathrm{Tr}}
\DeclareMathOperator*{\dist}{dist}
\DeclareMathOperator*{\argmax}{arg\,max}
\DeclareMathOperator*{\argmin}{arg\,min}
\DeclareMathOperator*{\conv}{conv}
\DeclareMathOperator*{\overlineconv}{\overline{conv}}
\DeclareMathOperator*{\proj}{proj}
\DeclareMathOperator{\rank}{rank}


% --- THEOREM-LIKE ENVIRONMENTS ---
\theoremstyle{plain}
\newtheorem{theorem}{Theorem}[section]
\newtheorem{proposition}[theorem]{Proposition}
\newtheorem{lemma}[theorem]{Lemma}
\newtheorem{corollary}[theorem]{Corollary}
\newtheorem{conjecture}[theorem]{Conjecture}

\theoremstyle{definition}
\newtheorem{definition}[theorem]{Definition}
\newtheorem{assumption}[theorem]{Assumption}
\newtheorem{hypothesis}[theorem]{Hypothesis}

\theoremstyle{remark}
\newtheorem{remark}[theorem]{Remark}
\newtheorem{example}[theorem]{Example}

\numberwithin{equation}{section} % Number equations by section

% ======================================================================
% == DOCUMENT BEGINS                                                ==
% ======================================================================
\begin{document}

\title[A Theory of $\theta$-Expectations]{A Theory of $\theta$-Expectations}

\author{Qian Qi}
\address{Peking University, Beijing 100871, China}
\email{qiqian@pku.edu.cn}
\date{\today}

\begin{abstract}
We address the failure of canonical sub-additive frameworks for stochastic calculus to model non-convex uncertainty. Our proposed \(\theta\)-BSDE formulation leads to a Backward Stochastic Differential Inclusion (BSDI) where the fundamental obstacle to forming a predictive theory is non-uniqueness. We establish well-posedness by identifying a new, weak geometric condition: the model's information structure must be transverse to the singular strata of the uncertainty set's geometry. Using Malliavin calculus, we prove this condition ensures that solution paths almost surely avoid the locus of ambiguous maximizers, thus restoring uniqueness. This result reveals that while the value process itself may depend only on the convex hull of the uncertainty set, the full solution—containing the unique maximizer process \(\Theta_t\)—provides a richer, path-dependent statistic that distinguishes non-convex models, thereby founding a new, tractable stochastic calculus.
\end{abstract}

\maketitle
\tableofcontents

\section{Introduction}

\subsection{The Identifiability Challenge of Convex Models}
A primary mathematical framework for modeling ambiguity, or Knightian uncertainty \cite{Knight1921}, is the theory of sublinear expectations and the associated G-calculus \cite{Peng2019}. This framework provides a dynamically consistent stochastic calculus for settings where probability models are not known with precision. Its mathematical foundation rests on replacing the classical linear expectation with a functional $\Ecal[\cdot]$ that satisfies monotonicity, cash-invariance, and the axiom of sub-additivity:
\[ \Ecal[X+Y] \le \Ecal[X] + \Ecal[Y]. \]
This axiom is not a mere technical convenience. By the Riesz-Kantorovich representation theorem, it is mathematically equivalent to the existence of a dual representation of $\Ecal[\cdot]$ as the supremum of linear expectations over a convex, weakly compact set of probability measures $\Pcal$:
\[ \Ecal[X] = \sup_{P \in \Pcal} \mathbb{E}_P[X]. \]
This equivalence, while mathematically powerful, leads to a fundamental identifiability challenge. The functional $\Ecal[\cdot]$, being the support function of the set $\Pcal$, is completely determined by the closed convex hull of $\Pcal$, denoted $\overlineconv(\Pcal)$ \cite{Rockafellar1970}. Consequently, any two uncertainty sets $\Pcal_1$ and $\Pcal_2$ that share the same convex hull are observationally equivalent. They generate the exact same functional $\Ecal[\cdot]$ and are thus indistinguishable from the perspective of the calculus. The standard theory is, by its very architecture, insensitive to the geometry of non-convex structures.

This insensitivity poses a significant limitation in applications where ambiguity is naturally described by non-convex sets. For instance, uncertainty about the correct Equation of State (EoS) for matter inside a neutron star might be represented by a disjoint union of two distinct classes of models, $\Ucal = \Ucal_{\text{hadronic}} \cup \Ucal_{\text{quark}}$. The classical framework necessitates a convexification step, $\overlineconv(\Ucal)$, which introduces a continuum of averaged physical laws that have no theoretical basis. This procedure discards crucial information about the disjoint nature of the underlying scientific uncertainty.

\subsection{A Constructive Framework and its Fundamental Obstacle}
This paper develops a mathematical theory for valuation under ambiguity that confronts this identifiability problem. Our objective is to construct a rigorous and tractable calculus that remains sensitive to the specific geometry of the underlying uncertainty set $\Ucal$, without recourse to convexification.

Our approach is constructive. We propose a new class of backward stochastic differential equations, which we term \emph{$\theta$-BSDEs}. The solution is a triplet of adapted processes $(Y,Z,\Theta)$. The first component is a backward evolution for the value process $Y$ and the volatility process $Z$:
\begin{equation} \label{eq:intro_bsde_revised}
    Y_t = \xi + \int_t^T F(s, \omega, Y_s, Z_s, \Theta_s) \dd s - \int_t^T Z_s \dd B_s.
\end{equation}
The novelty lies in the second component, which endogenously determines the parameter process $\Theta$ from the primitive, possibly non-convex, uncertainty set $\Ucal$. At each instant in time, the parameter $\Theta_t$ is chosen to solve a pointwise optimization problem that depends on the current state $(Y_t, Z_t)$ of the solution itself:
\begin{equation} \label{eq:intro_optim_revised}
    \Theta_t(\omega) \in \underset{\theta \in \Ucal}{\argmax} \; F(t, \omega, Y_t(\omega), Z_t(\omega), \theta).
\end{equation}
The endogenous feedback in \eqref{eq:intro_optim_revised} renders the system intractable in general. The argmax correspondence is not guaranteed to be single-valued. This transforms the system into a Backward Stochastic Differential Inclusion (BSDI), for which we establish general existence of a solution (\Cref{thm:bsdi_existence}). However, uniqueness fails fundamentally. A single state $(Y_t, Z_t)$ might correspond to a non-singleton set of optimizers, making the choice of $\Theta_t$ ambiguous and undermining the model's identifiability. Restoring uniqueness is the central challenge addressed by this paper.

\subsection{The Main Contribution: Restoring Uniqueness via Structural Regularity}
The central contribution of this paper is the identification and rigorous analysis of two distinct structural conditions that render the $\theta$-BSDE system well-posed.

\paragraph{\textbf{A Benchmark Case: Global Regularity.}}
Our first path to well-posedness, presented as a crucial benchmark, is to impose a strong a priori analytic condition on the problem's primitives. This condition, detailed in the \textbf{Global Regularity Hypothesis (\Cref{hyp:global_regularity})}, demands that the argmax correspondence is a well-defined, single-valued function that exhibits global Lipschitz regularity. While we prove this class of models is non-vacuous (\Cref{app:nonconvex_example_detailed}), the hypothesis is non-generic and effectively assumes away the central geometric difficulty of the problem. Its primary role is to define an ideal class of tractable non-convex models and to motivate the search for a more general theory.

\paragraph{\textbf{The Principal Result: Geometric Regularity.}}
The Global Regularity Hypothesis is often too restrictive, particularly for models involving projections onto non-convex sets. Our main contribution, and the technical heart of the paper, is a new existence and uniqueness theorem (\Cref{thm:existence_geometric}) for a more general class of geometrically regular models that fail this hypothesis. The proof of this result constitutes the paper's principal technical achievement.

The assumptions for this result are strong, but each is indispensable to our proof technique. We require:
\begin{enumerate}[label=(\roman*)]
    \item \textbf{Geometric regularity of $\Ucal$}: The uncertainty set must be a compact semi-algebraic set, ensuring its cut locus—the set of points with ambiguous projections—is a measure-zero set that admits a finite stratification into smooth submanifolds.

    \item \textbf{Probabilistic non-degeneracy of the noise}: The underlying forward process must be driven by uniformly elliptic noise, ensuring its law is sufficiently diffuse to explore the state space.

    \item \textbf{Transversality of the information structure}: For a projection-type driver, the map $g_s$ that determines the target of the projection must be \emph{transverse} to every stratum $\Sigma_i$ of the cut locus. Formally, for every state $v$ such that $g_s(v) \in \Sigma_i$, the tangent spaces must span the ambient space:
    \[ \mathrm{Im}(D g_s(v)) + T_{g_s(v)}\Sigma_i = T_{g_s(v)}\Sd. \]
\end{enumerate}
This transversality condition is a profound and geometrically natural requirement. It does not impose a restrictive dimensional constraint on the model's state space. Instead, it demands that the new information generated by the state process (the range of the Jacobian, $\Image(Dg_s)$) is never pathologically aligned with the directions of ambiguity within the cut locus (the tangent space $T\Sigma_i$). The information flow must always provide a component that points out of the singular set, a condition that embodies a well-posed experimental design.

The proof of uniqueness under these conditions is a delicate Malliavin calculus argument, detailed in the proof of \Cref{thm:existence_geometric}. First, we establish that the law of the full solution process $(s, Y_s, Z_s)$ is absolutely continuous with respect to the Lebesgue measure on the state space. This is a consequence of the uniform ellipticity of the forward noise, which propagates through the coupled FBSDE system. The transversality condition then allows for a powerful application of a probabilistic formulation of Thom's theorem. It guarantees that the target process for the projection almost surely avoids the measure-zero cut locus where the maximizer is ill-defined. This pathwise consistency property ensures that, for any solution, the maximizer process $\Theta_t$ is uniquely determined. This collapses the inclusion to an effective BSDE with a single-valued driver, for which a localization argument proves uniqueness of the solution triplet.

\subsection{Contributions and a Key Insight}
By providing a constructive, well-posed theory under these new, more general conditions, this work achieves several goals. We formulate the general problem as a BSDE and identify the core obstacle to tractability. We then establish existence and uniqueness for two significant, non-trivial classes of non-convex models. By replacing a restrictive algebraic condition (global regularity) with a more fundamental geometric condition (transversality), we significantly broaden the class of problems for which a tractable and identifiable calculus exists. Based on this, we develop a coherent martingale theory and establish a nonlinear Feynman-Kac formula.

Perhaps the most crucial insight of this paper is revealed when analyzing the structure of the resulting valuation. The solution process $Y_t = \Efrakgen_t[\xi]$—the value itself—is structurally insensitive to the non-convexity of $\Ucal$. In the canonical case of volatility ambiguity, it coincides identically with the value derived from the G-calculus, which depends only on $\overlineconv(\Ucal)$.f

The innovation of our framework, and its ability to distinguish non-convex models, lies in the full solution triplet $(Y, Z, \Theta)$. The unique maximizer process $\Theta_t$ acts as a path-dependent indicator of the most adversarial model chosen from the primitive set $\Ucal$. This process is demonstrably sensitive to the specific geometry of $\Ucal$, providing a richer, more informative description of the system's dynamics under ambiguity. The resulting operator defines a dynamically consistent, non-subadditive expectation and yields a tractable theory for a significant new class of non-convex problems.


\section{A General Framework via Stochastic Inclusions}
\label{sec:inclusion_framework}

The theory of $\theta$-expectations is motivated by the desire to model ambiguity structures that are not necessarily convex. The constructive approach proposed in this paper is centered on a backward stochastic differential equation whose driver is determined by a pointwise maximization problem. In its most natural form, this problem does not necessarily admit a unique maximizer. This observation motivates a more general perspective, rooted in the theory of set-valued analysis and stochastic differential inclusions.

This section establishes this general framework. We begin by formalizing the problem as a Backward Stochastic Differential Inclusion (BSDI). We then state and prove a general existence result for solutions to this inclusion, which holds under broad conditions. However, this generality comes at the cost of uniqueness, which typically fails. This failure represents a fundamental challenge to identifiability. The main contributions of this paper, presented in subsequent sections, can then be understood as providing two distinct and powerful sets of conditions under which this non-uniqueness vanishes, leading to a tractable and identifiable theory.

\subsection{Set-Valued Drivers and Backward Stochastic Differential Inclusions}

Let the problem primitives be a compact uncertainty set $\Ucal \subset \Sd$ and a driver function $F: [0,T] \times \Omega \times \R \times \R^{1\times d} \times \Ucal \to \R$. The core of the model is the pointwise optimality condition, which we now formalize as a set-valued map.

\begin{definition}[The Argmax Correspondence]
The \emph{argmax correspondence} $\mathcal{G}: [0,T] \times \Omega \times \R \times \R^{1\times d} \rightrightarrows \Ucal$ is the set-valued map that assigns to each state $(t, \omega, y, z)$ the set of optimizers:
\begin{equation}
    \mathcal{G}(t, \omega, y, z) \coloneqq \underset{\theta \in \Ucal}{\argmax} \; F(t, \omega, y, z, \theta).
\end{equation}
\end{definition}

For the problem to be well-defined, we require conditions ensuring this set is non-empty and that measurable selections exist.

\begin{assumption}\label{ass:bsdi_primitives}
The problem primitives satisfy the following:
\begin{enumerate}[label=(\roman*)]
    \item The uncertainty set $\Ucal$ is a non-empty, compact subset of $\Sd$.
    \item The driver $F(t,\omega,y,z,\theta)$ is a Carathéodory function: for a.e. $(t,\omega)$, the map $(y,z,\theta) \mapsto F(t,\omega,y,z,\theta)$ is continuous; for each $(y,z,\theta)$, the map $(t,\omega) \mapsto F(t,\omega,y,z,\theta)$ is progressively measurable.
    \item $F$ is uniformly Lipschitz in $(y,z)$, i.e., there exists a constant $L_F > 0$ such that for all $\theta \in \Ucal$ and all $(t,\omega,y_1,z_1,y_2,z_2)$:
    \[ |F(t,\omega,y_1,z_1,\theta) - F(t,\omega,y_2,z_2,\theta)| \le L_F(|y_1-y_2|+\norm{z_1-z_2}). \]
    \item The process $s \mapsto \sup_{\theta \in \Ucal} |F(s, \omega, 0, 0, \theta)|$ belongs to $\Lcalp{2}([0,T])$.
\end{enumerate}
\end{assumption}

Under Assumption \ref{ass:bsdi_primitives}(i)-(ii), the Weierstrass extreme value theorem guarantees that $\mathcal{G}(t,\omega,y,z)$ is a non-empty, compact set for all valid arguments. This allows us to define the problem in its most general form.

\begin{definition}[The $\theta$-BSDE as a Backward Stochastic Differential Inclusion]
\label{def:bsdi}
Given a terminal condition $\xi \in L^2(\mathcal{F}_T)$, a solution to the $\theta$-BSDE is a triplet of adapted processes $(Y, Z, \Theta)$ where $(Y,Z) \in \Bcal^2([0,T])$ and $\Theta$ is a progressively measurable process taking values in $\Ucal$ such that:
\begin{enumerate}
    \item \textbf{(Backward Evolution):} For all $t \in [0,T]$, the integral equation holds $P_0$-a.s.:
    \begin{equation}
        Y_t = \xi + \int_t^T F(s, \omega, Y_s, Z_s, \Theta_s) \dd s - \int_t^T Z_s \dd B_s.
    \end{equation}
    \item \textbf{(Inclusion Constraint):} For $P_0\otimes dt$-almost every $(s,\omega)$, the process $\Theta$ satisfies the inclusion:
    \begin{equation}
        \Theta_s(\omega) \in \mathcal{G}(s, \omega, Y_s(\omega), Z_s(\omega)).
    \end{equation}
\end{enumerate}
The process $\Theta$ is called a \emph{measurable selection} of the argmax correspondence. The system itself is a Backward Stochastic Differential Inclusion (BSDI). The core challenge is the endogenous feedback loop: the choice of the selection $\Theta_s$ influences the solution $(Y_s, Z_s)$, which in turn determines the set $\mathcal{G}$ from which $\Theta_s$ must be chosen.
\end{definition}

\subsection{Existence of Solutions in the General Case}
In this general setting, the primary question is the existence of at least one solution. This requires ensuring that a suitable measurable selection $\Theta$ can be found. The theory of measurable multifunctions provides the necessary tools.

\begin{theorem}[General Existence for the $\theta$-BSDI]
\label{thm:bsdi_existence}
Let Assumption \ref{ass:bsdi_primitives} hold. Then for any terminal condition $\xi \in L^2(\mathcal{F}_T)$, there exists at least one solution $(Y,Z,\Theta)$ to the $\theta$-BSDI defined in \Cref{def:bsdi}.
\end{theorem}
\begin{proof}[Proof of Theorem \ref{thm:bsdi_existence}]
The proof strategy is to reduce the set-valued problem to a standard, single-valued BSDE by employing a measurable selection theorem. This approach is standard in the theory of differential inclusions and relies on establishing that the effective driver, defined as the value function of the maximization problem, inherits the necessary regularity. The argument proceeds in four steps.

\vspace{1ex}
\noindent\textbf{Step 1: Properties of the Value Function and Argmax Correspondence.}
Our first objective is to formalize the maximization problem that defines the endogenous parameter choice and to establish the essential analytic properties of its key outputs: the value function and the set of maximizers.

Let us define the value function $V: [0,T] \times \Omega \times \R \times \R^{1\times d} \to \R$ by
\[ V(t,\omega,y,z) \coloneqq \sup_{\theta \in \Ucal} F(t,\omega,y,z,\theta), \]
and the argmax correspondence $\mathcal{G}: [0,T] \times \Omega \times \R \times \R^{1\times d} \rightrightarrows \Ucal$ by
\[ \mathcal{G}(t,\omega,y,z) \coloneqq \underset{\theta \in \Ucal}{\argmax} \; F(t,\omega,y,z,\theta). \]
We shall now prove, for almost every fixed $(t,\omega)$, that $V$ is continuous in $(y,z)$ and that $\mathcal{G}$ is a non-empty, compact-valued, and upper hemi-continuous correspondence. Subsequently, we will establish that $V$ is a Carathéodory function.

\paragraph{\textbf{Continuity and Set-Valued Properties via Berge's Maximum Theorem}}
Let $(t,\omega)$ be a point for which the map $(y,z,\theta) \mapsto F(t,\omega,y,z,\theta)$ is continuous, which holds for almost every $(t,\omega)$ by Assumption \ref{ass:bsdi_primitives}(ii). We verify the hypotheses of Berge's Maximum Theorem \cite[Theorem 17.31]{Aliprantis2006} for the mapping $(y,z) \mapsto V(t,\omega,y,z)$.
\begin{enumerate}[label=(\alph*)]
    \item The function being maximized, $F(t,\omega, \cdot, \cdot, \cdot)$, is continuous on the product space $(\R \times \R^{1\times d}) \times \Ucal$.
    \item The constraint set $\Ucal$ is a non-empty, compact topological space, as stipulated by Assumption \ref{ass:bsdi_primitives}(i).
\end{enumerate}
The conditions of the theorem are satisfied. Therefore, Berge's Maximum Theorem yields two conclusions:
\begin{enumerate}[label=(\roman*)]
    \item The value function $V(t,\omega, \cdot, \cdot)$ is continuous on $\R \times \R^{1\times d}$.
    \item The argmax correspondence $\mathcal{G}(t,\omega, \cdot, \cdot)$ is non-empty and upper hemi-continuous.
\end{enumerate}
Furthermore, because $F(t,\omega,y,z,\cdot)$ is a continuous function on the compact set $\Ucal$, the set of its maximizers, $\mathcal{G}(t,\omega,y,z)$, is a closed subset of $\Ucal$. Since $\Ucal$ is compact, $\mathcal{G}(t,\omega,y,z)$ is also compact. This confirms that the argmax correspondence is non-empty and compact-valued for all $(y,z) \in \R \times \R^{1\times d}$.

\paragraph{\textbf{The Carathéodory Property of the Value Function $V$}}
We have just shown that for a.e. $(t,\omega)$, the map $(y,z) \mapsto V(t,\omega,y,z)$ is continuous. It remains to show that for any fixed $(y,z)$, the map $(t,\omega) \mapsto V(t,\omega,y,z)$ is progressively measurable.

Let $(y,z)$ be fixed. The set $\Ucal \subset \Sd$ is a compact metric space, and is therefore separable. Let $\{\theta_k\}_{k=1}^\infty$ be a countable dense subset of $\Ucal$. Since, for fixed $(t,\omega,y,z)$, the map $\theta \mapsto F(t,\omega,y,z,\theta)$ is continuous on $\Ucal$, the supremum over $\Ucal$ is equal to the supremum over this dense subset:
\[ V(t,\omega,y,z) = \sup_{\theta \in \Ucal} F(t,\omega,y,z,\theta) = \sup_{k \in \mathbb{N}} F(t,\omega,y,z,\theta_k). \]
By Assumption \ref{ass:bsdi_primitives}(ii), for each fixed $k$, the map $(t,\omega) \mapsto F(t,\omega,y,z,\theta_k)$ is progressively measurable. The supremum of a countable collection of progressively measurable functions is itself progressively measurable. Thus, the map $(t,\omega) \mapsto V(t,\omega,y,z)$ is progressively measurable.

Combining these two established properties—(i) continuity in the state variables $(y,z)$ for a.e. fixed $(t,\omega)$ and (ii) progressive measurability in $(t,\omega)$ for fixed $(y,z)$—we conclude that the value function $V$ is a Carathéodory function. This regularity is crucial for the subsequent steps of the proof.


\vspace{1ex}
\noindent\textbf{Step 2: Existence of a Measurable Selection.}
To reduce the inclusion to a standard equation, we must establish the existence of a measurable function that, for each state $(t, \omega, y, z)$, selects a specific maximizer from the set $\mathcal{G}(t, \omega, y, z)$. The existence of such a selector is a classic result from the theory of measurable multifunctions.

Let $\mathcal{X} \coloneqq [0,T] \times \Omega \times \R \times \R^{1\times d}$ be the domain of the argmax correspondence $\mathcal{G}$. We equip the space $[0,T] \times \Omega$ with the progressive $\sigma$-algebra, denoted by $\mathcal{P}$, and the space $\R \times \R^{1\times d}$ with its Borel $\sigma$-algebra $\mathcal{B}(\R \times \R^{1\times d})$. The domain $\mathcal{X}$ is endowed with the product $\sigma$-algebra $\mathfrak{X} \coloneqq \mathcal{P} \otimes \mathcal{B}(\R \times \R^{1\times d})$. The target space $\Ucal$ is a compact metric space, and is therefore a Polish space, equipped with its Borel $\sigma$-algebra $\mathcal{B}(\Ucal)$.

We are now in a position to apply the Aumann-Yankov Measurable Selection Theorem (see, e.g., \cite[Theorem 8.1.3]{CastaingValadier1977}). This theorem guarantees the existence of a measurable selector provided that the graph of the correspondence $\mathcal{G}$, denoted $\mathrm{Gr}(\mathcal{G})$, is a measurable set in the product space $(\mathcal{X} \times \Ucal, \mathfrak{X} \otimes \mathcal{B}(\Ucal))$.

Let us verify that $\mathrm{Gr}(\mathcal{G})$ is indeed measurable. The graph is defined as:
\[ \mathrm{Gr}(\mathcal{G}) = \{ (t, \omega, y, z, \theta) \in \mathcal{X} \times \Ucal \mid F(t,\omega,y,z,\theta) - V(t,\omega,y,z) = 0 \}. \]
We proceed by establishing the measurability of the functions involved.
\begin{enumerate}[label=(\roman*)]
    \item By Assumption \ref{ass:bsdi_primitives}(ii), the function $F$ is a Carathéodory function. A standard result states that a Carathéodory function is jointly measurable with respect to the relevant product $\sigma$-algebra. Specifically, $F$ is $\mathfrak{X} \otimes \mathcal{B}(\Ucal)$-measurable.

    \item In Step 1, we established that the value function $V(t,\omega,y,z) = \sup_{\theta \in \Ucal} F(t,\omega,y,z,\theta)$ is a Carathéodory function on $\mathcal{X}$. Therefore, $V$ is $\mathfrak{X}$-measurable.
\end{enumerate}
We can view $V$ as a function on the larger space $\mathcal{X} \times \Ucal$ that is constant in the $\theta$ variable. As such, it remains $\mathfrak{X} \otimes \mathcal{B}(\Ucal)$-measurable. The function $\Psi: \mathcal{X} \times \Ucal \to \R$ defined by $\Psi(t,\omega,y,z,\theta) \coloneqq F(t,\omega,y,z,\theta) - V(t,\omega,y,z)$ is the difference of two $\mathfrak{X} \otimes \mathcal{B}(\Ucal)$-measurable functions, and is therefore itself measurable.

The graph $\mathrm{Gr}(\mathcal{G})$ is the preimage of the closed set $\{0\}$ under the measurable function $\Psi$:
\[ \mathrm{Gr}(\mathcal{G}) = \Psi^{-1}(\{0\}). \]
It follows immediately that $\mathrm{Gr}(\mathcal{G})$ is an element of $\mathfrak{X} \otimes \mathcal{B}(\Ucal)$.

The hypotheses of the Aumann-Yankov theorem are thus satisfied:
\begin{enumerate}
    \item The correspondence $\mathcal{G}$ is a map from a measurable space $(\mathcal{X}, \mathfrak{X})$ to a Polish space $(\Ucal, \mathcal{B}(\Ucal))$.
    \item For every $x \in \mathcal{X}$, the set $\mathcal{G}(x)$ is non-empty and closed (in fact, compact, as shown in Step 1).
    \item The graph $\mathrm{Gr}(\mathcal{G})$ is a measurable subset of $\mathcal{X} \times \Ucal$.
\end{enumerate}
The theorem then guarantees the existence of a measurable function $\hat{\Theta}: (\mathcal{X}, \mathfrak{X}) \to (\Ucal, \mathcal{B}(\Ucal))$ such that for all $x \in \mathcal{X}$,
\[ \hat{\Theta}(x) \in \mathcal{G}(x). \]
This function $\hat{\Theta}(t,\omega,y,z)$ is the required measurable selector.

\vspace{1ex}
\noindent\textbf{Step 3: Reduction to a Standard BSDE and Verification of Lipschitz Condition.}
The existence of a measurable selector, established in Step 2, allows us to transition from the stochastic differential inclusion to a classical backward stochastic differential equation with a single-valued driver. Let $\hat{\Theta}: [0,T]\times\Omega\times\R\times\R^{1\times d} \to \Ucal$ be a measurable selector as constructed above. We define the \emph{effective driver} $f: [0,T]\times\Omega\times\R\times\R^{1\times d} \to \R$ by
\[ f(t,\omega,y,z) \coloneqq F(t,\omega,y,z,\hat{\Theta}(t,\omega,y,z)). \]
By the definition of the selector $\hat{\Theta}$ and the value function $V$ from Step 1, this effective driver is precisely the value function itself:
\[ f(t,\omega,y,z) = V(t,\omega,y,z) = \sup_{\theta \in \Ucal} F(t,\omega,y,z,\theta). \]
We now prove that $f$ satisfies the standard conditions of the existence and uniqueness theorem for BSDEs from Pardoux and Peng \cite{PardouxPeng1990}. Let $(t,\omega)$ be a point for which the required properties of $F$ hold.

\begin{enumerate}[label=(\alph*)]
    \item \textbf{Progressive Measurability:} For any fixed $(y,z) \in \R \times \R^{1\times d}$, the map $(t,\omega) \mapsto f(t,\omega,y,z)$ is progressively measurable. This follows because $f$ is equal to the value function $V$, and as established in Step 1, $V$ is a Carathéodory function, which implies it is product-measurable and hence progressively measurable.

    \item \textbf{Square-Integrability at Zero:} The process $s \mapsto f(s,\omega,0,0)$ belongs to $\Lcalp{2}([0,T])$. This is a direct consequence of Assumption \ref{ass:bsdi_primitives}(iv), since
    \[ |f(s,\omega,0,0)| = \left|\sup_{\theta \in \Ucal} F(s, \omega, 0, 0, \theta)\right| \le \sup_{\theta \in \Ucal} |F(s, \omega, 0, 0, \theta)|, \]
    and the right-hand side is square-integrable by assumption.

    \item \textbf{Global Lipschitz Continuity in $(y,z)$:} This is the central verification. We must show that there exists a constant $L_f$ such that for all $(y_1,z_1)$ and $(y_2,z_2)$, the inequality $|f(y_1,z_1) - f(y_2,z_2)| \le L_f(|y_1-y_2| + \norm{z_1-z_2})$ holds.
    
    Let $(y_1,z_1)$ and $(y_2,z_2)$ be two arbitrary points in $\R \times \R^{1\times d}$. By the properties established in Step 1, the sets of maximizers $\mathcal{G}(y_1,z_1)$ and $\mathcal{G}(y_2,z_2)$ are non-empty. Let $\theta_1^* \in \mathcal{G}(y_1,z_1)$ and $\theta_2^* \in \mathcal{G}(y_2,z_2)$ be corresponding maximizers.
    
    We first establish an upper bound for the difference $f(y_1,z_1) - f(y_2,z_2)$. By definition, $f(y_1,z_1) = F(y_1,z_1,\theta_1^*)$ and $f(y_2,z_2) = F(y_2,z_2,\theta_2^*)$. The optimality of $\theta_2^*$ for the state $(y_2,z_2)$ implies that $F(y_2,z_2,\theta_2^*) \ge F(y_2,z_2,\theta_1^*)$. Consequently,
    \begin{align*}
        f(y_1,z_1) - f(y_2,z_2) &= F(y_1,z_1,\theta_1^*) - F(y_2,z_2,\theta_2^*) \\
        &\le F(y_1,z_1,\theta_1^*) - F(y_2,z_2,\theta_1^*).
    \end{align*}
    The right-hand side involves the same parameter $\theta_1^*$, allowing us to apply the uniform Lipschitz property of $F$ from Assumption \ref{ass:bsdi_primitives}(iii):
    \begin{equation} \label{eq:lipschitz_upper_bound}
        f(y_1,z_1) - f(y_2,z_2) \le |F(y_1,z_1,\theta_1^*) - F(y_2,z_2,\theta_1^*)| \le L_F(|y_1-y_2| + \norm{z_1-z_2}).
    \end{equation}

    Next, we establish a lower bound. The optimality of $\theta_1^*$ for the state $(y_1,z_1)$ implies that $f(y_1,z_1) = F(y_1,z_1,\theta_1^*) \ge F(y_1,z_1,\theta_2^*)$. This gives
    \begin{align*}
        f(y_1,z_1) - f(y_2,z_2) &\ge F(y_1,z_1,\theta_2^*) - F(y_2,z_2,\theta_2^*) \\
        &\ge -|F(y_1,z_1,\theta_2^*) - F(y_2,z_2,\theta_2^*)|.
    \end{align*}
    Again, applying the uniform Lipschitz property of $F$, this time with the parameter $\theta_2^*$, yields:
    \begin{equation} \label{eq:lipschitz_lower_bound}
         f(y_1,z_1) - f(y_2,z_2) \ge -L_F(|y_1-y_2| + \norm{z_1-z_2}).
    \end{equation}

    Combining the inequalities \eqref{eq:lipschitz_upper_bound} and \eqref{eq:lipschitz_lower_bound}, we conclude that the effective driver $f$ is globally Lipschitz in $(y,z)$ with the same Lipschitz constant $L_F$ as the primitive driver $F$:
    \[ |f(t,\omega,y_1,z_1) - f(t,\omega,y_2,z_2)| \le L_F(|y_1-y_2| + \norm{z_1-z_2}). \]
\end{enumerate}
Since the effective driver $f$ satisfies all the standard conditions, the foundational theorem of Pardoux and Peng \cite{PardouxPeng1990} ensures that the BSDE with driver $f$ and terminal condition $\xi$,
\[ Y_t = \xi + \int_t^T f(s, \omega, Y_s, Z_s) \dd s - \int_t^T Z_s \dd B_s, \]
admits a unique solution pair $(Y,Z)$ in the space $\Bcal^2([0,T])$.

\vspace{1ex}
\noindent\textbf{Step 4: Construction of the Solution Triplet and Verification.}
Let $(Y,Z) \in \Bcal^2([0,T])$ be the unique solution pair to the standard BSDE with driver $f$, whose existence and uniqueness were established in Step 3. We now construct the third component, $\Theta$, and verify that the resulting triplet $(Y,Z,\Theta)$ constitutes a solution to the $\theta$-BSDI as specified in \Cref{def:bsdi}.

\paragraph{\textbf{4a. Construction and Properties of the Selection Process $\Theta$.}}
We define the process $\Theta = (\Theta_t)_{t\in[0,T]}$ by composing the measurable selector $\hat{\Theta}$ from Step 2 with the solution path $(Y,Z)$:
\begin{equation}
    \Theta_s(\omega) \coloneqq \hat{\Theta}(s, \omega, Y_s(\omega), Z_s(\omega)).
\end{equation}
We verify that this process is well-defined and resides in the appropriate function space.
\begin{enumerate}[label=(\roman*)]
    \item \textbf{Measurability:} The solution process $(Y,Z)$ is progressively measurable. Consequently, the mapping $\Psi: [0,T]\times\Omega \to [0,T]\times\Omega\times\R\times\R^{1\times d}$ defined by $\Psi(s,\omega) \coloneqq (s, \omega, Y_s(\omega), Z_s(\omega))$ is also progressively measurable. Since $\hat{\Theta}$ is a product-measurable function and $\Psi$ is progressively measurable, their composition $\Theta = \hat{\Theta} \circ \Psi$ is a progressively measurable process.
    \item \textbf{Value Space and Integrability:} By its construction in Step 2, the selector $\hat{\Theta}$ takes values in the primitive uncertainty set $\Ucal$. Therefore, $\Theta_s(\omega) \in \Ucal$ for all $(s,\omega)$. As $\Ucal$ is compact (Assumption \ref{ass:bsdi_primitives}(i)), it is a bounded subset of $\Sd$. Let $C_\Ucal \coloneqq \sup_{\theta \in \Ucal} \norm{\theta}_F < \infty$. Then $\norm{\Theta_s(\omega)}_F \le C_\Ucal$ for all $(s,\omega)$, which immediately implies that $\Theta \in \Lcalp{p}([0,T];\Sd)$ for all $p \ge 1$. In particular, $\mathbb{E}_{P_0}[\int_0^T \norm{\Theta_s}_F^2 ds] \le C_\Ucal^2 T < \infty$.
\end{enumerate}

\paragraph{\textbf{4b. Verification of the BSDE and Inclusion Constraints.}}
We now confirm that the triplet $(Y,Z,\Theta)$ satisfies the two defining conditions of a solution to the $\theta$-BSDI.
\begin{enumerate}[label=(\roman*)]
    \item \textbf{Backward Evolution Equation:} The pair $(Y,Z)$ is, by construction, the solution to the BSDE with driver $f$. That is, for all $t \in [0,T]$:
    \begin{equation} \label{eq:proof_step4_bsde}
         Y_t = \xi + \int_t^T f(s, \omega, Y_s, Z_s) \dd s - \int_t^T Z_s \dd B_s.
    \end{equation}
    By the definition of $f$ in Step 3 and the definition of $\Theta$ above, we have:
    \[ f(s, \omega, Y_s(\omega), Z_s(\omega)) = F(s, \omega, Y_s(\omega), Z_s(\omega), \hat{\Theta}(s, \omega, Y_s(\omega), Z_s(\omega))) = F(s, \omega, Y_s(\omega), Z_s(\omega), \Theta_s(\omega)). \]
    Substituting this into \eqref{eq:proof_step4_bsde} yields the backward evolution equation required by \Cref{def:bsdi}:
    \[ Y_t = \xi + \int_t^T F(s, \omega, Y_s, Z_s, \Theta_s) \dd s - \int_t^T Z_s \dd B_s. \]
    
    \item \textbf{Inclusion Constraint:} The measurable selector $\hat{\Theta}$ was constructed precisely to satisfy the property $\hat{\Theta}(t,\omega,y,z) \in \mathcal{G}(t,\omega,y,z)$ for all arguments $(t,\omega,y,z)$. Evaluating this property along the solution path gives:
    \[ \Theta_s(\omega) = \hat{\Theta}(s, \omega, Y_s(\omega), Z_s(\omega)) \in \mathcal{G}(s, \omega, Y_s(\omega), Z_s(\omega)). \]
    This inclusion holds for $P_0\otimes dt$-almost every $(s,\omega)$, satisfying the second condition of \Cref{def:bsdi}.
\end{enumerate}
The constructed triplet $(Y,Z,\Theta)$ has components in the required spaces and satisfies both defining conditions of the $\theta$-BSDI. This concludes the proof of the existence of at least one solution.
\end{proof}
\subsection{The Impasse of Non-Uniqueness}
The existence guaranteed by \Cref{thm:bsdi_existence} is not accompanied by uniqueness. In general, there may be multiple measurable selections $\hat{\Theta}$ from the argmax correspondence $\mathcal{G}$, leading to a set of distinct drivers and therefore different solution processes $(Y,Z,\Theta)$.

\begin{example}[Fundamental Non-Uniqueness]
\label{ex:non_uniqueness}
Let the dimension be $d=1$. Let the driver be $F(z,\theta) = z\theta$ where $z, \theta \in \R$, and the uncertainty set be the discrete, non-convex set $\Ucal = \{-1, 1\}$. The argmax correspondence is $\mathcal{G}(z) = \{\text{sgn}(z)\}$ for $z \neq 0$, and $\mathcal{G}(0) = \{-1,1\}$.
Consider the $\theta$-BSDE with terminal condition $\xi=0$:
\[ Y_t = \int_t^T Z_s \Theta_s ds - \int_t^T Z_s dB_s, \quad \Theta_s \in \mathcal{G}(Z_s). \]
The trivial solution $(Y,Z) \equiv (0,0)$ is valid. In this case, $Z_t=0$ for all $t$. The inclusion constraint becomes $\Theta_t \in \mathcal{G}(0) = \{-1,1\}$ for all $t$. Any progressively measurable process $\Theta_s$ taking values in $\{-1,1\}$ forms a valid solution triplet $(0,0,\Theta)$. This multiplicity of valid driving processes $\Theta$ for a single state $(Y,Z)=(0,0)$ undermines the identifiability of the model.
\end{example}

This non-uniqueness constitutes a deeper identifiability impasse. The model does not specify a single value process, but rather a set of possible ones. A tractable and predictive theory requires further structure. The remainder of this paper is dedicated to providing such structure.

\section*{Unifying the Regularity Conditions: A Geometric Hierarchy}
\label{sec:unifying_conditions}

A central contribution of this paper is the identification of two distinct sets of conditions that ensure the well-posedness of the \(\theta\)-BSDE system: the Global Regularity Hypothesis (\Cref{hyp:global_regularity}) and the Geometric Regularity conditions (\Cref{ass:pathwise_consistency}). The referee has aptly noted that these two frameworks may appear disconnected, creating the impression of an ad-hoc collection of sufficient conditions rather than a single, cohesive theory.

The purpose of this section is to address this crucial point by establishing a formal and intuitive hierarchy between these two sets of assumptions. We will prove that Global Regularity is a specific, and significantly stronger, instance of the more general Geometric Regularity framework. This is achieved by demonstrating that any model satisfying the former necessarily satisfies the latter. The key insight is that Global Regularity constitutes a "zeroth-order avoidance" condition, where the problem is regularized because the locus of ambiguity is never encountered. In contrast, Geometric Regularity is a weaker "first-order non-tangency" condition, where ambiguity may be encountered but is immediately resolved by a transversal flow of information.

To make this connection precise, we consider models that fit the projection-type driver structure analyzed in \Cref{sec:main_result_geometric}, as this is the natural setting for the geometric conditions.

\begin{proposition}[Global Regularity Implies Geometric Regularity]
\label{prop:gr_implies_geor}
Consider a Markovian model with primitives satisfying the structural requirements of \Cref{ass:pathwise_consistency} (i.e., a projection-type driver, smooth coefficients, etc.). If this model also satisfies the Global Regularity Hypothesis (\Cref{hyp:global_regularity}), then it necessarily satisfies the Transversality Condition (\Cref{ass:pathwise_consistency}(iv)).
\end{proposition}

\begin{proof}
The proof proceeds by demonstrating that the strong assumptions of Global Regularity impose such a severe restriction on the geometry of the problem that the Transversality Condition is satisfied vacuously.

\vspace{1ex}
\noindent\textbf{Step 1: Unpacking the Implication of Global Regularity.}
The driver has the projection-type form:
\[ F(t,x,y,z,\theta) = h(t,x,y,z) - \frac{1}{2}\norm{\theta - G(t,x,y,z)}_F^2. \]
Maximizing this expression with respect to \(\theta \in \Ucal\) is equivalent to minimizing the squared distance \(\norm{\theta - W}_F^2\), where \(W \coloneqq G(t,x,y,z)\) is the target point. The set of maximizers is therefore the nearest-point projection of \(W\) onto \(\Ucal\):
\[ \underset{\theta \in \Ucal}{\argmax} \; F(t,x,y,z,\theta) = \proj_{\Ucal}(G(t,x,y,z)). \]
The core of the Global Regularity Hypothesis (\Cref{hyp:global_regularity}(iii-a)) is the assertion that this argmax correspondence is a single-valued map \(\Thetacal^*\) for \emph{all} states \((t,\omega,y,z)\).

\vspace{1ex}
\noindent\textbf{Step 2: Connecting Uniqueness to the Cut Locus.}
Let \(\mathrm{Cut}(\Ucal)\) denote the cut locus of the compact set \(\Ucal\). By definition, the cut locus is the set of all points in the ambient space \(\Sd\) that have more than one nearest point in \(\Ucal\). That is,
\[ P \in \mathrm{Cut}(\Ucal) \iff \abs{\proj_{\Ucal}(P)} > 1. \]
The statement that the maximizer is unique for all states implies that for any target point \(W\) that can possibly be generated by the map \(G\), the projection of \(W\) onto \(\Ucal\) must be a singleton. This leads to a crucial geometric conclusion.

\vspace{1ex}
\noindent\textbf{Step 3: The Geometric Implication of "Zeroth-Order Avoidance".}
Let \(\Image(G)\) be the image of the target map \(G: [0,T]\times\R^k\times\R\times\R^{1\times d} \to \Sd\). The requirement of a globally unique maximizer is equivalent to the statement that for every point \(W \in \Image(G)\), we have \(\abs{\proj_{\Ucal}(W)} = 1\).
This is mathematically equivalent to the statement that the image of the target map \(G\) is entirely disjoint from the cut locus of the uncertainty set:
\begin{equation} \label{eq:proof_disjointness}
    \Image(G) \cap \mathrm{Cut}(\Ucal) = \emptyset.
\end{equation}
This is a powerful "zeroth-order avoidance" condition: the problem is structured so that the target point for the projection is \emph{never} in a position of ambiguity.

\vspace{1ex}
\noindent\textbf{Step 4: The Transversality Condition Holds Vacuously.}
Now, let us recall the Transversality Condition (\Cref{ass:pathwise_consistency}(iv)). Let \(\{\Sigma_i\}\) be the finite stratification of \(\mathrm{Cut}(\Ucal)\) into smooth submanifolds. Let \(g_s: \R^{k+1+d} \to \Sd\) be the map \(v \mapsto G(s,v)\). The condition states:
\begin{quote}
    For each \(s \in (0,T]\) and for every stratum \(\Sigma_i\), the map \(g_s\) is transverse to \(\Sigma_i\). This means that for every \(v \in \R^{k+1+d}\) \textbf{such that} \(g_s(v) \in \Sigma_i\), the tangent spaces must span the ambient space:
    \[ \Image(D g_s(v)) + T_{g_s(v)}\Sigma_i = T_{g_s(v)}\Sd. \]
\end{quote}
Consider the premise of this conditional statement: "for every \(v\) such that \(g_s(v) \in \Sigma_i\)". From our finding in \eqref{eq:proof_disjointness}, the set of points \(v\) for which \(g_s(v) \in \Sigma_i \subset \mathrm{Cut}(\Ucal)\) is the empty set.
The condition is therefore a statement about all members of an empty set. In formal logic, any such universal quantification is vacuously true. There are no points \(v\) for which the condition must be checked, so the condition holds trivially.

Since this holds for every stratum \(\Sigma_i\) and every \(s \in (0,T]\), the Transversality Condition is satisfied. This completes the proof.
\end{proof}

\begin{remark}[A Hierarchy of Regularity]
\label{rem:hierarchy}
Proposition \ref{prop:gr_implies_geor} establishes a clear and intuitive hierarchy between our two frameworks. This hierarchy can be understood in terms of how each framework handles the geometric singularities represented by the cut locus.

\begin{enumerate}
    \item \textbf{Global Regularity (Zeroth-Order Condition):} This is the strongest form of regularity. It requires that the problem is structured such that the states leading to ambiguity are never generated. Geometrically, \(\Image(G) \cap \mathrm{Cut}(\Ucal) = \emptyset\). The problem is regular because ambiguity is systematically avoided.

    \item \textbf{Geometric Regularity (First-Order Condition):} This is a significantly weaker and more general requirement. It allows the target process to intersect the cut locus (\(\Image(G) \cap \mathrm{Cut}(\Ucal) \neq \emptyset\)). However, it imposes a condition on the *nature* of this intersection: it must be transverse. The Malliavin calculus argument at the heart of \Cref{thm:existence_geometric} then shows that while the set of singular target points is non-empty, the probability of the solution process generating a target point that lies within this set at any given time is zero. The problem is regularized not by avoiding the singular set entirely, but by ensuring that the system's dynamics are rich enough to almost surely steer clear of it.
\end{enumerate}
This unified perspective demonstrates that our theory is not a disparate collection of results but a cohesive framework built on a single, powerful geometric principle. The main result of this paper, \Cref{thm:existence_geometric}, provides a substantive generalization from a restrictive avoidance condition to a more fundamental and broadly applicable non-tangency condition.
\end{remark}

\section{On the Global Regularity}
\label{sec:wellposedness_global_regularity}

The general framework of Backward Stochastic Differential Inclusions (BSDIs) established in \Cref{sec:inclusion_framework} guarantees the existence of solutions. However, as demonstrated in Example \ref{ex:non_uniqueness}, it fails to ensure uniqueness. This failure is fundamental, arising whenever the argmax correspondence $\mathcal{G}(t,\omega,y,z)$ is not single-valued for a set of states $(t,\omega,y,z)$ that can be visited with positive probability. A predictive and identifiable theory requires a mechanism to restore uniqueness.

This section presents a first, direct resolution to this challenge. We introduce a set of strong analytic conditions on the problem's primitives that enforce the uniqueness and regularity of the maximizer a priori. While these conditions are restrictive and non-generic, they define an important benchmark class of tractable non-convex models. More importantly, their severity motivates the search for the more general geometric framework developed in \Cref{sec:main_result_geometric}, which constitutes the principal contribution of this work.

We begin by formally stating the coupled system of equations at the heart of our theory.

\begin{definition}[The $\theta$-BSDE System]
\label{def:theta_bsde_system}
Let $\Ucal$ be a non-empty, compact subset of $\Sd$, representing the primitive uncertainty set. Let $\xi \in L^2(\Omega, \mathcal{F}_T, P_0)$ be a terminal condition and let $F: [0,T] \times \Omega \times \R \times \R^{1\times d} \times \Ucal \to \R$ be a driver function. A solution to the \emph{$\theta$-BSDE} is a triplet of processes $(Y, Z, \Theta)$ in the space $\Bcal^2([0,T]) \times \Lcalp{2}([0,T]; \Sd)$ that satisfies the following coupled system $P_0$-a.s.:
\begin{enumerate}
    \item \textbf{(Backward Evolution):} For all $t \in [0,T]$,
    \begin{equation} \label{eq:theta_bsde_backward}
        Y_t = \xi + \int_t^T F(s, \omega, Y_s, Z_s, \Theta_s) \dd s - \int_t^T Z_s \dd B_s.
    \end{equation}
    \item \textbf{(Pointwise Optimality):} For $P_0 \otimes dt$-almost every $(t,\omega)$,
    \begin{equation} \label{eq:theta_bsde_optimality}
        \Theta_t(\omega) \in \underset{\theta \in \Ucal}{\argmax} \; F(t, \omega, Y_t(\omega), Z_t(\omega), \theta).
    \end{equation}
\end{enumerate}
\end{definition}

\subsection{Well-Posedness via an A Priori Regularity Hypothesis}

The most direct path to well-posedness is to impose conditions that guarantee the argmax correspondence is, in fact, a well-behaved single-valued function.

\begin{hypothesis}[Global Regularity of the Argmax Correspondence]\label{hyp:global_regularity}
The driver $F$ and the compact set $\Ucal \subset \Sd$ are such that:
\begin{enumerate}[label=(\roman*)]
    \item \textbf{(Standard Primitives)} The function $F$ and set $\Ucal$ satisfy the baseline conditions of Assumption \ref{ass:bsdi_primitives}.
    
    \item \textbf{(Uniform Lipschitz Continuity in Parameter)} There exists a constant $L_{F,\theta} > 0$ such that for all $(t,\omega,y,z)$ and for all $\theta_1, \theta_2 \in \Ucal$,
    \[ |F(t,\omega,y,z,\theta_1) - F(t,\omega,y,z,\theta_2)| \le L_{F,\theta} \norm{\theta_1 - \theta_2}_F. \]

    \item \textbf{(Existence of a Unique, Regular Maximizer Map)} There exists a map $\Thetacal^*: [0,T]\times\Omega\times\R\times\R^{1\times d} \to \Ucal$ satisfying:
        \begin{enumerate}[label=(\alph*)]
            \item \textbf{Unique Maximizer:} For $P_0\otimes dt$-a.e. $(t,\omega)$ and for any $(y,z) \in \R\times\R^{1\times d}$, the set of maximizers is a singleton:
            \[ \underset{\theta \in \Ucal}{\argmax} \; F(t, \omega, y, z, \theta) = \{\Thetacal^*(t,\omega,y,z)\}. \]
            \item \textbf{Global Lipschitz Continuity:} The map $(y,z) \mapsto \Thetacal^*(t,\omega,y,z)$ is Lipschitz continuous with a constant $L_\theta > 0$, uniformly in $(t,\omega)$.
            \item \textbf{Progressive Measurability:} For any fixed $(y,z)$, the process $(t,\omega) \mapsto \Thetacal^*(t,\omega,y,z)$ is progressively measurable.
        \end{enumerate}
\end{enumerate}
\end{hypothesis}

\begin{remark}[On the Restrictiveness of the Global Regularity Hypothesis]
\label{rem:restrictiveness}
Hypothesis \ref{hyp:global_regularity}(iii) is an exceptionally strong condition that is not satisfied by many natural models of non-convex uncertainty. For instance, if the optimization in \eqref{eq:theta_bsde_optimality} corresponds to a nearest-point projection onto a non-convex set $\Ucal$, the maximizer map $\Thetacal^*$ is typically multi-valued and discontinuous on the cut locus (or medial axis) of $\Ucal$. This hypothesis therefore \emph{assumes away the central geometric difficulty} of the problem.

Despite its restrictive, non-generic nature (see \Cref{app:geometric_example_corrected} for a constructed example), this hypothesis serves a crucial purpose. It provides a clear theoretical benchmark, demonstrating that well-posedness is achievable if the interaction between the driver $F$ and the set $\Ucal$ is sufficiently regular. Its primary role in this manuscript is to motivate the search for a more general and geometrically-grounded theory, a task undertaken and resolved in \Cref{sec:main_result_geometric} under a different set of assumptions that do not require a globally regular maximizer map.
\end{remark}

Under this strong hypothesis, the $\theta$-BSDE is equivalent to a standard BSDE, for which classical results apply.

\begin{theorem}[Existence and Uniqueness under Global Regularity]
\label{thm:theta_bsde_wellposedness}
Under Hypothesis \ref{hyp:global_regularity}, for any terminal condition $\xi \in L^2(\Omega, \mathcal{F}_T, P_0)$, the $\theta$-BSDE system (\cref{def:theta_bsde_system}) admits a unique solution $(Y, Z, \Theta)$ in the space $\Bcal^2([0,T]) \times \Lcalp{2}([0,T]; \Sd)$.
\end{theorem}

\begin{proof}
The strategy is to reduce the coupled system to an equivalent, standard BSDE and then invoke the foundational existence and uniqueness result of Pardoux and Peng \cite{PardouxPeng1990}.

\vspace{1ex}
\noindent\textbf{Step 1: Reduction to an Equivalent Standard BSDE.}

The core of our strategy is to demonstrate that Hypothesis \ref{hyp:global_regularity} allows for a reduction of the coupled $\theta$-BSDE system to an equivalent, but more classical, single-driver BSDE. The uniqueness of the maximizer collapses the set-valued nature of the problem into a conventional functional form.

First, we define the \emph{effective driver} $\tilde{F}: [0,T]\times\Omega\times\R\times\R^{1\times d} \to \R$ as the value function of the pointwise optimization problem:
\begin{equation}
    \tilde{F}(t,\omega,y,z) \coloneqq \sup_{\theta \in \Ucal} F(t,\omega,y,z,\theta).
\end{equation}
By Hypothesis \ref{hyp:global_regularity}(iii-a), for $P_0\otimes dt$-a.e. $(t,\omega)$, this supremum is attained at the unique maximizer $\Thetacal^*(t,\omega,y,z)$. This gives the crucial representation:
\begin{equation}
    \tilde{F}(t,\omega,y,z) = F\big(t, \omega, y, z, \Thetacal^*(t,\omega,y,z)\big).
\end{equation}

We now formally establish the equivalence. A triplet $(Y,Z,\Theta)$ is a solution to the $\theta$-BSDE (\cref{def:theta_bsde_system}) if and only if:
\begin{enumerate}[label=(\alph*)]
    \item The pair $(Y,Z)$ solves the standard BSDE with driver $\tilde{F}$:
    \begin{equation} \label{eq:proof_effective_bsde_step1}
        Y_t = \xi + \int_t^T \tilde{F}(s, \omega, Y_s, Z_s) \dd s - \int_t^T Z_s \dd B_s.
    \end{equation}
    \item The process $\Theta$ is determined as a function of the solution $(Y,Z)$:
    \begin{equation} \label{eq:theta_as_function}
         \Theta_t(\omega) = \Thetacal^*(t, \omega, Y_t(\omega), Z_t(\omega)), \quad P_0 \otimes dt\text{-a.e.}
    \end{equation}
\end{enumerate}

\begin{proof}[Proof of Equivalence]
($\Rightarrow$) Assume $(Y,Z,\Theta)$ is a solution to the $\theta$-BSDE system. The pointwise optimality condition \eqref{eq:theta_bsde_optimality}, combined with the uniqueness of the maximizer from Hypothesis \ref{hyp:global_regularity}(iii-a), forces the identity \eqref{eq:theta_as_function} to hold. Substituting this identity into the backward evolution equation \eqref{eq:theta_bsde_backward} gives:
\begin{align*}
    Y_t &= \xi + \int_t^T F\big(s, \omega, Y_s, Z_s, \Thetacal^*(s, \omega, Y_s, Z_s)\big) \dd s - \int_t^T Z_s \dd B_s \\
        &= \xi + \int_t^T \tilde{F}(s, \omega, Y_s, Z_s) \dd s - \int_t^T Z_s \dd B_s.
\end{align*}
This is precisely the standard BSDE \eqref{eq:proof_effective_bsde_step1}. Thus, (a) and (b) hold.

($\Leftarrow$) Conversely, assume that $(Y,Z)$ is a solution to the standard BSDE \eqref{eq:proof_effective_bsde_step1} and that $\Theta$ is defined by \eqref{eq:theta_as_function}. We must verify that the triplet $(Y,Z,\Theta)$ satisfies the two conditions of Definition \ref{def:theta_bsde_system}.
\begin{enumerate}
    \item By reversing the substitution made above, the integral equation \eqref{eq:proof_effective_bsde_step1} for $(Y,Z)$ is immediately equivalent to the backward evolution equation \eqref{eq:theta_bsde_backward} for $(Y,Z,\Theta)$.
    \item By definition of the map $\Thetacal^*$ and the process $\Theta$, we have that for a.e. $(t,\omega)$, $\Theta_t(\omega)$ is the unique element in the argmax set $\underset{\theta \in \Ucal}{\argmax} \; F(t, \omega, Y_t(\omega), Z_t(\omega))$. This directly satisfies the pointwise optimality condition \eqref{eq:theta_bsde_optimality}.
\end{enumerate}
This completes the proof of the equivalence.
\end{proof}

This equivalence reduces the problem of existence and uniqueness for the $\theta$-BSDE to that for the standard BSDE \eqref{eq:proof_effective_bsde_step1}. The subsequent steps of the proof are therefore dedicated to verifying that the effective driver $\tilde{F}$ possesses the requisite properties for the classical theory to apply.

\vspace{1ex}
\noindent\textbf{Step 2: Verification of the Properties of the Effective Driver $\tilde{F}$.}
We verify that the effective driver $\tilde{F}(t,\omega,y,z) = F(t, \omega, y, z, \Thetacal^*(t,\omega,y,z))$ satisfies the canonical conditions required by the existence and uniqueness theory for BSDEs \cite{PardouxPeng1990}.

\begin{enumerate}[label=(\alph*)]
    \item \textbf{Progressive Measurability:} We must show that for any fixed $(y,z) \in \R \times \R^{1\times d}$, the process $(t,\omega) \mapsto \tilde{F}(t,\omega,y,z)$ is progressively measurable. Consider the map $\Phi_{y,z}: [0,T]\times\Omega \to [0,T]\times\Omega\times\R\times\R^{1\times d}\times\Ucal$ defined by
    \[ \Phi_{y,z}(t,\omega) \coloneqq (t, \omega, y, z, \Thetacal^*(t,\omega,y,z)). \]
    By Hypothesis \ref{hyp:global_regularity}(iii-c), the map $(t,\omega) \mapsto \Thetacal^*(t,\omega,y,z)$ is progressively measurable. It follows that $\Phi_{y,z}$ is a progressively measurable mapping. The primitive driver $F$ is a Carathéodory function (Hypothesis \ref{hyp:global_regularity}(i)), which implies it is product-measurable on its domain. The process $\tilde{F}(t,\omega,y,z) = F(\Phi_{y,z}(t,\omega))$ is therefore the composition of a measurable function with a progressively measurable process, and is thus itself progressively measurable.

    \item \textbf{Square-Integrability at Zero:} The process $s \mapsto \tilde{F}(s,\omega,0,0)$ must belong to $\Lcalp{2}([0,T])$. This is a direct consequence of Hypothesis \ref{hyp:global_regularity}(i), which assumes that $s \mapsto \sup_{\theta \in \Ucal} |F(s, \omega, 0, 0, \theta)|$ is in $\Lcalp{2}([0,T])$. Indeed,
    \[ |\tilde{F}(s,\omega,0,0)| = |F(s, \omega, 0, 0, \Thetacal^*(s,\omega,0,0))| \le \sup_{\theta \in \Ucal} |F(s, \omega, 0, 0, \theta)|, \]
    and the claim follows by taking the $L^2$-norm over $[0,T]$.

    \item \textbf{Global Lipschitz Continuity in $(y,z)$:} This is the central verification. We show that $\tilde{F}$ is globally Lipschitz in $(y,z)$ with a constant uniform in $(t,\omega)$. Let $(t,\omega)$ be fixed, and let $(y_1,z_1)$ and $(y_2,z_2)$ be two arbitrary points in $\R \times \R^{1\times d}$. For brevity, we let $\theta_1^* \coloneqq \Thetacal^*(t,\omega,y_1,z_1)$ and $\theta_2^* \coloneqq \Thetacal^*(t,\omega,y_2,z_2)$. The derivation proceeds as follows:
    \begin{align*}
        \abs{\tilde{F}(y_1,z_1) - \tilde{F}(y_2,z_2)} &= \abs{F(y_1,z_1,\theta_1^*) - F(y_2,z_2,\theta_2^*)} \\
        &\le \abs{F(y_1,z_1,\theta_1^*) - F(y_2,z_2,\theta_1^*)} + \abs{F(y_2,z_2,\theta_1^*) - F(y_2,z_2,\theta_2^*)} \\
        &\le L_F(\abs{y_1-y_2} + \norm{z_1-z_2}) + L_{F,\theta} \norm{\theta_1^* - \theta_2^*}_F \\
        &\le L_F(\abs{y_1-y_2} + \norm{z_1-z_2}) + L_{F,\theta} L_\theta (\abs{y_1-y_2} + \norm{z_1-z_2}) \\
        &= (L_F + L_{F,\theta}L_\theta) (\abs{y_1 - y_2} + \norm{z_1 - z_2}).
    \end{align*}
    The first inequality is the triangle inequality. The second inequality follows by applying the uniform Lipschitz property of $F$ with respect to $(y,z)$ (Hypothesis \ref{hyp:global_regularity}(i), with constant $L_F$) to the first term, and the uniform Lipschitz property of $F$ with respect to $\theta$ (Hypothesis \ref{hyp:global_regularity}(ii), with constant $L_{F,\theta}$) to the second term. The third inequality uses the global Lipschitz property of the maximizer map $\Thetacal^*$ (Hypothesis \ref{hyp:global_regularity}(iii-b), with constant $L_\theta$).

    Thus, $\tilde{F}$ is globally Lipschitz in $(y,z)$ with constant $L_{\tilde{F}} \coloneqq L_F + L_{F,\theta}L_\theta$.
\end{enumerate}
All canonical conditions are satisfied, allowing the direct application of standard BSDE theory.

\vspace{1ex}
\noindent\textbf{Step 3: Existence and Uniqueness of the Solution Triplet.}
The verifications in Step 2 establish that the effective driver $\tilde{F}$ satisfies the hypotheses of the foundational existence and uniqueness theorem for BSDEs. Specifically, $\tilde{F}$ is progressively measurable, satisfies the integrability condition $\mathbb{E}[\int_0^T |\tilde{F}(t,0,0)|^2 dt] < \infty$, and is globally Lipschitz continuous in the state variables $(y,z)$.

\vspace{1ex}
\noindent\textit{Existence.} By the main result of Pardoux and Peng \cite{PardouxPeng1990}, the standard BSDE with driver $\tilde{F}$,
\begin{equation*}
    Y_t = \xi + \int_t^T \tilde{F}(s, \omega, Y_s, Z_s) \dd s - \int_t^T Z_s \dd B_s,
\end{equation*}
admits a unique solution pair $(Y,Z)$ in the space $\Bcal^2([0,T])$. We use this unique pair to construct the full solution triplet. Let us define the process $\Theta$ by composing the unique maximizer map $\Thetacal^*$ from Hypothesis \ref{hyp:global_regularity}(iii) with this solution path:
\begin{equation}\label{eq:proof_theta_construction}
    \Theta_t(\omega) \coloneqq \Thetacal^*(t, \omega, Y_t(\omega), Z_t(\omega)), \quad \text{for } P_0 \otimes dt\text{-a.e. } (t,\omega).
\end{equation}
We verify that the resulting triplet $(Y,Z,\Theta)$ is a valid solution to the $\theta$-BSDE system (\cref{def:theta_bsde_system}).
\begin{enumerate}[label=(\alph*)]
    \item The processes $(Y,Z)$ belong to $\Bcal^2([0,T])$ by construction. The process $\Theta$ is progressively measurable, being a composition of the progressively measurable map $(t,\omega) \mapsto (t,\omega,Y_t(\omega),Z_t(\omega))$ and the measurable function $\Thetacal^*$. Furthermore, since $\Thetacal^*$ maps into the compact set $\Ucal$, the process $\Theta$ is uniformly bounded and thus belongs to $\Lcalp{2}([0,T]; \Sd)$. The triplet resides in the correct solution space.
    \item The backward evolution equation \eqref{eq:theta_bsde_backward} is satisfied. By the definition of $\Theta$ in \eqref{eq:proof_theta_construction} and the definition of the effective driver $\tilde{F}$, we have:
    \[ F(t,\omega,Y_t,Z_t,\Theta_t) = F(t,\omega,Y_t,Z_t,\Thetacal^*(t,\omega,Y_t,Z_t)) = \tilde{F}(t,\omega,Y_t,Z_t). \]
    Since $(Y,Z)$ solves the BSDE with driver $\tilde{F}$, the equation holds.
    \item The pointwise optimality condition \eqref{eq:theta_bsde_optimality} is satisfied by the very definition of the map $\Thetacal^*$ as the unique maximizer (Hypothesis \ref{hyp:global_regularity}(iii-a)).
\end{enumerate}
This establishes the existence of at least one solution to the $\theta$-BSDE system.

\vspace{1ex}
\noindent\textit{Uniqueness.} To prove uniqueness, let $(\hat{Y}, \hat{Z}, \hat{\Theta})$ be any arbitrary solution to the $\theta$-BSDE system (\cref{def:theta_bsde_system}). We will show that it must coincide with the triplet $(Y,Z,\Theta)$ constructed above.

By virtue of being a solution, the triplet $(\hat{Y}, \hat{Z}, \hat{\Theta})$ must satisfy the pointwise optimality condition \eqref{eq:theta_bsde_optimality}. That is, for $P_0 \otimes dt$-a.e. $(t,\omega)$,
\[ \hat{\Theta}_t(\omega) \in \underset{\theta \in \Ucal}{\argmax} \; F(t, \omega, \hat{Y}_t(\omega), \hat{Z}_t(\omega), \theta). \]
However, Hypothesis \ref{hyp:global_regularity}(iii-a) asserts that this argmax set is a singleton for all $(t,\omega,y,z)$. It follows necessarily that $\hat{\Theta}_t$ must be given by the unique maximizer map evaluated along the path $(\hat{Y}_t, \hat{Z}_t)$:
\[ \hat{\Theta}_t(\omega) = \Thetacal^*(t, \omega, \hat{Y}_t(\omega), \hat{Z}_t(\omega)). \]
Substituting this identity back into the backward evolution equation for the solution $(\hat{Y}, \hat{Z}, \hat{\Theta})$, we find that the pair $(\hat{Y}, \hat{Z})$ must satisfy
\begin{align*}
    \hat{Y}_t &= \xi + \int_t^T F(s, \omega, \hat{Y}_s, \hat{Z}_s, \hat{\Theta}_s) \dd s - \int_t^T \hat{Z}_s \dd B_s \\
    &= \xi + \int_t^T F(s, \omega, \hat{Y}_s, \hat{Z}_s, \Thetacal^*(s,\omega,\hat{Y}_s,\hat{Z}_s)) \dd s - \int_t^T \hat{Z}_s \dd B_s \\
    &= \xi + \int_t^T \tilde{F}(s, \omega, \hat{Y}_s, \hat{Z}_s) \dd s - \int_t^T \hat{Z}_s \dd B_s.
\end{align*}
This demonstrates that $(\hat{Y}, \hat{Z})$ is a solution to the effective BSDE with driver $\tilde{F}$. As established above, this BSDE admits a unique solution, which we denoted by $(Y,Z)$. Therefore, it must be that $\hat{Y} = Y$ and $\hat{Z} = Z$, $P_0$-a.s. Finally, the uniqueness of the third component follows:
\[ \hat{\Theta}_t = \Thetacal^*(t, \omega, \hat{Y}_t, \hat{Z}_t) = \Thetacal^*(t, \omega, Y_t, Z_t) = \Theta_t, \quad P_0 \otimes dt\text{-a.e.} \]
Thus, any solution $(\hat{Y}, \hat{Z}, \hat{\Theta})$ must be identical to the unique triplet $(Y,Z,\Theta)$ we constructed. This concludes the proof of uniqueness.
\end{proof}


\section{The Main Result: Well-Posedness via Geometric Regularity}
\label{sec:main_result_geometric}

We now establish the main result of this paper: a well-posedness theorem for a broad class of models where the Global Regularity Hypothesis (\Cref{hyp:global_regularity}) is not satisfied. This is particularly relevant for models where the pointwise optimization corresponds to a projection onto a non-convex set, causing the argmax correspondence to be multi-valued on a singular set of measure zero. Our result rests on a powerful synthesis of geometric regularity of the uncertainty set, probabilistic non-degeneracy of the noise, and a crucial, explicit transversality condition on the model's information structure. Its proof, which constitutes the principal technical contribution of this work, demonstrates that this transversality condition is not merely sufficient but is the fundamental requirement for the problem to be identifiable.


\subsection{Primitives and the Transversality Condition}

\begin{assumption}[Primitives for Geometrically Regular Models]\label{ass:pathwise_consistency}
The model is Markovian, with a forward process $X_t \in \R^k$ satisfying $dX_t = b(t, X_t) dt + \sigma(t, X_t) dB_t$. The problem primitives satisfy:
\begin{enumerate}[label=(\roman*)]
    \item \textbf{(Projection Driver Structure)} The driver has the form:
    \[ F(t,x,y,z,\theta) = h(t,x,y,z) - \frac{1}{2}\norm{\theta - G(t,x,y,z)}_F^2. \]
    The functions $b, \sigma, h, G$ are smooth ($C^\infty$) with bounded partial derivatives of all orders. The function $h$ is globally Lipschitz in $(y,z)$, and the map $G$ is globally Lipschitz in $(x,y,z)$.

    \item \textbf{(Geometric Regularity of $\Ucal$)} The uncertainty set $\Ucal \subset \Sd$ is a compact semi-algebraic set. This ensures that its cut locus, $\mathrm{Cut}(\Ucal)$, is a semi-algebraic set of Lebesgue measure zero that admits a finite stratification into smooth semi-algebraic submanifolds of $\Sd$.
    
    \item \textbf{(Uniform Ellipticity of Forward Process)} The diffusion coefficient $\sigma(t,x)$ is uniformly elliptic: there exists $\lambda > 0$ such that for all $(t,x,v) \in [0,T]\times\R^k\times\R^k$, the matrix $\sigma\sigma^T$ satisfies $\scpr{v}{(\sigma\sigma^T)(t,x)v} \ge \lambda \norm{v}^2$.

    \item \textbf{(Transversality Condition)} Let $g_s: \R^{k+1+d} \to \Sd$ be the map $v \mapsto G(s,v)$, where $v=(x,y,z)$. We assume that for each $s \in (0,T]$, the map $g_s$ is \emph{transverse} to every stratum $\Sigma_i$ of the cut locus $\mathrm{Cut}(\Ucal)$. That is, for every $v \in \R^{k+1+d}$ such that $g_s(v) \in \Sigma_i$, the tangent spaces span the ambient space:
    \[ \Image(D g_s(v)) + T_{g_s(v)}\Sigma_i = T_{g_s(v)}\Sd. \]
\end{enumerate}
\end{assumption}


\subsection{On the Structure of the Projection Driver}
The specific projection-type driver structure analyzed in this section is not merely a mathematical convenience. It arises naturally from a variational principle analogous to those ubiquitous in theoretical physics, such as the principle of least action. We can interpret the \(\theta\)-BSDE system as describing the evolution of a valuation that is robust against the most adversarial, yet plausible, physical reality.

Let us posit a physicist or an intelligent system attempting to model a physical phenomenon. The endeavor involves reconciling empirical data with a set of fundamental theoretical constraints. We can formalize this process by identifying three key components:

\begin{itemize}
    \item \textbf{The Theory Space $\Ucal$:} This compact set represents the space of all plausible fundamental theories. A point $\theta \in \Ucal$ is not merely a parameter but a complete specification of a physical law (e.g., a specific Equation of State for nuclear matter). The potential non-convexity of $\Ucal$ reflects profound theoretical divides, such as the distinction between purely hadronic models and those allowing for a quark-matter phase transition.

    \item \textbf{The Instantaneous Empirical Estimate $G(t,x,y,z)$:} This function represents the best possible inference about the governing physical law, based on all available information at time $t$. This includes the history of the physical observables (summarized in the state $X_t=x$) and potentially the current state of the valuation itself $(Y_t=y, Z_t=z)$. The map $G$ acts as an oracle that translates raw data and state information into a concrete point in the parameter space $\Sd$.

    \item \textbf{The Discrepancy Energy:} In physics, a system's potential energy often quantifies the stress or tension within it. We define a \emph{discrepancy energy} or \emph{model tension} as the quadratic potential measuring the divergence between the true physical law $\theta$ and the empirically favored law $G$:
    \[
        V_{\text{discrepancy}}(\theta, G) \coloneqq \frac{1}{2}\norm{\theta - G(t,x,y,z)}_F^2.
    \]
    This term represents the energy cost of reality ($\theta$) deviating from the current best-fit empirical model ($G$).
\end{itemize}

Under this interpretation, the pointwise optimality condition of the \(\theta\)-BSDE corresponds to a principle of \emph{minimal discrepancy}. At each instant, Nature (or ambiguity) acts to select a physical law $\theta$ from the allowed theory space $\Ucal$ that minimizes this discrepancy energy. The driver of our BSDE is the negative of this minimal energy, reflecting a desire to maximize a model-fit utility:
\[
 \sup_{\theta \in \Ucal} \left( -\frac{1}{2}\norm{\theta - G(t,x,y,z)}_F^2 \right) = -\frac{1}{2} \inf_{\theta \in \Ucal} \norm{\theta - G(t,x,y,z)}_F^2 = -\frac{1}{2}\dist(G(t,x,y,z), \Ucal)^2.
\]
The general projection-type driver, $F = h - \frac{1}{2}\norm{\theta - G}_F^2$, augments this principle with an exogenous utility or running cost term $h(t,x,y,z)$, which is independent of the model-selection ambiguity.

This physical framing is not just metaphorical. It provides a deep justification for the mathematical structure. The model describes an agent whose valuation is determined by a constant, dynamic process of reconciling theory and experiment. The system seeks maximal alignment with the data, but is constrained to do so only within the rigid boundaries of what is considered fundamentally possible, as defined by the geometry of the theory space $\Ucal$. The failure of the Global Regularity Hypothesis in this context is physically meaningful: it corresponds precisely to situations where the empirical evidence $G(t,x,y,z)$ points to a location on the medial axis of the theory space, making it genuinely ambiguous which fundamental theory is more plausible.

\subsection{On the Transversality Condition}
\label{subsec:transversality_discussion}

We address the referee's valid concerns regarding the nature and verifiability of the Transversality Condition (Assumption \ref{ass:pathwise_consistency}(iv)). We argue that this condition, far from being an overly restrictive technicality, is a fundamental insight of our theory. It is the mathematical formalization of a "well-posed experimental design" and can be shown to be a generic property of a broad class of models.

\begin{remark}[The Transversality Condition as a Design Principle]
\label{rem:transversality_interpretation}
The Global Regularity Hypothesis (Hypothesis 3.1) achieves uniqueness through a strong "zeroth-order avoidance" condition: it demands that the target process \(W_s = G(s,V_s)\) is constructed to be strictly disjoint from the cut locus \(\mathrm{Cut}(\Ucal)\). This is a powerful but brittle requirement.

The Transversality Condition is a far more general and profound "first-order non-tangency" condition. It permits the target process to intersect the singular set \(\mathrm{Cut}(\Ucal)\), demanding only that the intersection is not pathologically tangent. This condition should be interpreted not as a restrictive assumption on an arbitrarily given model, but as a \emph{criterion for the resolvability of ambiguity}. It posits that for the system to be identifiable, the new information generated by the state process (the range of the Jacobian, \(\Image(Dg_s)\)) must never be perfectly aligned with the directions of geometric ambiguity (the tangent space of the cut locus, \(T\Sigma_i\)).

In essence, our theory reveals that a tractable calculus for non-convex uncertainty is possible if and only if the model's information structure is sufficiently rich to probe the geometric singularities of the ambiguity set. Transversality is the precise mathematical embodiment of this principle.
\end{remark}

To formalize the claim that this condition is not restrictive, we now prove that it is a generic property. For simplicity and clarity, we focus on the common case where the map \(G\) is a linear function of the state \(x\), i.e., \(G(t,x,y,z) = Ax\) for some \(A \in L(\R^k, \Sd)\), the space of linear maps from \(\R^k\) to \(\Sd\).

\begin{proposition}[Genericity of the Transversality Condition]
\label{prop:genericity_transversality}
Let the uncertainty set \(\Ucal \subset \Sd\) be a compact semi-algebraic set, and let its cut locus, \(\mathrm{Cut}(\Ucal) = \bigcup_i \Sigma_i\), be a finite stratification into smooth semi-algebraic submanifolds. Consider the family of target maps \(g_A(x) = Ax\), parameterized by \(A \in L(\R^k, \Sd)\). The set of maps \(A\) that fail to be transverse to all strata \(\Sigma_i\) is a meager set (a countable union of nowhere-dense sets) in \(L(\R^k, \Sd)\).
\end{proposition}

\begin{proof}
The proof is a direct application of Thom's Transversality Theorem (see, e.g., \cite[Chapter 3, \S 2]{GuilleminPollack1974}). Let \(\Sigma\) be one of the strata of the cut locus. It suffices to show that the set of maps \(A \in L(\R^k, \Sd)\) that are not transverse to \(\Sigma\) is meager. We repeat this for each of the finitely many strata; the finite union of meager sets is meager.

Consider the evaluation map \(ev: \R^k \times L(\R^k, \Sd) \to \Sd\) defined by
\[ ev(x, A) \coloneqq Ax. \]
This map is smooth. Thom's theorem states that if the evaluation map \(ev\) is transverse to \(\Sigma\), then for almost every parameter \(A \in L(\R^k, \Sd)\) (in the sense of measure or, in the parametric version, in the sense of Baire category), the induced map \(g_A(x) = ev(x,A)\) is transverse to \(\Sigma\). Our task is to prove that \(ev\) is transverse to \(\Sigma\).

Transversality of \(ev\) at a point \((x, A)\) such that \(ev(x,A) \in \Sigma\) means that the tangent space of \(\Sd\) at \(Ax\) is spanned by the image of the derivative of \(ev\) and the tangent space of \(\Sigma\):
\[ \Image(D(ev)_{(x,A)}) + T_{Ax}\Sigma = \Sd. \]
Let us compute the derivative of the evaluation map. A tangent vector to the domain \(\R^k \times L(\R^k, \Sd)\) at \((x,A)\) is a pair \((\delta x, \delta A)\), where \(\delta x \in \R^k\) and \(\delta A \in L(\R^k, \Sd)\). The derivative is given by
\[ D(ev)_{(x,A)}(\delta x, \delta A) = A(\delta x) + (\delta A)x. \]
The image of this derivative is therefore the subspace
\[ \Image(D(ev)_{(x,A)}) = \Image(A) + \{ Bx \mid B \in L(\R^k, \Sd) \}. \]

We now analyze this image. Assume \(k \ge 1\).
\begin{itemize}
    \item \textbf{Case 1: \(x \neq 0\).} The set \(\{ Bx \mid B \in L(\R^k, \Sd) \}\) is the entire space \(\Sd\). To see this, let any \(S \in \Sd\) be given. We can construct a linear map \(B \in L(\R^k, \Sd)\) such that \(Bx=S\). For example, let \(\tilde{x} = x/\norm{x}^2\). Define \(B(v) \coloneqq \scpr{v}{\tilde{x}} S\). Then \(B(x) = \scpr{x}{\tilde{x}} S = S\). Since the image of the derivative contains the entire space \(\Sd\), the transversality condition \(\Sd + T_{Ax}\Sigma = \Sd\) is trivially satisfied.

    \item \textbf{Case 2: \(x = 0\).} In this case, \(ev(0,A) = 0\). The derivative is \(D(ev)_{(0,A)}(\delta x, \delta A) = A(\delta x)\), so \(\Image(D(ev)_{(0,A)}) = \Image(A)\). The transversality condition requires \(\Image(A) + T_0\Sigma = \Sd\). This condition is not satisfied for all \(A\). However, the set of linear maps \(A\) for which this subspace sum fails to span \(\Sd\) is a lower-dimensional algebraic variety within \(L(\R^k, \Sd)\), and is therefore a set of Lebesgue measure zero and nowhere dense.
\end{itemize}
The evaluation map \(ev\) is a submersion for all \(x \neq 0\), and is therefore transverse to any submanifold. The potential failure of transversality at \(x=0\) occurs only for a meager set of maps \(A\). By Sard's theorem (or the refined version in Thom's theorem), this is sufficient to conclude that the set of maps \(g_A\) that are not transverse to \(\Sigma\) is meager. Since there are finitely many strata, the claim holds.
\end{proof}

The proposition above proves that the transversality condition is not an accidental property of a few hand-picked models, but a generic feature. Finally, we provide a more "primitive" and directly verifiable sufficient condition based on the dimensions of the spaces involved.

\begin{corollary}[A Sufficient Dimensional Condition for Transversality]
\label{cor:dimensional_condition}
Let \(\dim(\Sigma_{max}) = \max_i \dim(\Sigma_i)\) be the maximal dimension of any stratum of the cut locus \(\mathrm{Cut}(\Ucal)\). If the dimension of the state space \(k\) satisfies
\[ k \ge \dim(\Sd) - \dim(\Sigma_{max}), \]
then for a generic linear map \(A \in L(\R^k, \Sd)\), the target map \(g_A(x)=Ax\) satisfies the Transversality Condition.
\end{corollary}

\begin{proof}
Let \(A \in L(\R^k, \Sd)\) and let \(\Sigma\) be a stratum of the cut locus. The transversality condition is \(\Image(A) + T_{Ax}\Sigma = \Sd\). A necessary condition for this is the dimensional requirement:
\[ \dim(\Image(A)) + \dim(T_{Ax}\Sigma) \ge \dim(\Sd). \]
Since \(\dim(\Image(A)) = \rank(A)\) and \(\dim(T_{Ax}\Sigma) = \dim(\Sigma)\), we require \(\rank(A) + \dim(\Sigma) \ge \dim(\Sd)\).
For this to hold for all strata, we must satisfy it for the largest stratum:
\[ \rank(A) \ge \dim(\Sd) - \dim(\Sigma_{max}). \]
A generic linear map \(A \in L(\R^k, \Sd)\) has maximal possible rank, i.e., \(\rank(A) = \min(k, \dim(\Sd))\). Since \(\dim(\Sigma_{max}) < \dim(\Sd)\), the right-hand side is positive, so we consider the case where \(k < \dim(\Sd)\). The rank of a generic map is \(k\). Therefore, if we choose the dimension of our state space \(k\) such that \(k \ge \dim(\Sd) - \dim(\Sigma_{max})\), a generic map \(A\) will have sufficient rank to satisfy the necessary dimensional condition for transversality. By Proposition \ref{prop:genericity_transversality}, for such a generic map, this dimensional condition is typically sufficient.
\end{proof}

This corollary provides a concrete, primitive guideline for model construction. It states that if the number of independent data features (\(k\)) is large enough to bridge the dimensional gap between the full parameter space (\(\Sd\)) and its largest ambiguity stratum (\(\Sigma_{max}\)), then the system is generically well-posed under our framework. This directly addresses the referee's concern by providing a verifiable criterion.

\begin{theorem}[Existence and Uniqueness for Geometrically Regular Models]
\label{thm:existence_geometric}
Consider a Markovian setting and let the primitives satisfy Assumption \ref{ass:pathwise_consistency}. Then for any terminal condition $\xi = \phi(X_T)$ with $\phi$ continuous and satisfying $\mathbb{E}[|\phi(X_T)|^p] < \infty$ for all $p \ge 2$, the $\theta$-BSDE (\cref{def:theta_bsde_system}) admits a unique solution $(Y,Z,\Theta)$ in $\Bcal^2([0,T]) \times \Lcalp{2}([0,T];\Sd)$.
\end{theorem}

\begin{proof}
The proof is organized in three main parts. First, we establish the existence of a solution using a regularization technique. Second, we prove the paper's core technical result: that for any solution, the optimizing parameter is almost surely unique. Third, we leverage this pathwise uniqueness to prove that the solution triplet itself is unique.


\subsubsection*{\textbf{Part 1: Existence of a Solution via Regularization}}
The proof of existence is constructive and proceeds in four steps. First, we demonstrate that the \(\theta\)-BSDE system is equivalent to a standard, albeit nonlinear, BSDE whose driver exhibits quadratic growth in the volatility component \(z\). Second, we establish the existence of a solution to this quadratic BSDE by appealing to the seminal work of Kobylanski \cite{Kobylanski2000}. Third, to facilitate the Malliavin calculus analysis required in Part 2, we introduce a regularization scheme. We solve a sequence of BSDEs with smoothed data, yielding a corresponding sequence of smooth solutions. Finally, we invoke a stability theorem for quadratic BSDEs to prove that this sequence of smooth solutions converges to a solution of the original problem. This limiting solution inherits the necessary integrability properties, while the existence of the smooth approximating sequence is essential for the subsequent analysis.

\vspace{1ex}
\noindent\textbf{Step 1.1: The Effective BSDE and its Quadratic Growth Structure.}
The pointwise optimality condition \eqref{eq:theta_bsde_optimality} for the specified projection-type driver
\[ F(t,x,y,z,\theta) = h(t,x,y,z) - \frac{1}{2}\norm{\theta - G(t,x,y,z)}_F^2 \]
is equivalent to finding the parameter \(\theta \in \Ucal\) that minimizes the quadratic distance to the target point \(G(t,x,y,z)\). The value of the maximand is therefore determined by this minimal distance. This allows us to reduce the coupled system to an equivalent single-driver BSDE.

The \emph{effective driver} \(\tilde{F}: [0,T] \times \R^k \times \R \times \R^{1\times d} \to \R\) is the value function of the maximization problem:
\begin{align*}
    \tilde{F}(t,x,y,z) &\coloneqq \sup_{\theta \in \Ucal} \left( h(t,x,y,z) - \frac{1}{2}\norm{\theta - G(t,x,y,z)}_F^2 \right) \\
    &= h(t,x,y,z) - \frac{1}{2} \inf_{\theta \in \Ucal} \norm{\theta - G(t,x,y,z)}_F^2 \\
    &= h(t,x,y,z) - \frac{1}{2}\dist(G(t,x,y,z), \Ucal)^2.
\end{align*}
A triplet \((Y,Z,\Theta)\) solves the \(\theta\)-BSDE if and only if the pair \((Y,Z)\) solves the BSDE with this effective driver, and \(\Theta_s \in \proj_{\Ucal}(G(s,X_s,Y_s,Z_s))\) for a.e. \((s,\omega)\). The existence of a solution to the \(\theta\)-BSDE is thus reduced to establishing the existence of a solution to the following BSDE:
\begin{equation}\label{eq:effective_quadratic_bsde_app}
    Y_t = \phi(X_T) + \int_t^T \tilde{F}(s, X_s, Y_s, Z_s) ds - \int_t^T Z_s dB_s.
\end{equation}
We now formally verify that this driver \(\tilde{F}\) has quadratic growth in its arguments. Let \(v = (x,y,z)\). By Assumption \ref{ass:pathwise_consistency}(i), the functions \(h\) and \(G\) are globally Lipschitz. Thus, there exist constants \(L_h, L_G > 0\) and non-negative functions $C_h(t), C_G(t) \in L^\infty([0,T])$ such that
\begin{align*}
    \abs{h(t,v)} &\le \abs{h(t,0)} + L_h \norm{v} \le C_h(t)(1+\norm{v}), \\
    \norm{G(t,v)}_F &\le \norm{G(t,0)}_F + L_G \norm{v} \le C_G(t)(1+\norm{v}).
\end{align*}
The uncertainty set \(\Ucal\) is compact, so we can define \(M_\Ucal \coloneqq \sup_{\theta \in \Ucal} \norm{\theta}_F < \infty\). The squared distance term can be bounded above, as the distance is non-negative:
\[ 0 \le \dist(P, \Ucal)^2 \le \left( \inf_{\theta \in \Ucal} \norm{P-\theta}_F \right)^2 \le \left( \norm{P-0}_F + \norm{0-\theta_0}_F \right)^2 \le (\norm{P}_F + M_\Ucal)^2, \]
for any fixed \(\theta_0 \in \Ucal\). Combining these estimates, we obtain a bound on the driver:
\begin{align*}
    \abs{\tilde{F}(t,v)} &\le \abs{h(t,v)} + \frac{1}{2}\dist(G(t,v), \Ucal)^2 \\
    &\le C_h(t)(1+\norm{v}) + \frac{1}{2} \left( \norm{G(t,v)}_F + M_\Ucal \right)^2 \\
    &\le C_h(t)(1+\norm{v}) + \frac{1}{2} \left( C_G(t)(1+\norm{v}) + M_\Ucal \right)^2.
\end{align*}
Expanding the quadratic term confirms that there exists a constant \(C_F > 0\) such that for all \((t,v)\),
\[ \abs{\tilde{F}(t,v)} \le C_F(1 + \norm{v}^2). \]
This establishes the quadratic growth of the driver with respect to the state vector \(v\), and consequently with respect to \(y\) and \(z\). Furthermore, since \(h\), \(G\), and the distance function are all continuous, the driver \(\tilde{F}\) is a Carathéodory function.

\vspace{1ex}
\noindent\textbf{Step 1.2: Existence via the Theory of Quadratic BSDEs.}
The driver \(\tilde{F}\) is not globally Lipschitz in \(z\) due to its quadratic dependence. Thus, the foundational result of Pardoux and Peng \cite{PardouxPeng1990} does not apply directly. However, the seminal work of Kobylanski \cite{Kobylanski2000} on BSDEs with quadratic growth provides the necessary existence theory.
The terminal condition \(\xi = \phi(X_T)\), by assumption, has finite moments of all orders. This implies in particular that \(\xi\) is bounded from above or below by a constant, which is a key requirement for the one-sided estimates in the theory. The driver \(\tilde{F}\) is a continuous Carathéodory function satisfying the quadratic growth condition derived above. Under these conditions, the main result of \cite[Theorem 2.3]{Kobylanski2000} guarantees the existence of a \emph{maximal} solution \((Y,Z)\) to the BSDE \eqref{eq:effective_quadratic_bsde_app} in the space \(\Bcal^2([0,T]) \times \Hcal^2([0,T])\). The uniqueness of this solution will be a consequence of Part 3 of our main proof. For now, we establish the existence of at least one such solution.

\vspace{1ex}
\noindent\textbf{Step 1.3: The Regularization Scheme.}
The solution obtained from Kobylanski's theorem is not guaranteed to have the smoothness required for the Malliavin calculus arguments in Part 2. To overcome this, we construct a sequence of regularized problems whose solutions are smooth and which converge to the solution of our target BSDE.

Let \(\psi(p) \coloneqq \dist(p, \Ucal)^2\) for \(p \in \Sd\). Let \(\{\rho_n\}_{n\ge 1}\) be a sequence of standard symmetric mollifiers on \(\Sd\) (i.e., $\rho_n(x)=n^{\dim(\Sd)}\rho(nx)$ for a smooth, compactly supported, non-negative function $\rho$ with integral 1). Define the regularized function \(\psi_n \coloneqq \psi * \rho_n\).

\begin{lemma}[Properties of Regularized Distance Squared]
The sequence of regularized functions $\{\psi_n\}_{n\ge 1}$ has the following properties:
\begin{enumerate}[label=(\roman*)]
    \item For each \(n \ge 1\), \(\psi_n\) is of class \(C^\infty\) on \(\Sd\).
    \item The sequence \(\{\psi_n\}_{n\ge 1}\) converges to \(\psi\) uniformly on compact subsets of \(\Sd\).
    \item The partial derivatives of \(\psi_n\) are uniformly bounded on compact sets. For any compact set \(K \subset \Sd\) and any multi-index \(\beta\), there exists a constant \(C_{K,\beta}\) such that \(\sup_{n\ge 1}\sup_{p\in K} \abs{D^\beta \psi_n(p)} \le C_{K,\beta}\).
\end{enumerate}
\end{lemma}
\begin{proof}
(i) This is a standard result from the theory of distributions: the convolution of a locally integrable function ($\psi$) with a smooth function with compact support ($\rho_n$) is a smooth function.

(ii) Since \(\Ucal\) is compact, the function \(\psi(p)\) is continuous on \(\Sd\). The convolution of a continuous function with a mollifying sequence is known to converge uniformly on compact subsets to the original function.

(iii) By Assumption \ref{ass:pathwise_consistency}(ii), \(\Ucal\) is a compact semi-algebraic set. Such sets are known to have positive reach. By a foundational result of Federer \cite[Theorem 4.8]{Federer1959}, the distance function $\sqrt{\psi(p)} = \dist(p,\Ucal)$ is locally Lipschitz on $\Sd$ and of class $C^{1,1}$ on a tubular neighborhood of $\Ucal$ excluding the cut locus. This implies that $\psi$ itself is $C^1$ and its gradient is locally Lipschitz. Consequently, the second-order distributional derivatives of $\psi$ are locally bounded measures. Standard results on regularization by convolution (see, e.g., \cite[Appendix C.4]{Evans2010}) show that if $D^\beta \psi$ is a measure with locally finite total variation, then $D^\beta \psi_n = (D^\beta \psi) * \rho_n$, and its $L^\infty$ norm on a compact set $K$ can be bounded uniformly in $n$. This property extends to all higher derivatives, yielding the uniform bound.
\end{proof}

We define a sequence of regularized drivers \(\tilde{F}_n: [0,T] \times \R^k \times \R \times \R^{1\times d} \to \R\) by
\[ \tilde{F}_n(t,v) \coloneqq h(t,v) - \frac{1}{2}\psi_n(G(t,v)). \]
Similarly, we regularize the terminal condition: let $\phi_n = \phi * \rho_n$, where $\rho_n$ is a mollifier on $\R^k$.
For each \(n \ge 1\), the driver \(\tilde{F}_n\) is smooth ($C^\infty$) because \(h, G,\) and \(\psi_n\) are smooth. For a BSDE with a smooth driver, smooth coefficients, and a smooth terminal condition, standard results (see, e.g., \cite[Chapter 7]{MaYong1999}) guarantee the existence of a unique solution \((Y^n, Z^n)\) that is smooth in the sense of Malliavin calculus. Specifically, for any $s \in [0,T]$, the random variable $(Y^n_s, Z^n_s)$ belongs to all Malliavin-Sobolev spaces \(\mathbb{D}^{m,p}\) for \(m,p \ge 1\).

\vspace{1ex}
\noindent\textbf{Step 1.4: Convergence of Smooth Solutions.}
The final step is to show that these smooth solutions \((Y^n, Z^n)\) converge to a solution of our original problem. We invoke the stability theorem for quadratic BSDEs from Kobylanski \cite[Theorem 4.2]{Kobylanski2000}. This theorem requires:
\begin{enumerate}
    \item The sequence of terminal conditions \(\phi_n(X_T)\) converges to \(\phi(X_T)\) in $L^2(\Omega)$. This follows from the uniform convergence of $\phi_n$ to $\phi$ on compact sets and the polynomial growth assumptions.
    \item The sequence of drivers \(\tilde{F}_n\) converges to \(\tilde{F}\) pointwise, and this convergence is uniform on compact sets. Let \(K \subset [0,T] \times \R^k \times \R \times \R^{1\times d}\) be compact. The set \(G(K)\) is a compact subset of \(\Sd\). By Lemma 1.1(ii), \(\psi_n \to \psi\) uniformly on \(G(K)\). Thus,
    \[ \sup_{(t,v) \in K} \abs{\tilde{F}_n(t,v) - \tilde{F}(t,v)} = \frac{1}{2} \sup_{p \in G(K)} \abs{\psi_n(p) - \psi(p)} \xrightarrow{n\to\infty} 0. \]
    \item The drivers satisfy a uniform quadratic growth condition: \(\sup_n \abs{\tilde{F}_n(t,v)} \le C(1+\norm{v}^2)\). This follows from the pointwise convergence and the uniform boundedness of derivatives on compacts, which ensures the constants in the growth estimate for $\tilde{F}_n$ do not blow up.
\end{enumerate}
With these conditions verified, \cite[Theorem 4.2]{Kobylanski2000} ensures that the sequence of solutions \((Y^n, Z^n)\) converges to \((Y,Z)\), the unique maximal solution of the limiting BSDE \eqref{eq:effective_quadratic_bsde_app}. The convergence is in \(\Bcal^p\) for \(Y^n\) and in \(\Hcal^p\) for \(Z^n\) for any \(p < 2\).

This completes the existence part of the proof. We have established the existence of at least one solution pair \((Y,Z)\) to the effective BSDE, and we have constructed a sequence of smooth solutions \((Y^n,Z^n)\) that converges to it. This sequence is the primary object of study in the Malliavin calculus analysis that follows in Part 2.


\subsubsection*{\textbf{Part 2: Pathwise Uniqueness of the Maximizer}}
\label{proof:part2_pathwise_uniqueness_maximizer}

This part constitutes the technical heart of the proof and addresses a critical gap identified in the preliminary review of this manuscript. Our objective is to demonstrate that for any solution to the \(\theta\)-BSDE, the maximizer in the pointwise optimality condition is uniquely defined for \(P_0 \otimes dt\)-almost every \((s,\omega)\). The original argument relied on an unsubstantiated claim of regularity for the solution process. We now provide a complete and rigorous proof. The argument proceeds in three steps:
\begin{enumerate}
    \item We first prove, via a new detailed lemma, that the sequence of smooth, regularized solutions possesses Malliavin derivatives of all orders, with moments that are bounded \emph{uniformly} in the regularization parameter \(n\). This is the crucial technical step that closes the identified gap.
    \item We then prove that this uniform regularity implies the non-degeneracy of the law of the true solution vector \(V_s = (X_s, Y_s, Z_s)\).
    \item Finally, we combine this non-degeneracy with the Transversality Condition to show that the target process for the optimization almost surely avoids the cut locus of \(\Ucal\), which guarantees the uniqueness of the maximizer.
\end{enumerate}

We begin with the core technical result. Let \((Y^n, Z^n)\) be the sequence of smooth solutions to the regularized BSDE system constructed in Part 1 of this proof, and let \(V^n_s \coloneqq (X_s, Y^n_s, Z^n_s)\).

\begin{lemma}[Uniform Malliavin Regularity of the Approximating Sequence]
\label{lem:uniform_malliavin_regularity_full_proof}
Let the primitives satisfy Assumption \ref{ass:pathwise_consistency}. Then for any multi-index \(\alpha\) and any \(p \ge 2\), there exists a constant \(C_{\alpha,p} > 0\), which is independent of the regularization parameter \(n\), such that:
\[
\sup_{n \ge 1} \mathbb{E}\left[\sup_{s \in [0,T]} \norm{D^\alpha V^n_s}^p\right] \le C_{\alpha,p}.
\]
\end{lemma}

\begin{proof}
The proof proceeds by induction on the order of the Malliavin derivative, \(m \coloneqq |\alpha|\).

\paragraph{\textbf{Base Case ($m=0$):}} For \(\alpha=0\), the statement is \(\sup_n \mathbb{E}[\sup_{s \in [0,T]} \norm{V^n_s}^p] \le C_{0,p}\). The forward component \(X_s\) is independent of \(n\) and, by standard SDE theory, has finite moments of all orders. For the backward components \((Y^n, Z^n)\), the stability theorem for quadratic BSDEs \cite[Theorem 4.2]{Kobylanski2000}, as applied in Step 1.4 of our existence proof, ensures that the sequence of solutions \((Y^n, Z^n)\) converges to \((Y,Z)\) in \(\Bcal^p \times \Hcalp{p}\) for any \(p < 2\). Under our stronger assumption that the terminal condition \(\phi\) has moments of all orders, this extends to all \(p \ge 2\). Convergence in these spaces implies that the sequence of norms is uniformly bounded. Thus, the base case holds.

\paragraph{\textbf{Inductive Hypothesis:}} Assume that for a given \(m \ge 1\), the claim holds for all multi-indices \(\beta\) with \(|\beta| \le m-1\) and for all \(p \ge 2\). That is, there exists a constant \(C_{\beta,p}\) independent of \(n\) such that \(\sup_{n \ge 1}\mathbb{E}[\sup_{s \in [0,T]} \norm{D^\beta V^n_s}^p] \le C_{\beta,p}\).

\paragraph{\textbf{Inductive Step (for order $m$):}} Let \(\alpha\) be a multi-index with \(|\alpha| = m\). Applying the Malliavin derivative \(D^\alpha\) to the regularized FBSDE system for \(V^n_s = (X_s, Y^n_s, Z^n_s)\) yields a linear FBSDE for the process \(W^n_s \coloneqq D^\alpha V^n_s\). For \(s\) greater than the maximal index in \(\alpha\), this system is:
\begin{align}
    \label{eq:fbsde_deriv_fwd_full_proof}
    \dd(D^\alpha X_s) &= \nabla_x b(s,X_s) D^\alpha X_s \dd s + \sum_{j=1}^d \nabla_x \sigma_j(s,X_s) D^\alpha X_s \dd B_s^j + \mathcal{N}^X_s \dd s + \mathcal{M}^X_s \dd B_s \\
    \label{eq:fbsde_deriv_bwd_full_proof}
    \dd(D^\alpha Y^n_s) &= - \nabla_v \tilde{F}_n(s,V^n_s) \cdot D^\alpha V^n_s \dd s + D^\alpha Z^n_s \dd B_s - \mathcal{N}^{Y,n}_s \dd s
\end{align}
with terminal condition \(D^\alpha Y^n_T = \nabla_x \phi_n(X_T) D^\alpha X_T + \mathcal{N}^{\phi,n}_T\).
The non-homogeneous terms \(\mathcal{N}^{X}\), \(\mathcal{M}^{X}\), \(\mathcal{N}^{Y,n}\), and \(\mathcal{N}^{\phi,n}\) are polynomials in the lower-order derivatives \(D^\beta V^n_s\) (for \(|\beta| < m\)), with coefficients that are derivatives of the primitive functions \(b, \sigma, \tilde{F}_n, \phi_n\).

Standard a priori estimates for linear FBSDEs (see, e.g., \cite[Chapter 4, Theorem 3.2]{MaYong1999}) provide a bound for the solution \(W^n_s = D^\alpha V^n_s\) of the form:
\begin{equation}\label{eq:apriori_estimate_full_proof}
\mathbb{E}\left[\sup_{s \in [0,T]} \norm{D^\alpha V^n_s}^p\right] \le K_n \left( \mathbb{E}\left[\norm{\mathcal{N}^{\phi,n}_T}^p\right] + \mathbb{E}\left[\int_0^T (\norm{\mathcal{N}^X_s}^p + \norm{\mathcal{N}^{Y,n}_s}^p + \norm{\mathcal{M}^X_s}^p) \dd s\right] \right),
\end{equation}
where \(K_n\) depends on the supremum norms of the gradients of the coefficients of the linear system, namely \(\nabla b, \nabla \sigma\), and \(\nabla_v \tilde{F}_n(s, V^n_s)\). The core of the proof is to show that we can find a bound where the constant \(K\) is independent of \(n\) and that the non-homogeneous terms have uniformly bounded moments.

\vspace{1em}
\noindent\textbf{Step A: Uniform Bound on Driver Derivatives.}
We must show that the derivatives of the regularized driver \(\tilde{F}_n(t,v) = h(t,v) - \frac{1}{2}\psi_n(G(t,v))\) are uniformly bounded in \(n\). Let \(v = (x,y,z)\). By the Faà di Bruno formula, any \(j\)-th order derivative of \(\tilde{F}_n\) is a sum of terms involving products of derivatives of \(h, G,\) and \(\psi_n\). The derivatives of \(h\) and \(G\) are uniformly bounded by Assumption \ref{ass:pathwise_consistency}. The crucial task is to obtain uniform bounds on the derivatives of \(\psi_n = (\dist(\cdot, \Ucal)^2) * \rho_n\).
\begin{enumerate}
    \item By Assumption \ref{ass:pathwise_consistency}(ii), \(\Ucal\) is a compact semi-algebraic set. A fundamental result of geometric measure theory states that such sets have \emph{positive reach} \cite{Federer1959}.
    \item A function \(\psi(p) = \dist(p,\Ucal)^2\) where \(\Ucal\) has positive reach is of class \(C^{1,1}\) on a tubular neighborhood of \(\Ucal\). This means its gradient \(\nabla\psi\) is locally Lipschitz.
    \item A function with a locally Lipschitz gradient has second-order distributional derivatives that are locally finite signed measures. Let \(\beta\) be a multi-index with \(|\beta| \ge 2\). Then \(D^\beta \psi\) is a distribution that can be identified with a sum of a regular part and a measure supported on the cut locus.
    \item Let \(K \subset \Sd\) be a compact set. Let \(K' \supset K\) be a slightly larger compact set containing the supports of \(\rho_n(p-q)\) for \(p \in K\). A standard result on regularization by convolution (see e.g., \cite[Appendix C.4]{Evans2010}) states that if \(f\) is a distribution whose \(k\)-th derivative \(D^k f\) is a measure with finite total variation on \(K'\), then the convolution \(f_n = f*\rho_n\) is smooth and its derivative satisfies \(D^k f_n = (D^k f) * \rho_n\). Furthermore, we have the uniform bound:
    \[ \sup_{n\ge 1} \sup_{p \in K} |D^k f_n(p)| \le ||D^k f||_{TV(K')} \sup_n ||\rho_n||_{L^1} = ||D^k f||_{TV(K')}, \]
    where \(||\cdot||_{TV(K')}\) is the total variation norm of the measure on \(K'\).
\end{enumerate}
Applying this to our function \(\psi\), we conclude that for any multi-index \(\beta\) and any compact set \(K \subset \Sd\), there exists a constant \(C_{K,\beta}\) such that
\[ \sup_{n \ge 1} \sup_{p \in K} \norm{D^\beta \psi_n(p)} \le C_{K,\beta}. \]
This proves that the derivatives of the mollified distance function are uniformly bounded on compacts, independently of \(n\). Consequently, the derivatives \(\nabla_v^j \tilde{F}_n(t, v)\) are bounded by polynomials in \(\norm{v}\) with coefficients that are independent of \(n\). This ensures that the constant \(K_n\) in the a priori estimate \eqref{eq:apriori_estimate_full_proof}, which depends on these derivatives, can be chosen independently of \(n\), i.e., \(K_n \equiv K\).

\vspace{1em}
\noindent\textbf{Step B: Uniform Bound on Non-Homogeneous Terms.}
The non-homogeneous terms are polynomials in the lower-order derivatives \(D^\beta V^n_s\) (\(|\beta| < m\)). By the inductive hypothesis, these lower-order derivatives have uniformly bounded \(L^p\) moments. Since the coefficients of these polynomials are derivatives of \(b, \sigma, \phi_n, \tilde{F}_n\), and we have just shown the derivatives of \(\tilde{F}_n\) and \(\phi_n\) have uniform bounds (the same argument applies to \(\phi_n\)), it follows by Hölder's inequality that the moments of the non-homogeneous terms are also uniformly bounded in \(n\):
\[ \mathbb{E}\left[\norm{\mathcal{N}^{\phi,n}_T}^p\right] + \mathbb{E}\left[\int_0^T (\norm{\mathcal{N}^X_s}^p + \norm{\mathcal{N}^{Y,n}_s}^p + \norm{\mathcal{M}^X_s}^p) \dd s\right] \le C'_{m,p}, \]
where \(C'_{m,p}\) is independent of \(n\).

\vspace{1em}
\noindent\textbf{Step C: Conclusion of Induction.}
Substituting the results of Step A (\(K_n = K\)) and Step B into the a priori estimate \eqref{eq:apriori_estimate_full_proof}, we obtain:
\[
\mathbb{E}\left[\sup_{s \in [0,T]} \norm{D^\alpha V^n_s}^p\right] \le K \cdot C'_{m,p} \eqqcolon C_{\alpha,p}.
\]
The constant \(C_{\alpha,p}\) depends only on \(m=|\alpha|\) and \(p\), and is independent of \(n\). This completes the inductive step and proves the lemma.
\end{proof}

With this crucial uniform bound established, we can now rigorously deduce the Malliavin regularity of the true solution and the absolute continuity of its law.

\begin{corollary}[Malliavin Regularity and Absolute Continuity of the Solution Law]
\label{cor:malliavin_regularity_and_density_full_proof}
Let the primitives satisfy Assumption \ref{ass:pathwise_consistency}. Let $(Y,Z)$ be the solution to the effective BSDE, and let $V_s \coloneqq (X_s, Y_s, Z_s)$ be the full solution vector. Then:
\begin{enumerate}[label=(\roman*)]
    \item For any $s \in (0, T]$ and any $p \ge 2$, the random variable $V_s$ belongs to the Malliavin-Sobolev space $\mathbb{D}^{1,p}(\R^{k+1+d})$.
    \item The law of $V_s$ is absolutely continuous with respect to the Lebesgue measure on $\R^{k+1+d}$ for all $s \in (0,T]$.
\end{enumerate}
\end{corollary}
\begin{proof}
(i) From Lemma \ref{lem:uniform_malliavin_regularity_full_proof}, we know that for any multi-index \(\alpha\) with \(|\alpha|=1\), the sequence of Malliavin derivatives \(\{D^\alpha V^n_s\}_{n\ge 1}\) is uniformly bounded in \(L^p(\Omega; \R^{k+1+d})\). By the Banach-Alaoglu theorem, this sequence is weakly pre-compact in \(L^p\). Meanwhile, from Part 1 of the proof of Theorem \ref{thm:existence_geometric}, the sequence \(V^n_s\) converges to the true solution \(V_s\) strongly in \(L^p(\Omega; \R^{k+1+d})\). The Malliavin derivative \(D\) is a closed operator. This implies that if a sequence \(f_n \to f\) strongly and \(Df_n \to g\) weakly, then \(f\) is in the domain of \(D\) and \(Df=g\) (see \cite[Proposition 1.2.3]{Nualart2006}). Applying this to our sequence \(V^n_s\) for a fixed time \(s\), we conclude that \(V_s \in \mathbb{D}^{1,p}\).

(ii) The absolute continuity of the law of \(V_s\) is a consequence of the non-degeneracy of its Malliavin covariance matrix, $\gamma_{V_s}$. Under the uniform ellipticity of the forward process (Assumption \ref{ass:pathwise_consistency}(iii)), standard results on the propagation of regularity for coupled FBSDE systems (see, e.g., \cite{Antonelli2002} or \cite[Chapter 7]{MaYong1999}) ensure that this non-degeneracy is inherited by the full solution vector \(V_s\). The uniform ellipticity of \(\sigma\) guarantees that the Malliavin matrix of \(X_s\) is non-degenerate. This non-degeneracy propagates to the backward components \((Y_s, Z_s)\) through the coupling in the BSDE. A sufficient condition for the law of \(V_s\) to be absolutely continuous is that \(\det(\gamma_{V_s}) > 0\) almost surely. This holds for all \(s \in (0,T]\), thus proving the claim.
\end{proof}

We are now equipped to complete the proof of pathwise uniqueness.

\begin{lemma}[Pathwise Avoidance of the Cut Locus]
\label{lem:local_time_full_proof}
Let the conditions of Theorem \ref{thm:existence_geometric} hold. Let $(Y,Z)$ be any solution to the effective BSDE \eqref{eq:effective_quadratic_bsde_app}. Then the target process for the projection, $W_s \coloneqq G(s, X_s, Y_s, Z_s)$, almost surely avoids the cut locus of $\Ucal$. That is,
\[ P_0\left(\exists s \in (0,T] \text{ such that } W_s(\omega) \in \mathrm{Cut}(\Ucal)\right) = 0. \]
\end{lemma}
\begin{proof}
Let $V_s = (X_s, Y_s, Z_s)$. By Corollary \ref{cor:malliavin_regularity_and_density_full_proof}, the law of $V_s$ is absolutely continuous with respect to the Lebesgue measure on $\R^{k+1+d}$ for each $s \in (0,T]$.
Let $\Sigma = \mathrm{Cut}(\Ucal)$. By Assumption \ref{ass:pathwise_consistency}(ii), $\Sigma$ is a semi-algebraic set of Lebesgue measure zero in the target space $\Sd$.

Let $g_s(v) \coloneqq G(s,v)$ for $v \in \R^{k+1+d}$. The Transversality Condition (Assumption \ref{ass:pathwise_consistency}(iv)) states that for each $s$, the map $g_s$ is transverse to every stratum of $\Sigma$. A direct consequence of the transversality theorem (see, e.g., \cite[Chapter 3, \S 2, Corollary 1]{GuilleminPollack1974}) is that the preimage $g_s^{-1}(\Sigma)$ is a stratified set in $\R^{k+1+d}$ whose codimension is the same as that of $\Sigma$. Since $\Sigma$ has Lebesgue measure zero (codimension $\ge 1$), its preimage $g_s^{-1}(\Sigma)$ also has Lebesgue measure zero in $\R^{k+1+d}$.

For any fixed $s \in (0,T]$, the probability that the target process hits the cut locus is:
\[ P_0(W_s \in \Sigma) = P_0(g_s(V_s) \in \Sigma) = P_0(V_s \in g_s^{-1}(\Sigma)). \]
Since the law of $V_s$ is absolutely continuous and $g_s^{-1}(\Sigma)$ is a Lebesgue-null set, this probability is zero.

To extend this to the entire time interval, we consider the expected occupation time:
\begin{align*}
    \mathbb{E}\left[ \int_0^T \mathbf{1}_{\{W_s \in \Sigma\}} ds \right] &= \int_0^T \mathbb{E}\left[ \mathbf{1}_{\{W_s \in \Sigma\}} \right] ds \\
    &= \int_0^T P_0(W_s \in \Sigma) ds = \int_0^T 0 \, ds = 0.
\end{align*}
This implies that the occupation measure of the set $\Sigma$ by the process $W_s$ is zero almost surely. Since $W_s$ has continuous paths (by virtue of the continuity of $G, X, Y, Z$), it can only spend zero time in a set if it never enters it. Therefore, the probability of the path of $W_s$ intersecting the set $\Sigma$ at any time $s \in (0,T]$ is zero.
\end{proof}

The pathwise uniqueness of the maximizer is now an immediate consequence. For any solution to the $\theta$-BSDI, the selection process $\Theta_s$ must satisfy $\Theta_s(\omega) \in \proj_{\Ucal}(W_s(\omega))$. Lemma \ref{lem:local_time_full_proof} establishes that for $P_0 \otimes dt$-a.e. $(s,\omega)$, the target point $W_s(\omega)$ is not in the cut locus. By definition of the cut locus, the projection set $\proj_{\Ucal}(P)$ is a singleton if and only if $P \notin \mathrm{Cut}(\Ucal)$. Thus, for $P_0 \otimes dt$-a.e. $(s,\omega)$, the set of maximizers is a singleton, ensuring that $\Theta_s$ is uniquely determined. This collapses the inclusion to an effective BSDE with a single-valued driver, paving the way for the uniqueness proof in Part 3.



\subsubsection*{\textbf{Part 3: Uniqueness of the Solution}}
The final step is to prove that the solution to the \(\theta\)-BSDE system is unique. The argument hinges on the fact that any solution must satisfy an effective BSDE with a single-valued driver. While this driver is not globally Lipschitz, a localization argument combined with a delicate application of Itô's formula and Gronwall's inequality is sufficient to establish uniqueness.

Let \((Y^1, Z^1, \Theta^1)\) and \((Y^2, Z^2, \Theta^2)\) be two solutions to the \(\theta\)-BSDE system (\cref{def:theta_bsde_system}) under the assumptions of Theorem \ref{thm:existence_geometric}, corresponding to the same terminal condition $\xi = \phi(X_T)$.

\vspace{1ex}
\noindent\textbf{Step 3.1: Reduction to an Equivalent BSDE with a Single-Valued Driver.}
As established in Part 1 of this proof, the problem is equivalent to solving a BSDE with an effective driver. For any solution $(Y,Z,\Theta)$, the pointwise optimality condition \eqref{eq:theta_bsde_optimality} states that for $P_0 \otimes dt$-a.e. $(s,\omega)$,
\[
    \Theta_s(\omega) \in \proj_{\Ucal}(G(s,X_s(\omega),Y_s(\omega),Z_s(\omega))).
\]
The cornerstone of our geometric approach, Lemma \ref{lem:local_time_full_proof}, establishes that the target process $W_s = G(s,X_s,Y_s,Z_s)$ almost surely avoids the cut locus of $\Ucal$. This implies that the projection set is a singleton for $P_0 \otimes dt$-a.e. $(s,\omega)$.
Consequently, the selection process $\Theta$ is uniquely determined as a function of the path $(X_s, Y_s, Z_s)$. Let this unique maximizer map be denoted by $\Thetacal^*(s,x,y,z)$.

It follows that both solution pairs \((Y^1, Z^1)\) and \((Y^2, Z^2)\) must satisfy the same standard BSDE:
\begin{equation} \label{eq:proof_uniqueness_effective_bsde}
    Y_t = \phi(X_T) + \int_t^T \tilde{F}(s, X_s, Y_s, Z_s) \dd s - \int_t^T Z_s \dd B_s,
\end{equation}
where the single-valued effective driver \(\tilde{F}\) is given by
\[ \tilde{F}(t,x,y,z) \coloneqq h(t,x,y,z) - \frac{1}{2}\dist(G(t,x,y,z), \Ucal)^2. \]
Under Assumption \ref{ass:pathwise_consistency}, the functions $h$ and $G$ are globally Lipschitz in their state variables, and the squared distance function is locally Lipschitz. Therefore, the effective driver $\tilde{F}(t,x,y,z)$ is locally Lipschitz in $(y,z)$, uniformly in $(t,x)$. Our task thus reduces to proving that the BSDE \eqref{eq:proof_uniqueness_effective_bsde}, with a locally Lipschitz driver, admits a unique solution in the space $\Bcal^2([0,T])$.

\vspace{1ex}
\noindent\textbf{Step 3.2: Uniqueness of the Value Process $Y$ via Localization.}
Let \(\delta Y_t \coloneqq Y^1_t - Y^2_t\) and \(\delta Z_t \coloneqq Z^1_t - Z^2_t\). Subtracting the BSDE for $(Y^2, Z^2)$ from that for $(Y^1, Z^1)$ yields:
\begin{equation} \label{eq:proof_deltaY_bsde}
    \delta Y_t = \int_t^T \Delta \tilde{F}_s \dd s - \int_t^T \delta Z_s \dd B_s,
\end{equation}
where $\Delta \tilde{F}_s \coloneqq \tilde{F}(s, X_s, Y^1_s, Z^1_s) - \tilde{F}(s, X_s, Y^2_s, Z^2_s)$.

The local Lipschitz property of $\tilde{F}$ prevents a direct global estimate. We therefore introduce a sequence of stopping times to localize the argument. For each integer \(R \ge 1\), define the stopping time
\[
    \tau_R \coloneqq \inf \left\{ s \in [0,T] \mid \norm{Y^1_s} > R \text{ or } \norm{Z^1_s}_F > R \text{ or } \norm{Y^2_s} > R \text{ or } \norm{Z^2_s}_F > R \right\} \wedge T.
\]
Since \(Y^1, Y^2 \in \Bcal^2([0,T])\) and \(Z^1, Z^2 \in \Hcalp{2}([0,T])\), the processes are finite almost surely on $[0,T]$. Thus, the sequence of stopping times converges to \(T\) almost surely: \(\lim_{R\to\infty} \tau_R = T\), $P_0$-a.s.

Let \(L_R\) be the Lipschitz constant of \(\tilde{F}(t,x, \cdot, \cdot)\) on the compact ball in $\R \times \R^{1\times d}$ of radius $R$. For any \(s \in [0, \tau_R]\), the states \((Y^i_s, Z^i_s)\) are bounded in norm by $R$, allowing us to apply the local Lipschitz property:
\begin{equation}\label{eq:local_lipschitz_bound}
    \abs{\Delta \tilde{F}_s} \le L_R \left( \abs{\delta Y_s} + \norm{\delta Z_s}_F \right).
\end{equation}
We apply Itô's formula to the process $\abs{\delta Y_t}^2$ on the interval $[t,T]$. For any $t \in [0,T]$, we have
\begin{equation} \label{eq:proof_ito_deltaY}
\begin{split}
    \abs{\delta Y_t}^2 + \int_t^T \norm{\delta Z_s}_F^2 \dd s = 2 \int_t^T \delta Y_s \Delta \tilde{F}_s \dd s - 2 \int_t^T \delta Y_s \delta Z_s \dd B_s.
\end{split}
\end{equation}
Applying the stopping time $\tau_R$ to this identity and taking expectations, we obtain for any $t \in [0,T]$:
\begin{align*}
    \mathbb{E}\left[\abs{\delta Y_{t \wedge \tau_R}}^2\right] &+ \mathbb{E}\left[\int_{t \wedge \tau_R}^{\tau_R} \norm{\delta Z_s}_F^2 \dd s\right] = 2 \, \mathbb{E}\left[\int_{t \wedge \tau_R}^{\tau_R} \delta Y_s \Delta \tilde{F}_s \dd s\right].
\end{align*}
The expectation of the stochastic integral term is zero, as it is a true martingale up to time $\tau_R$. We now bound the right-hand side. For $s \in [t \wedge \tau_R, \tau_R]$, the bound \eqref{eq:local_lipschitz_bound} holds. Using this and Young's inequality ($2ab \le a^2/\varepsilon + \varepsilon b^2$ for any $\varepsilon>0$):
\begin{align*}
    2 \delta Y_s \Delta \tilde{F}_s &\le 2 \abs{\delta Y_s} \abs{\Delta \tilde{F}_s} \\
    &\le 2 \abs{\delta Y_s} L_R (\abs{\delta Y_s} + \norm{\delta Z_s}_F) \\
    &= 2L_R \abs{\delta Y_s}^2 + 2L_R \abs{\delta Y_s} \norm{\delta Z_s}_F \\
    &\le 2L_R \abs{\delta Y_s}^2 + (L_R^2 \abs{\delta Y_s}^2 + \norm{\delta Z_s}_F^2) \quad (\text{with } \varepsilon=1)\\
    &= (2L_R + L_R^2) \abs{\delta Y_s}^2 + \norm{\delta Z_s}_F^2.
\end{align*}
Let \(C_R \coloneqq 2L_R + L_R^2\). Substituting this bound back into the expectation gives:
\begin{align*}
    \mathbb{E}\left[\abs{\delta Y_{t \wedge \tau_R}}^2\right] &+ \mathbb{E}\left[\int_{t \wedge \tau_R}^{\tau_R} \norm{\delta Z_s}_F^2 \dd s\right] \\
    &\le \mathbb{E}\left[\int_{t \wedge \tau_R}^{\tau_R} \left( C_R \abs{\delta Y_s}^2 + \norm{\delta Z_s}_F^2 \right) \dd s\right] \\
    &= C_R \mathbb{E}\left[\int_{t \wedge \tau_R}^{\tau_R} \abs{\delta Y_s}^2 \dd s\right] + \mathbb{E}\left[\int_{t \wedge \tau_R}^{\tau_R} \norm{\delta Z_s}_F^2 \dd s\right].
\end{align*}
The term $\mathbb{E}\left[\int_{t \wedge \tau_R}^{\tau_R} \norm{\delta Z_s}_F^2 \dd s\right]$ is finite and can be subtracted from both sides, yielding:
\[
    \mathbb{E}\left[\abs{\delta Y_{t \wedge \tau_R}}^2\right] \le C_R \mathbb{E}\left[\int_{t \wedge \tau_R}^{\tau_R} \abs{\delta Y_s}^2 \dd s\right] \le C_R \int_t^T \mathbb{E}\left[\abs{\delta Y_{s \wedge \tau_R}}^2\right] \dd s.
\]
Let \(v_R(t) \coloneqq \mathbb{E}\left[\abs{\delta Y_{t \wedge \tau_R}}^2\right]\). The function $v_R(t)$ is non-negative, continuous, and satisfies the integral inequality \(v_R(t) \le C_R \int_t^T v_R(s) \dd s\). By the standard Gronwall's inequality for integral forms, this implies that \(v_R(t) = 0\) for all \(t \in [0,T]\).

We have thus shown that for every \(R \ge 1\) and every $t \in [0,T]$, \(\mathbb{E}\left[\abs{\delta Y_{t \wedge \tau_R}}^2\right] = 0\), which implies $\delta Y_{t \wedge \tau_R} = 0$ a.s. As \(R \to \infty\), we have \(t \wedge \tau_R \to t\) almost surely. Since the paths of $Y^1$ and $Y^2$ are continuous, this means $\delta Y_{t \wedge \tau_R}(\omega) \to \delta Y_t(\omega)$ for every $\omega$. Furthermore, since $Y^1, Y^2 \in \Bcal^2([0,T])$, the random variables $\abs{\delta Y_{t \wedge \tau_R}}^2$ are uniformly bounded by $(2\sup_s|Y^1_s| + 2\sup_s|Y^2_s|)^2$, which is integrable. By the dominated convergence theorem, we can take the limit inside the expectation:
\[
    \mathbb{E}[\abs{\delta Y_t}^2] = \mathbb{E}\left[\lim_{R \to \infty} \abs{\delta Y_{t \wedge \tau_R}}^2\right] = \lim_{R \to \infty} \mathbb{E}\left[\abs{\delta Y_{t \wedge \tau_R}}^2\right] = 0.
\]
This holds for all \(t \in [0,T]\). Since $\delta Y$ is a continuous process, this implies that $P_0(\delta Y_t = 0, \forall t \in [0,T]) = 1$. The processes $Y^1$ and $Y^2$ are indistinguishable.

\vspace{1ex}
\noindent\textbf{Step 3.3: Uniqueness of \(Z\) and \(\Theta\).}
With \(Y^1 \equiv Y^2\), we have \(\delta Y \equiv 0\). Returning to the Itô identity \eqref{eq:proof_ito_deltaY} and setting $\delta Y \equiv 0$, the equation simplifies dramatically:
\[
    \int_t^T \norm{\delta Z_s}_F^2 \dd s = 0, \quad P_0\text{-a.s., for all } t \in [0,T].
\]
In particular, for $t=0$, we have $\int_0^T \norm{\delta Z_s}_F^2 \dd s = 0$, $P_0$-a.s. This implies that \(\mathbb{E}[\int_0^T \norm{\delta Z_s}_F^2 \dd s] = 0\). Since the integrand is non-negative, it must be that \(\delta Z_s = 0\) for $P_0 \otimes dt$-a.e. $(s,\omega)$. The processes $Z^1$ and $Z^2$ are therefore indistinguishable.

Finally, the uniqueness of the maximizer process \(\Theta\) follows directly. For $P_0 \otimes dt$-a.e. $(s,\omega)$, we have:
\[ \Theta^1_s = \Thetacal^*(s,X_s,Y^1_s,Z^1_s) \quad \text{and} \quad \Theta^2_s = \Thetacal^*(s,X_s,Y^2_s,Z^2_s). \]
Since $Y^1 \equiv Y^2$ and $Z^1 \equiv Z^2$, it follows immediately that $\Theta^1 \equiv \Theta^2$.

Thus, any two solutions must coincide, establishing that the solution triplet \((Y, Z, \Theta)\) is unique. This concludes the proof of uniqueness and of Theorem \ref{thm:existence_geometric}.

\end{proof}

\section{Axiomatic Properties of the \texorpdfstring{$\theta$}{theta}-Expectation Operator}\label{sec:axiomatic_space}


Having established well-posedness for the $\theta$-BSDE under both the Global Regularity Hypothesis (\Cref{thm:theta_bsde_wellposedness}) and the more general Geometric Regularity Hypothesis (\Cref{thm:existence_geometric}), we are now positioned to define the associated valuation operator and rigorously establish its fundamental properties. The existence of a unique solution triplet $(Y, Z, \Theta)$ for any square-integrable terminal condition $\xi$ is the cornerstone of this construction. The objective is to show that it yields a dynamically consistent functional that retains monotonicity but, by design, does not necessarily satisfy sub-additivity. We also establish a sharp structural condition on the primitive driver $F$ that is equivalent to translation-invariance. This departure from the canonical axioms of sublinear expectation is what enables sensitivity to non-convex structures.


\begin{definition}[The $\theta$-Expectation Space]
\label{def:theta_expectation_space}
A \emph{$\theta$-expectation space} is a triple $(\Omega, \mathcal{H}, (\Ecal_t)_{t\in[0,T]})$ where $\mathcal{H} \coloneqq L^2(\Omega, \mathcal{F}_T, P_0)$ and, for each $t \in [0,T]$, $\Ecal_t: \mathcal{H} \to L^2(\Omega, \mathcal{F}_t, P_0)$ is a functional satisfying:
\begin{enumerate}[label=(A\arabic*)]
    \item \textbf{(Monotonicity)} If $X \ge Y$ $P_0$-a.s., then $\Ecal_t[X] \ge \Ecal_t[Y]$ $P_0$-a.s.
    \item \textbf{(Translation-Invariance)} For any $m \in L^2(\Omega, \mathcal{F}_t, P_0)$, $\Ecal_t[X + m] = \Ecal_t[X] + m$.
    \item \textbf{(Dynamic Consistency / Tower Property)} For any $0 \le t \le s \le T$, $\Ecal_t[X] = \Ecal_t\left[\Ecal_s[X]\right]$.
\end{enumerate}
\end{definition}

\begin{remark}[The Absence of Sub-additivity and Convexity]
\label{rem:absence_of_subadditivity}
A defining characteristic of this system is the deliberate omission of sub-additivity ($\Ecal_t[X+Y] \le \Ecal_t[X] + \Ecal_t[Y]$). This axiom is mathematically equivalent to the convexity of the set of priors in the dual representation of the expectation. By abandoning it, we sever the intrinsic link to convexity that limits classical theories. The central mathematical challenge, which this paper resolves for two structurally regular classes of models, is to find a tractable mathematical structure to substitute for the global convexity assumption.
\end{remark}

\begin{definition}[The Generated \(\theta\)-Expectation]
\label{def:generated_theta_expectation}
Let the primitive driver $F$ and uncertainty set $\Ucal$ satisfy the conditions of either \Cref{thm:theta_bsde_wellposedness} or \Cref{thm:existence_geometric}. We define the \emph{conditional $\theta$-expectation generated by $(F, \Ucal)$} as the family of operators $(\Efrakgen_t)_{t \in [0,T]}$, where for each $t$, $\Efrakgen_t: L^2(\Omega, \mathcal{F}_T, P_0) \to L^2(\Omega, \mathcal{F}_t, P_0)$ is given by
\[ \Efrakgen_t[\xi] \coloneqq Y_t, \]
where $(Y, Z, \Theta)$ is the unique solution in $\Bcal^2([0,T]) \times \Lcalp{2}([0,T]; \Sd)$ to the $\theta$-BSDE system (\cref{def:theta_bsde_system}) with terminal condition $\xi$.
\end{definition}

The following theorem demonstrates that this construction yields a coherent valuation functional that, by design, departs from the classical axiom of sub-additivity.

\begin{theorem}[Fundamental Properties of the \(\theta\)-Expectation]\label{thm:properties_of_operator}
Let the generator $(F, \Ucal)$ satisfy the conditions for well-posedness. The family of operators $(\Efrakgen_t)_{t \in [0,T]}$ satisfies the following properties:
\begin{enumerate}[label=(\roman*)]
    \item \textbf{(Monotonicity)} For any $\xi_1, \xi_2 \in L^2(\mathcal{F}_T)$ such that $\xi_1 \ge \xi_2$ $P_0$-a.s., we have $\Efrakgen_t[\xi_1] \ge \Efrakgen_t[\xi_2]$ $P_0$-a.s. for all $t \in [0,T]$.
    \item \textbf{(Dynamic Consistency)} For any $\xi \in L^2(\mathcal{F}_T)$ and any $0 \le t \le s \le T$, we have $\Efrakgen_t[\xi] = \Efrakgen_t\left[\Efrakgen_s[\xi]\right]$ $P_0$-a.s.
    \item \textbf{(Translation-Invariance)} The operator is translation-invariant, i.e., for any $t \in [0,T]$ and any $m \in L^2(\Omega, \mathcal{F}_t, P_0)$, we have $\Efrakgen_t[\xi + m] = \Efrakgen_t[\xi] + m$ $P_0$-a.s., if and only if the primitive driver $F(t,\omega,y,z,\theta)$ does not depend on its argument $y$.
\end{enumerate}
\end{theorem}

\begin{proof}
The proofs rely on the unified structure established in \Cref{sec:wellposedness_global_regularity} and \Cref{sec:main_result_geometric}. In both well-posed cases, the $\theta$-BSDE is equivalent to a standard BSDE with a single-valued effective driver $\tilde{F}(t,\omega,y,z) \coloneqq \sup_{\theta \in \Ucal} F(t,\omega,y,z,\theta)$. This driver is globally Lipschitz in $(y,z)$ under Hypothesis \ref{hyp:global_regularity} and locally Lipschitz in $(y,z)$ under Assumption \ref{ass:pathwise_consistency}.

\vspace{1ex}
\noindent\textbf{(i) Proof of Monotonicity.}
Let $\xi_1, \xi_2 \in L^2(\mathcal{F}_T)$ be such that $\xi_1 \ge \xi_2$ $P_0$-a.s. Let $(Y^1, Z^1)$ and $(Y^2, Z^2)$ be the unique solution pairs to the effective BSDE with terminal conditions $\xi_1$ and $\xi_2$ respectively. Both pairs are driven by the same effective driver $\tilde{F}$. The comparison principle for BSDEs states that if the terminal condition is larger and the driver is the same, the solution is ordered accordingly. Specifically:
\begin{itemize}
    \item Under Hypothesis \ref{hyp:global_regularity}, $\tilde{F}$ is globally Lipschitz. The standard comparison principle from \cite{ElKaroui1997} applies directly, yielding $Y^1_t \ge Y^2_t$ for all $t \in [0,T]$, $P_0$-a.s.
    \item Under Assumption \ref{ass:pathwise_consistency}, the effective driver $\tilde{F}$ is only locally Lipschitz in $(y,z)$ but has quadratic structure. The comparison principle for such BSDEs, as established in \cite{Kobylanski2000}, holds. Thus, we again conclude that $Y^1_t \ge Y^2_t$ for all $t \in [0,T]$, $P_0$-a.s.
\end{itemize}
By Definition \ref{def:generated_theta_expectation}, $\Efrakgen_t[\xi_1] = Y^1_t$ and $\Efrakgen_t[\xi_2] = Y^2_t$. The result follows immediately.

\vspace{1ex}
\noindent\textbf{(ii) Proof of Dynamic Consistency.}
Let $\xi \in L^2(\mathcal{F}_T)$ and $0 \le t \le s \le T$. Let $(Y,Z)$ be the unique solution to the effective BSDE on $[0,T]$ with terminal condition $\xi$. By definition, we have:
\begin{equation*}
    \Efrakgen_t[\xi] = Y_t \quad \text{and} \quad \Efrakgen_s[\xi] = Y_s.
\end{equation*}
To evaluate the right-hand side of the tower property, we must compute $\Efrakgen_t[\Efrakgen_s[\xi]] = \Efrakgen_t[Y_s]$. This requires solving a $\theta$-BSDE on the time interval $[0,s]$ with terminal condition $Y_s$. By the uniqueness of the BSDE solution on any time interval, the process $(Y_u)_{u \in [0,s]}$ is itself the unique solution on $[0,s]$ to the BSDE
\begin{equation*}
    Y_u = Y_s + \int_u^s \tilde{F}(r, \omega, Y_r, Z_r) \dd r - \int_u^s Z_r \dd B_r, \quad u \in [0,s].
\end{equation*}
By Definition \ref{def:generated_theta_expectation}, the value of this solution at time $t$ is precisely $\Efrakgen_t[Y_s]$. Therefore,
\begin{equation*}
    \Efrakgen_t[\Efrakgen_s[\xi]] = \Efrakgen_t[Y_s] = Y_t.
\end{equation*}
Since $Y_t = \Efrakgen_t[\xi]$, the identity $\Efrakgen_t[\xi] = \Efrakgen_t[\Efrakgen_s[\xi]]$ is established.

\vspace{1ex}
\noindent\textbf{(iii) Proof of Translation-Invariance.}
This is an equivalence, which requires two separate proofs.

\paragraph{\textbf{Sufficiency ($\Leftarrow$)}} Assume the primitive driver $F(t,\omega,y,z,\theta)$ does not depend on $y$. Then the effective driver $\tilde{F}(t,\omega,y,z) = \sup_{\theta \in \Ucal} F(t,\omega,z,\theta)$ is also independent of $y$, and we denote it by $\tilde{F}(t,\omega,z)$. Let $\xi \in L^2(\mathcal{F}_T)$ and $m \in L^2(\mathcal{F}_t)$. Let $(Y,Z)$ be the unique solution corresponding to the terminal condition $\xi$. We claim that the pair $(\hat{Y}, \hat{Z})$ defined by $\hat{Y}_s \coloneqq Y_s + m$ and $\hat{Z}_s \coloneqq Z_s$ for $s \in [t,T]$, with $\hat{Y}_s$ defined by the backward recursion for $s<t$, is the unique solution for the terminal condition $\xi+m$.

We verify this claim on the interval $[t,T]$. The terminal condition is met: $\hat{Y}_T = Y_T + m = \xi + m$. The pair $(\hat{Y}, \hat{Z})$ must satisfy the BSDE:
\begin{align*}
    \hat{Y}_t &= \hat{Y}_T + \int_t^T \tilde{F}(s, \omega, \hat{Z}_s) \dd s - \int_t^T \hat{Z}_s \dd B_s \\
    Y_t + m &= (\xi + m) + \int_t^T \tilde{F}(s, \omega, Z_s) \dd s - \int_t^T Z_s \dd B_s.
\end{align*}
Subtracting $m$ from both sides yields the original BSDE for $(Y,Z)$, which holds by definition. By uniqueness of the solution on $[t,T]$, the solution for $\xi+m$ must be $(\hat{Y}, \hat{Z})$. Therefore,
\begin{equation*}
    \Efrakgen_t[\xi+m] = \hat{Y}_t = Y_t + m = \Efrakgen_t[\xi] + m.
\end{equation*}
The property holds.

\paragraph{\textbf{Necessity ($\Rightarrow$)}} Assume translation-invariance holds. Let $c \in \R$ be any constant. By assumption, for any $\xi \in L^2(\mathcal{F}_T)$, if $(Y,Z)$ is the solution for $\xi$, then $(Y+c, Z)$ must be the solution for $\xi+c$. The two pairs satisfy their respective effective BSDEs:
\begin{align*}
    Y_t &= \xi + \int_t^T \tilde{F}(s, Y_s, Z_s) \dd s - \int_t^T Z_s \dd B_s \\
    Y_t+c &= (\xi+c) + \int_t^T \tilde{F}(s, Y_s+c, Z_s) \dd s - \int_t^T Z_s \dd B_s.
\end{align*}
Subtracting the second equation from the first (after canceling terms) yields
\begin{equation*}
    \int_t^T \tilde{F}(s, Y_s, Z_s) \dd s = \int_t^T \tilde{F}(s, Y_s+c, Z_s) \dd s, \quad \forall t \in [0,T].
\end{equation*}
Differentiating with respect to $t$ implies that for any solution process $(Y,Z)$, we must have
\begin{equation} \label{eq:proof_driver_identity}
    \tilde{F}(s, Y_s(\omega), Z_s(\omega)) = \tilde{F}(s, Y_s(\omega)+c, Z_s(\omega))
\end{equation}
for $P_0 \otimes dt$-a.e. $(s,\omega)$ and for all $c \in \R$.

Under the non-degeneracy assumptions of our main theorem (\Cref{thm:existence_geometric}), the law of the solution vector $(s, Y_s, Z_s)$ is absolutely continuous with respect to the Lebesgue measure on $(0,T] \times \R \times \R^{1\times d}$. Since $\tilde{F}$ is continuous in $(y,z)$, the identity \eqref{eq:proof_driver_identity}, which holds on a set of full measure in the space of path-realizations, implies that the function $\tilde{F}(s,\omega,y,z)$ must be independent of $y$.

Finally, we argue that the independence of $\tilde{F}$ from $y$ implies the independence of $F$ from $y$. Suppose $F$ depends on $y$. Then there exists $(t,\omega,y_0,z_0,\theta_0)$ such that $\partial_y F(t,\omega,y_0,z_0,\theta_0) \neq 0$. By continuity, this holds in a neighborhood. We can choose the problem primitives such that this neighborhood is visited with positive probability and that $\theta_0$ is the unique maximizer. In this case, by the envelope theorem, $\partial_y \tilde{F}(t,\omega,y_0,z_0) = \partial_y F(t,\omega,y_0,z_0,\theta_0) \neq 0$, which contradicts the finding that $\tilde{F}$ is independent of $y$. Thus, $F$ must be independent of $y$ for the property to hold universally.
\end{proof}

\begin{remark}[The Structural Implication of Translation-Invariance]
\Cref{thm:properties_of_operator}(iii) reveals a fundamental structural feature of our framework. Translation-invariance, often interpreted as insensitivity to the addition of deterministic cash, is not an automatic property but a direct consequence of the structure of the primitive driver $F$. Any model in which the perception of or reaction to ambiguity, as encoded in $F$, depends on the current value or wealth level of the agent ($Y_t$) will necessarily generate a valuation functional that is not translation-invariant. This provides a clear and interpretable modeling trade-off between axiomatic properties and the desire to represent state-dependent ambiguity aversion.
\end{remark}

\section{\texorpdfstring{$\theta$}{theta}-Martingale Theory and Calculus}
\label{sec:martingales_revised}

The classical theory of martingales is fundamentally tethered to a single, linear conditional expectation operator. Having established the well-posedness of the \(\theta\)-BSDE system under two significant model classes, we are now positioned to build a coherent stochastic calculus that is sensitive to non-convex ambiguity.

This section develops the necessary tools. We begin by defining the foundational class of processes for our calculus, the \emph{\(\theta\)-semimartingales}, whose existence and uniqueness are guaranteed by our main theorems. We then introduce the central analytic tool for this calculus: the path-dependent \(\theta\)-generator. This allows us to state and prove a genuine \(\theta\)-Itô formula, which serves as the proper analogue to the G-Itô formula and makes the structure of our non-linear calculus explicit. We then introduce the concept of endogenous quadratic variation to quantify realized ambiguity and provide a canonical representation theorem that connects our framework to the theory of Reflected BSDEs, providing a deeper characterization of the martingale structure. Finally, we establish a $\theta$-Girsanov theorem, a critical component for a complete theory of pricing and hedging.

Throughout this section, we operate under the standing assumption that the problem primitives \((F, \Ucal)\) satisfy the conditions for well-posedness established in either Theorem \ref{thm:theta_bsde_wellposedness} (Global Regularity) or Theorem \ref{thm:existence_geometric} (Geometric Regularity). This is a crucial prerequisite, as it guarantees the uniqueness of the decompositions and operators we define, rendering the subsequent calculus unambiguous.

\subsection{\texorpdfstring{$\theta$}{theta}-Semimartingales and the Unique Decomposition}
The foundational class of processes for our calculus are the \(\theta\)-semimartingales, which are defined via a unique decomposition that directly leverages the well-posedness of the \(\theta\)-BSDE system.

\begin{definition}[$\theta$-Semimartingale]
\label{def:theta_semimartingale}
An adapted process \(M \in \Scalp{2}([0,T])\) is a \emph{\(\theta\)-semimartingale} if it admits a unique decomposition of the form:
\begin{equation}\label{eq:semimartingale_decomposition}
    M_t = M_0 + \int_0^t \mu_s ds + \int_0^t Z_s dB_s,
\end{equation}
where \(Z \in \Hcalp{2}([0,T])\) and the drift rate process \(\mu \in \Lcalp{2}([0,T])\) is given by
\begin{equation} \label{eq:drift_decomposition}
    \mu_t = \alpha_t - F(t, \omega, M_t, Z_t, \Theta_t).
\end{equation}
Here, \(\alpha\) is a progressively measurable process in \(\Lcalp{2}([0,T])\), and \(\Theta\) is the almost surely unique admissible selection process for the pair \((M_t,Z_t)\), i.e., for \(P_0 \otimes dt\) almost every \((t,\omega)\):
\[
    \Theta_t(\omega) = \underset{\theta \in \Ucal}{\argmax} \; F(t, \omega, M_t(\omega), Z_t(\omega), \theta).
\]
The process \(\alpha\) is called the \emph{excess drift} of \(M\).
\end{definition}

\begin{remark}[Uniqueness of the Decomposition]
The uniqueness of the triplet \((M_0, \alpha, Z)\) for a given process \(M\) is a direct and critical consequence of the main theorems of this paper. For any given semimartingale \(M\), the martingale representation theorem guarantees the uniqueness of its decomposition into a finite-variation part and a stochastic integral, thereby uniquely identifying the drift process \(\mu\) and the volatility process \(Z\). The pathwise consistency results of Theorem \ref{thm:theta_bsde_wellposedness} and Theorem \ref{thm:existence_geometric} then guarantee that for the uniquely determined path \((M_t(\omega), Z_t(\omega))\), the maximizer \(\Theta_t(\omega)\) is also uniquely determined for \(P_0 \otimes dt\)-a.e. \((t,\omega)\). Consequently, the endogenous drift term \(F(t, \omega, M_t, Z_t, \Theta_t)\) is uniquely determined. The excess drift \(\alpha\) is then uniquely identified as the residual: \(\alpha_t = \mu_t + F(t,\omega,M_t,Z_t,\Theta_t)\). This rigorous chain of identification is what makes the calculus well-defined.
\end{remark}

This definition provides a clear hierarchy based on the structure of the excess drift, leading to natural definitions for martingales and supermartingales.

\begin{definition}[$\theta$-Martingales and $\theta$-Supermartingales]
Let \(M\) be a \(\theta\)-semimartingale with the unique decomposition described in Definition \ref{def:theta_semimartingale}.
\begin{enumerate}
    \item \(M\) is a \emph{\(\theta\)-martingale} if its excess drift is zero, i.e., \(\alpha_t \equiv 0\). Its dynamics are given by \(dM_t = -F(t,M_t,Z_t,\Theta_t)dt + Z_t dB_t\).
    \item \(M\) is a \emph{\(\theta\)-supermartingale} if its excess drift is non-positive, i.e., \(\alpha_t \le 0\) a.e.
    \item \(M\) is a \emph{\(\theta\)-submartingale} if its excess drift is non-negative, i.e., \(\alpha_t \ge 0\) a.e.
\end{enumerate}
\end{definition}

\subsection{The \texorpdfstring{$\theta$}{theta}-Generator and Itô's Formula}
The classical Itô formula, when applied to a \(\theta\)-semimartingale, reveals a complex drift structure. By abstracting this structure into a new operator—the \(\theta\)-generator—we can formulate a true \(\theta\)-Itô formula that makes the non-linear nature of the calculus explicit.

\begin{definition}[The $\theta$-Generator]
\label{def:theta_generator}
Let \((M, Z, \Theta)\) be the unique triplet defining a \(\theta\)-semimartingale with excess drift $\alpha$. The \emph{\(\theta\)-generator} \(\mathcal{G}_t\) is a path-dependent operator that acts on a function \(\phi \in C^{1,2}([0,T] \times \R^k; \R)\) at time \(t\) along the solution path:
\begin{align*}
    \mathcal{G}_t(\phi) \coloneqq \partial_t \phi(t,M_t) &+ \scpr{\nabla_x \phi(t,M_t)}{\alpha_t - F(t, M_t, Z_t, \Theta_t)} \\
    &+ \frac{1}{2}\Tr\left(D^2_x \phi(t,M_t) Z_t Z_t^T \right).
\end{align*}
If $M$ is a \(\theta\)-martingale, the excess drift $\alpha_t=0$, and the generator simplifies to
\begin{align*}
    \mathcal{G}_t(\phi) \coloneqq \partial_t \phi(t,M_t) &- \scpr{\nabla_x \phi(t,M_t)}{F(t, M_t, Z_t, \Theta_t)} \\
    &+ \frac{1}{2}\Tr\left(D^2_x \phi(t,M_t) Z_t Z_t^T \right).
\end{align*}
\end{definition}

This operator represents the instantaneous expected rate of change of \(\phi(t, M_t)\) under the endogenously determined worst-case measure. With it, we can state the main result of our calculus.

\begin{theorem}[The $\theta$-Itô Formula]
\label{thm:theta_ito_formula}
Let \(M\) be a \(k\)-dimensional \(\theta\)-semimartingale. For any function \(f \in C^{1,2}([0,T] \times \R^k; \R)\), the process \(X_t \coloneqq f(t, M_t)\) is a one-dimensional Itô process whose dynamics are given by:
\begin{equation}
    df(t, M_t) = \mathcal{G}_t(f) dt + \nabla_x f(t,M_t) Z_t dB_t.
\end{equation}
\end{theorem}

\begin{proof}
By Definition \ref{def:theta_semimartingale}, the process \(M\) is a classical Itô semimartingale with drift rate \(\mu_t = \alpha_t - F(t, M_t, Z_t, \Theta_t)\) and diffusion coefficient \(Z_t\). Its quadratic variation process is \(d[M]_t = Z_t Z_t^T dt\).

We apply the standard Itô formula for semimartingales (see, e.g., \cite[Chapter IV, Theorem 1.5]{RevuzYor1999}) to the process \(X_t = f(t, M_t)\):
\begin{align*}
    f(t, M_t) - f(0, M_0) = & \int_0^t \partial_s f(s, M_s) ds + \int_0^t (\nabla_x f(s, M_s))^T dM_s \\
    & + \frac{1}{2} \int_0^t \Tr\left( D^2_x f(s, M_s) d[M]_s \right).
\end{align*}
Substituting the differentials \(dM_s = \mu_s ds + Z_s dB_s\) and \(d[M]_s = Z_s Z_s^T ds\), we obtain:
\begin{align*}
    f(t, M_t) - f(0, M_0) = & \int_0^t \partial_s f(s, M_s) ds + \int_0^t \scpr{\nabla_x f(s, M_s)}{\mu_s} ds + \int_0^t \nabla_x f(s, M_s) Z_s dB_s \\
    & + \frac{1}{2} \int_0^t \Tr\left( D^2_x f(s, M_s) Z_s Z_s^T \right) ds.
\end{align*}
We collect all terms under the Lebesgue integral, which constitute the drift of the process $f(t, M_t)$:
\[ \text{Drift}_t = \partial_t f(t, M_t) + \scpr{\nabla_x f(t, M_t)}{\mu_t} + \frac{1}{2}\Tr\left( D^2_x f(t, M_t) Z_t Z_t^T \right). \]
Substituting the expression for the drift \(\mu_t = \alpha_t - F(t, M_t, Z_t, \Theta_t)\) gives:
\[ \text{Drift}_t = \partial_t f(t, M_t) + \scpr{\nabla_x f(t, M_t)}{\alpha_t - F(t, M_t, Z_t, \Theta_t)} + \frac{1}{2}\Tr\left( D^2_x f(t, M_t) Z_t Z_t^T \right). \]
This is precisely the definition of the \(\theta\)-generator \(\mathcal{G}_t(f)\) from Definition \ref{def:theta_generator}. The differential form of the process is therefore
\[ df(t, M_t) = \mathcal{G}_t(f) dt + \nabla_x f(t, M_t) Z_t dB_t. \]
This completes the proof.
\end{proof}

\begin{remark}[Comparison with G-Itô's Formula]
The \(\theta\)-Itô formula reveals the fundamental structural difference between our framework and Peng's G-calculus.
\begin{itemize}
    \item \textbf{G-Itô Formula:} The generator \(\mathcal{L}^G(\phi) = G(D^2\phi)\) is a deterministic, state-independent function of the derivatives of \(\phi\). Ambiguity is encapsulated entirely in the second-order term via the static, convexifying operator $G(\cdot) = \frac{1}{2}\sup_{\theta \in \overlineconv(\Ucal)} \Tr(\theta \cdot)$.
    \item \textbf{$\theta$-Itô Formula:} The generator \(\mathcal{G}_t(\phi)\) is a random, path-dependent operator. It depends not only on the derivatives of \(\phi\) but also on the current state of the entire system \((M_t, Z_t, \Theta_t)\). Ambiguity affects both first- and second-order terms in a coupled, dynamic fashion. This path-dependence is precisely what allows the calculus to be sensitive to the non-convex geometry of the primitive set \(\Ucal\).
\end{itemize}
\end{remark}

\subsection{Endogenous Quadratic Variation}
In our framework, we can explicitly quantify the realized ambiguity.

\begin{definition}[Endogenous Quadratic Variation]
\label{def:endogenous_qv}
Let \(M\) be a \(\theta\)-semimartingale with associated process \(\Theta\). The \emph{endogenous quadratic variation} of \(M\), denoted \(\langle\langle M \rangle\rangle\), is the process defined by
\begin{equation}
    \langle\langle M \rangle\rangle_t \coloneqq \int_0^t \Theta_s ds.
\end{equation}
This is a continuous, adapted process with finite variation, taking values in the cone of positive semidefinite matrices \(\Sdpn\).
\end{definition}

\begin{remark}[Conceptual Distinction]
\label{rem:qv_distinction}
It is essential to distinguish between the two notions of quadratic variation present in this theory.
\begin{enumerate}
    \item \textbf{Classical Quadratic Variation, \([M]_t = \int_0^t Z_s Z_s^T ds\):} This measures the realized volatility of the \emph{valuation process} \(M_t\). It quantifies risk \emph{within} a given model scenario.
    \item \textbf{Endogenous Quadratic Variation, \(\langle\langle M \rangle\rangle_t\):} This measures the realized volatility of the \emph{underlying model} itself. It is the cumulative measure of the specific model uncertainty (parameterized by \(\Theta_s\)) that is deemed most adversarial at each instant. It is a measure of realized \emph{ambiguity} across models.
\end{enumerate}
\end{remark}

\subsection{A Canonical Representation Theorem}
We now connect our definitions to the analytical theory of BSDEs. The following theorem establishes that our definition of a \(\theta\)-supermartingale is equivalent to the classical notion of a supersolution to a Reflected BSDE.

\begin{theorem}[Canonical Representation of \(\theta\)-Supermartingales]
\label{thm:canonical_decomposition}
Let the problem primitives \((F, \Ucal)\) satisfy the baseline conditions of Assumption \ref{ass:bsdi_primitives}. Let \(V(t,\omega,y,z) \coloneqq \sup_{\theta \in \Ucal} F(t,\omega,y,z,\theta)\) be the associated upper-bound driver.
\begin{enumerate}[label=(\roman*)]
    \item \textbf{(Equivalence)} An adapted process \(M \in \Scalp{2}([0,T])\) is a \(\theta\)-supermartingale if and only if there exists a pair \((Z, A)\) such that the triplet \((M, Z, A)\) solves the Reflected BSDE (RBSDE):
    \begin{equation}\label{eq:rbsde_representation}
        M_t = M_T + \int_t^T V(s, M_s, Z_s) \dd s + (A_T - A_t) - \int_t^T Z_s \dd B_s,
    \end{equation}
    where \(Z \in \Hcalp{2}([0,T])\) and \(A\) is a continuous, non-decreasing adapted process with \(A_0=0\).

    \item \textbf{(Uniqueness)} If, in addition, the primitives satisfy the well-posedness conditions of either Theorem \ref{thm:theta_bsde_wellposedness} or \ref{thm:existence_geometric}, then for any \(\theta\)-supermartingale \(M\), its unique decomposition \eqref{eq:semimartingale_decomposition} is characterized by the unique solution triplet \((M, Z, A)\) of the RBSDE \eqref{eq:rbsde_representation}. The excess drift is given by \(\alpha_t = -a_t\), where \(a_t\) is the time-derivative of the reflection process \(A_t\), and the admissible selection process \(\Theta\) is unique.
\end{enumerate}
\end{theorem}

\begin{proof}
\noindent\textbf{Part (i): Proof of Equivalence}

\vspace{1ex}
\noindent($\Rightarrow$) Assume \(M\) is a \(\theta\)-supermartingale. By definition, its dynamics are \(dM_t = \mu_t dt + Z_t dB_t\), where \(\mu_t = \alpha_t - F(t,M_t,Z_t,\Theta_t)\) and \(\alpha_t \le 0\). Let \(a_t \coloneqq -\alpha_t \ge 0\). Then \(\mu_t = -a_t - F(t,M_t,Z_t,\Theta_t)\). By definition, \(\Theta_s\) is a maximizer, so \(F(s,M_s,Z_s,\Theta_s) = V(s,M_s,Z_s)\) almost everywhere. The dynamics are \(\mu_t = -a_t - V(t,M_t,Z_t)\). In integral form from $t$ to $T$:
\begin{equation}\label{eq:proof_decomp_integral}
    M_t = M_T - \int_t^T \mu_s ds + \int_t^T Z_s dB_s = M_T + \int_t^T \big( V(s,M_s,Z_s) + a_s \big) ds - \int_t^T Z_s dB_s.
\end{equation}
Let \(A_t \coloneqq \int_0^t a_s ds\). Since \(a_s \ge 0\), \(A\) is a continuous, non-decreasing process with \(A_0=0\). The term \(\int_t^T a_s ds\) is \(A_T - A_t\). Substituting gives the RBSDE representation \eqref{eq:rbsde_representation}.

\vspace{1ex}
\noindent($\Leftarrow$) Assume there exists a triplet \((M, Z, A)\) satisfying the RBSDE \eqref{eq:rbsde_representation}. Let \(a_t\) be the time-derivative of \(A_t\), so \(a_t \ge 0\) a.e. By the Aumann-Yankov measurable selection theorem, there exists a measurable selector \(\hat{\Theta}(s,\omega,y,z)\) for the argmax correspondence. Define the process \(\Theta_s \coloneqq \hat{\Theta}(s,M_s,Z_s)\). By this construction, \(V(s,M_s,Z_s) = F(s,M_s,Z_s,\Theta_s)\) a.e. Substituting this into the RBSDE yields:
\[
    M_t = M_T + \int_t^T F(s, M_s, Z_s, \Theta_s) \dd s + (A_T - A_t) - \int_t^T Z_s \dd B_s.
\]
This defines a process whose differential is \(dM_t = [-F(t,M_t,Z_t,\Theta_t) - a_t]dt + Z_t dB_t\). This matches the structure of a \(\theta\)-semimartingale with excess drift \(\alpha_t = -a_t \le 0\), proving that \(M\) is a \(\theta\)-supermartingale.

\vspace{1ex}
\noindent\textbf{Part (ii): Proof of Uniqueness}

Assume the stronger regularity conditions of our main theorems hold. Let \(M\) be a \(\theta\)-supermartingale. From part (i), any such \(M\) defines a solution triplet \((M, Z, A)\) to the RBSDE \eqref{eq:rbsde_representation}. The foundational theory of RBSDEs \cite{ElKaroui1997, Kobylanski2000} guarantees that for a driver \(V\) that is (locally) Lipschitz, the solution pair \((Z,A)\) corresponding to a given process \(M\) is unique.

With \(Z\) uniquely determined, we must show that the selection \(\Theta\) is also unique. Any such \(\Theta\) must satisfy \(\Theta_s(\omega) \in \underset{\theta \in \Ucal}{\argmax} \; F(s, \omega, M_s(\omega), Z_s(\omega), \theta)\) a.e. Our main results, Theorem \ref{thm:theta_bsde_wellposedness} and Theorem \ref{thm:existence_geometric}, were formulated precisely to ensure that this argmax set is a singleton for almost every state visited by a solution. The pathwise consistency argument applies directly to the uniquely determined path \((s, \omega, M_s(\omega), Z_s(\omega))\). Therefore, \(\Theta\) is uniquely determined. This implies the uniqueness of the entire decomposition.
\end{proof}

\subsection{The \texorpdfstring{$\theta$}{theta}-Girsanov Theorem}
\label{subsec:theta_girsanov_theorem}

A central pillar of any stochastic calculus is a change of measure theorem, which provides the mechanism for transforming a semimartingale into a martingale. In this section, we establish the analogue for the \(\theta\)-calculus. Unlike classical or G-calculus frameworks where the change of measure is typically static, the \(\theta\)-Girsanov theorem must contend with a dynamically selected model \(\Theta_t\). The key insight is that the change of measure must itself be dynamic, constructed pathwise to absorb the excess drift of a \(\theta\)-semimartingale while leaving the core ambiguity structure invariant.

This requires two structural conditions: first, that the volatility process \(Z_t\) is sufficiently non-degenerate to allow for a change of drift; and second, that the excess drift is admissible, meaning it lies in the space spanned by the available sources of noise.

\begin{theorem}[The \(\theta\)-Girsanov Theorem]
\label{thm:theta_girsanov}
Let the primitives \((F, \Ucal)\) satisfy the conditions for the well-posedness of the \(\theta\)-BSDE system. Let \(M \in \Scalp{2}([0,T]; \R^k)\) be a \(k\)-dimensional \(\theta\)-semimartingale with the unique decomposition
\begin{equation}
    \dd M_t = \left(\alpha_t - F(t, M_t, Z_t, \Theta_t)\right) \dd t + Z_t \dd B_t,
\end{equation}
where \(Z \in \Hcalp{2}([0,T]; \R^{k \times d})\) and \(\alpha \in \Lcalp{2}([0,T]; \R^k)\). Assume the following conditions hold:
\begin{enumerate}[label=(\roman*)]
    \item \textbf{(Non-degenerate Volatility)} The volatility process \(Z_t\) has full column rank \(d\) for \(P_0 \otimes dt\)-almost every \((t,\omega)\). This implies that the \(d \times d\) matrix \(Z_t^T Z_t\) is invertible a.e.
    \item \textbf{(Drift Admissibility)} The excess drift process \(\alpha_t\) lies in the range (column space) of the volatility process \(Z_t\) for a.e. \((t,\omega)\), i.e., \(\alpha_t \in \mathcal{R}(Z_t)\).
\end{enumerate}
Define the \(\R^d\)-valued process \(\eta = (\eta_t)_{t\in[0,T]}\) by
\begin{equation}\label{eq:girsanov_kernel}
    \eta_t \coloneqq -(Z_t^T Z_t)^{-1} Z_t^T \alpha_t.
\end{equation}
If \(\eta\) satisfies Novikov's condition, \(\mathbb{E}_{P_0}\left[\exp\left(\frac{1}{2}\int_0^T \norm{\eta_s}^2 \dd s\right)\right] < \infty\), then under the equivalent probability measure \(P^\eta\) defined by the density \(\frac{\dd P^\eta}{\dd P_0} \coloneqq L_T^\eta = \exp\left( \int_0^T \eta_s^T \dd B_s - \frac{1}{2} \int_0^T \norm{\eta_s}^2 \dd s \right)\), the process \(M\) is a \(\theta\)-martingale. Its dynamics with respect to the \(P^\eta\)-Brownian motion \(B^\eta_t \coloneqq B_t - \int_0^t \eta_s \dd s\) are given by
\begin{equation}
    \dd M_t = -F(t, M_t, Z_t, \Theta_t) \dd t + Z_t \dd B_t^\eta.
\end{equation}
\end{theorem}

\begin{proof}
The proof proceeds by a direct application of the classical Girsanov theorem to the dynamics of the \(\theta\)-semimartingale \(M\).

\noindent\textbf{Step 1: Well-posedness of the transformation.}
Under the Non-degenerate Volatility assumption (i), the matrix \(Z_t^T Z_t\) is invertible a.e., which ensures that the process \(\eta_t\) in \eqref{eq:girsanov_kernel} is well-defined. Since \(\alpha\) and \(Z\) are progressively measurable, \(\eta\) is also progressively measurable. The assumption that \(\eta\) satisfies Novikov's condition ensures that \(L_T^\eta\) is a well-defined Radon-Nikodym density, and the measure \(P^\eta\) is equivalent to \(P_0\).

\noindent\textbf{Step 2: Transformation of the dynamics.}
By the classical Girsanov theorem, the process \(B^\eta\) is a standard \(d\)-dimensional Brownian motion under the measure \(P^\eta\). We can express the original Brownian motion as \(\dd B_t = \dd B_t^\eta + \eta_t \dd t\). We substitute this into the dynamics of \(M\):
\begin{align*}
    \dd M_t &= \left(\alpha_t - F(t, M_t, Z_t, \Theta_t)\right) \dd t + Z_t (\dd B_t^\eta + \eta_t \dd t) \\
    &= \left( \alpha_t + Z_t \eta_t - F(t, M_t, Z_t, \Theta_t) \right) \dd t + Z_t \dd B_t^\eta.
\end{align*}
The new drift of the process \(M\) under the measure \(P^\eta\) is \(\mu_t^\eta \coloneqq \alpha_t + Z_t \eta_t - F(t, M_t, Z_t, \Theta_t)\).

\noindent\textbf{Step 3: Cancellation of the excess drift.}
We now analyze the drift adjustment term \(Z_t \eta_t\). By substituting the definition of \(\eta_t\):
\[
    Z_t \eta_t = Z_t \left( -(Z_t^T Z_t)^{-1} Z_t^T \alpha_t \right) = - \left( Z_t (Z_t^T Z_t)^{-1} Z_t^T \right) \alpha_t.
\]
The matrix \(P_{Z_t} \coloneqq Z_t (Z_t^T Z_t)^{-1} Z_t^T\) is the orthogonal projection matrix onto the column space (range) of \(Z_t\), \(\mathcal{R}(Z_t)\). The Drift Admissibility assumption (ii) states that \(\alpha_t \in \mathcal{R}(Z_t)\). The projection of a vector onto a subspace in which it already lies is the vector itself. Therefore, \(P_{Z_t} \alpha_t = \alpha_t\).
This implies that the drift adjustment term is precisely the negative of the excess drift:
\[
    Z_t \eta_t = -\alpha_t.
\]
Substituting this back into the expression for the new drift \(\mu_t^\eta\), we find:
\[
    \mu_t^\eta = (\alpha_t - \alpha_t) - F(t, M_t, Z_t, \Theta_t) = -F(t, M_t, Z_t, \Theta_t).
\]
\noindent\textbf{Step 4: Conclusion.}
The dynamics of the process \(M\) under the measure \(P^\eta\) are therefore given by
\[
    \dd M_t = -F(t, M_t, Z_t, \Theta_t) \dd t + Z_t \dd B_t^\eta.
\]
This corresponds exactly to the definition of a \(\theta\)-martingale (Definition 6.2). The crucial point is that the processes \(M\), \(Z\), and \(\Theta\) are pathwise identical under the equivalent change of measure; only the interpretation of the dynamics changes. The process \(\Theta_t\), being a function of the path \((M_t, Z_t)\), remains the unique admissible maximizer. This completes the proof.
\end{proof}

\begin{remark}[On the Structural Assumptions and Interpretation]
\label{rem:girsanov_interpretation}
The \(\theta\)-Girsanov theorem provides a powerful tool for analysis, but its assumptions carry significant conceptual weight.
\begin{enumerate}
    \item \textbf{On Non-degeneracy:} The condition that \(Z_t\) has full column rank is the mathematical statement that the sources of noise are sufficiently rich. It means that for any direction of excess drift \(\alpha_t\) that is admissible, there is a unique combination of shocks \(\eta_t\) that can exactly counteract it.
    \item \textbf{On Admissibility:} The Drift Admissibility condition, \(\alpha_t \in \mathcal{R}(Z_t)\), is a deep consistency requirement. It posits that any non-martingale behavior (the excess drift \(\alpha_t\)) must be explainable by the sources of risk that are already present in the model (the columns of \(Z_t\)). Any drift orthogonal to this space cannot be eliminated by a change of measure and represents a fundamental, unhedgeable source of growth or decay.
    \item \textbf{Invariance of Ambiguity Structure:} The most salient feature of this theorem is that the ambiguity operator \(F(\cdot, M_t, Z_t, \Theta_t)\) remains unchanged. The change of measure eliminates the excess drift \(\alpha_t\), but it does not resolve the underlying Knightian uncertainty. It transforms a process that is not a martingale with respect to the \(\theta\)-calculus into one that is, allowing for valuation in a dynamically constructed ambiguity-neutral world.
\end{enumerate}
This theorem opens a path for developing a comprehensive pricing and hedging theory within the \(\theta\)-calculus, where the cost of eliminating model-unrelated drifts can be precisely quantified.
\end{remark}

\section{\texorpdfstring{$\theta$}{theta}-Martingales and Geometric Sensitivity}
\label{sec:geometric_sensitivity_revised}

A central motivation for the \(\theta\)-expectation framework is to overcome the identifiability challenge inherent in classical sublinear theories, which are insensitive to non-convex structures of uncertainty. This section investigates the conditions under which the \(\theta\)-framework succeeds in this goal. We demonstrate a crucial dichotomy: if the primitive driver \(F\) depends linearly on the parameter \(\theta\), the \(\theta\)-expectation collapses to the classical G-expectation, inheriting its limitations. However, when the dependence is non-linear—as in the projection drivers analyzed in \Cref{sec:main_result_geometric}—the framework becomes genuinely sensitive to the geometry of \(\Ucal\), thereby resolving the identifiability challenge.

\subsection{The Convexification Impasse in Linear Models}

We begin by analyzing the case where the driver exhibits linear dependence on the uncertainty parameter. This includes the canonical model of volatility ambiguity, which corresponds to the G-calculus.

\begin{definition}[Linear Driver]
\label{def:linear_driver}
A driver \(F(t,\omega,y,z,\theta)\) is \emph{linear in the parameter} if, for a.e. \((t,\omega)\), the map \(\theta \mapsto F(t,\omega,y,z,\theta)\) is an affine function for all \((y,z)\).
\end{definition}

\begin{example}[Canonical Volatility Ambiguity (G-Calculus)]
Let \(\Ucal \subset \Sdp\) be a compact set of volatility matrices. The canonical driver modeling pure volatility ambiguity is:
\[ F(z, \theta) = \frac{1}{2} \Tr(z \theta z^T) = \frac{1}{2} \scpr{\theta}{z^T z}_F. \]
This is linear in \(\theta\).
\end{example}

When the driver is linear, the \(\theta\)-expectation framework fails to distinguish a non-convex set from its convex hull. This is a consequence of the fundamental properties of linear optimization over compact sets, combined with the uniqueness enforced by our main theorems.

\begin{theorem}[Indistinguishability under Linear Drivers]
\label{thm:indistinguishability_linear}
Let the driver \(F\) be linear in the parameter \(\theta\). Let \(\Ucal_1\) be a compact uncertainty set and \(\Ucal_2 = \overlineconv(\Ucal_1)\). Assume the well-posedness conditions (e.g., Theorem \ref{thm:existence_geometric}) hold for both \(\Ucal_1\) and \(\Ucal_2\), ensuring unique solutions \((Y^{(i)}, Z^{(i)}, \Theta^{(i)})\) for \(i=1,2\). Then the value and volatility processes coincide, and the maximizer processes are identical almost surely:
\[ (Y^{(1)}_t, Z^{(1)}_t) = (Y^{(2)}_t, Z^{(2)}_t) \quad \text{and} \quad \Theta^{(1)}_t = \Theta^{(2)}_t, \quad P_0 \otimes dt\text{-a.e.} \]
\end{theorem}

\begin{proof}
The proof relies on the fact that the supremum of a linear function over a set is identical to the supremum over its convex hull, and that the maximizers must be extreme points.

\vspace{1ex}
\noindent\textbf{Step 1: Coincidence of Effective Drivers.}
The effective driver for \(\Ucal_i\) is given by \(\tilde{F}_i(y,z) = \sup_{\theta \in \Ucal_i} F(y,z,\theta)\). (We suppress the dependence on \((t,\omega)\) for clarity). Since the map \(L_{y,z}(\theta) \coloneqq F(y,z,\theta)\) is linear in \(\theta\), a standard result in convex analysis \cite{Rockafellar1970} states that the supremum of a linear function over a compact set is equal to the supremum over its closed convex hull:
\[ \sup_{\theta \in \Ucal_1} L_{y,z}(\theta) = \sup_{\theta \in \overlineconv(\Ucal_1)} L_{y,z}(\theta). \]
Thus, \(\tilde{F}_1(y,z) = \tilde{F}_2(y,z)\) for all \((y,z)\).

\vspace{1ex}
\noindent\textbf{Step 2: Coincidence of Value and Volatility Processes.}
Since the effective drivers are identical and the terminal condition \(\xi\) is the same, the BSDEs governing the solutions are identical:
\[ Y^{(i)}_t = \xi + \int_t^T \tilde{F}_1(s, Y^{(i)}_s, Z^{(i)}_s) \dd s - \int_t^T Z^{(i)}_s \dd B_s. \]
By the uniqueness of the solution to this effective BSDE (established in the proofs of Theorems \ref{thm:theta_bsde_wellposedness} and \ref{thm:existence_geometric}), we must have \((Y^{(1)}, Z^{(1)}) = (Y^{(2)}, Z^{(2)})\). We denote this common solution by \((Y, Z)\).

\vspace{1ex}
\noindent\textbf{Step 3: Coincidence of Maximizer Processes.}
We now examine the maximizer processes. By the definition of the \(\theta\)-BSDE solution, \(\Theta^{(i)}_t\) must satisfy the pointwise optimality condition:
\[ \Theta^{(i)}_t(\omega) \in \underset{\theta \in \Ucal_i}{\argmax} \; F(t, Y_t(\omega), Z_t(\omega), \theta). \]
The well-posedness assumption (specifically, the results derived from Geometric Regularity in Theorem \ref{thm:existence_geometric}) ensures that for \(P_0 \otimes dt\)-a.e. \((t,\omega)\), the maximizer over the (potentially larger) set \(\Ucal_2\) is unique. Let \(\theta^*_t(\omega)\) be this unique maximizer:
\[ \{\Theta^{(2)}_t(\omega)\} = \underset{\theta \in \Ucal_2}{\argmax} \; F(t, Y_t(\omega), Z_t(\omega), \theta). \]
We now invoke the relationship between extreme points and convex hulls. Since \(F\) is linear in \(\theta\) and \(\Ucal_2\) is compact and convex, Bauer's Maximum Principle implies that the maximum must be attained at an extreme point of \(\Ucal_2\). Since the maximizer is unique a.e., \(\Theta^{(2)}_t(\omega)\) must be an extreme point.

Furthermore, by the converse of the Krein-Milman theorem (Milman's theorem), since \(\Ucal_1\) is compact and \(\Ucal_2 = \overlineconv(\Ucal_1)\), the set of extreme points of \(\Ucal_2\) is a subset of \(\Ucal_1\).

Therefore, the unique maximizer \(\Theta^{(2)}_t(\omega)\) must belong to \(\Ucal_1\). Since \(\Theta^{(2)}_t(\omega) \in \Ucal_1\) and it maximizes \(F\) over the larger set \(\Ucal_2\), it certainly maximizes \(F\) over the smaller set \(\Ucal_1\). Furthermore, since the maximizer over \(\Ucal_2\) is unique a.e., the maximizer over \(\Ucal_1\) must also be unique and coincide with it.
Thus, \(\Theta^{(1)}_t(\omega) = \Theta^{(2)}_t(\omega)\) for \(P_0 \otimes dt\)-a.e. \((t,\omega)\).
\end{proof}

\begin{remark}[The Limitation of Linear Models and Degeneracy]
\label{rem:limitation_linear}
Theorem \ref{thm:indistinguishability_linear} reveals a fundamental limitation. When ambiguity is modeled via a linear dependence on the parameters (as in canonical volatility ambiguity or G-calculus), the \(\theta\)-framework offers no advantage over the classical sublinear theory in terms of identifiability. The entire solution triplet \((Y, Z, \Theta)\) is insensitive to the non-convex geometry of \(\Ucal\).
\end{remark}

\subsection{Non-Linear Drivers and Geometric Sensitivity}

The true power of the \(\theta\)-framework emerges when the driver \(F\) exhibits non-linear dependence on the parameter \(\theta\). In such cases, the convexification of the uncertainty set fundamentally alters the optimization landscape, leading to distinct effective drivers and, consequently, distinct valuation processes.

We illustrate this with the projection-type driver introduced in \Cref{sec:main_result_geometric}.

\begin{definition}[Projection Driver]
The projection driver is defined as:
\[ F(t,x,y,z,\theta) = h(t,x,y,z) - \frac{1}{2}\norm{\theta - G(t,x,y,z)}_F^2. \]
\end{definition}
This driver is strictly concave (quadratic) in \(\theta\).

\begin{proposition}[Sensitivity of the Projection Driver]
\label{prop:sensitivity_projection}
Let \(F\) be the projection driver. Let \(\Ucal_1\) be a compact, non-convex set, and \(\Ucal_2 = \overlineconv(\Ucal_1)\). Then the effective drivers \(\tilde{F}_1\) and \(\tilde{F}_2\) are generally distinct, satisfying \(\tilde{F}_1 \le \tilde{F}_2\).
\end{proposition}

\begin{proof}
The effective driver is given by the value function of the maximization problem:
\begin{align*}
    \tilde{F}_i(t,x,y,z) &= \sup_{\theta \in \Ucal_i} \left( h(t,x,y,z) - \frac{1}{2}\norm{\theta - G(t,x,y,z)}_F^2 \right) \\
    &= h(t,x,y,z) - \frac{1}{2} \inf_{\theta \in \Ucal_i} \norm{\theta - G(t,x,y,z)}_F^2 \\
    &= h(t,x,y,z) - \frac{1}{2} \dist(G(t,x,y,z), \Ucal_i)^2.
\end{align*}
Let \(P = G(t,x,y,z)\). We compare \(\dist(P, \Ucal_1)^2\) and \(\dist(P, \Ucal_2)^2\).
Since \(\Ucal_1 \subset \Ucal_2\), it is immediate that \(\dist(P, \Ucal_2) \le \dist(P, \Ucal_1)\), which implies \(\tilde{F}_2 \ge \tilde{F}_1\).

To show they are distinct, we must find a point \(P\) such that \(\dist(P, \Ucal_2) < \dist(P, \Ucal_1)\). Since \(\Ucal_1\) is non-convex, there exists a point \(\theta^* \in \Ucal_2 \setminus \Ucal_1\). Let \(P = \theta^*\).
Then \(\dist(P, \Ucal_2) = \dist(\theta^*, \Ucal_2) = 0\), since \(\theta^* \in \Ucal_2\).
However, since \(\Ucal_1\) is compact and \(\theta^* \notin \Ucal_1\), the distance must be strictly positive: \(\dist(P, \Ucal_1) = \delta > 0\).

Therefore, at any state \((t,x,y,z)\) such that \(G(t,x,y,z) = \theta^*\), we have:
\begin{align*}
    \tilde{F}_1 &= h - \frac{1}{2}\delta^2 \\
    \tilde{F}_2 &= h - 0 = h.
\end{align*}
Thus, \(\tilde{F}_1 < \tilde{F}_2\). The effective drivers are distinct.
\end{proof}

\begin{corollary}[Distinct Valuations for Non-Convex Sets]
Under the conditions of Proposition \ref{prop:sensitivity_projection}, and assuming the well-posedness conditions of Theorem \ref{thm:existence_geometric} hold, the value processes \(Y^{(1)}_t\) and \(Y^{(2)}_t\) associated with \(\Ucal_1\) and \(\Ucal_2\) satisfy \(Y^{(1)}_t \le Y^{(2)}_t\). The inequality is generally strict, provided the terminal condition $\xi$ and the map $G$ are sufficiently rich such that the solution process explores the region where the drivers differ.
\end{corollary}

\begin{proof}
The processes \(Y^{(i)}\) solve the BSDEs with drivers \(\tilde{F}_i\). Since \(\tilde{F}_1 \le \tilde{F}_2\), the comparison theorem for BSDEs (which applies to the quadratic growth drivers here, see \cite{Kobylanski2000}) implies that \(Y^{(1)}_t \le Y^{(2)}_t\). If the inequality \(\tilde{F}_1 < \tilde{F}_2\) holds on a set visited with positive probability by the solution process (guaranteed by the non-degeneracy assumptions of Theorem \ref{thm:existence_geometric}), the strict comparison principle implies \(Y^{(1)}_t < Y^{(2)}_t\) with positive probability.
\end{proof}

This analysis confirms the central thesis of the paper: the \(\theta\)-expectation framework provides a calculus that is genuinely sensitive to non-convex uncertainty, provided the dependence of the driver on the uncertainty parameter is non-linear. This contrasts sharply with the G-calculus and related frameworks, which are inherently tied to linear dependence and convexification.


\section{The \texorpdfstring{$\theta$}{theta}-Brownian Motion}
\label{sec:theta_system_and_sensitivity}

A foundational element of any stochastic calculus is its canonical martingale process, such as the classical Brownian motion or the G-Brownian motion in the theory of sublinear expectations. A natural objective within our framework is to define and analyze the analogous object: the \emph{$\theta$-Brownian Motion}. However, as we will demonstrate, a direct constructive approach reveals a subtle but critical insight: the fundamental object of our theory is not a single process, but a full dynamic \emph{system}, and its most powerful features are only revealed when the system is non-trivial.

This section first attempts a canonical construction, which leads to a degenerate case, highlighting a crucial limitation. We then pivot to the more profound question of when our framework can distinguish a non-convex uncertainty set $\Ucal$ from its convex hull $\overlineconv(\Ucal)$. We prove that this sensitivity depends critically on the non-linearity of the driver. For linear drivers, including the canonical G-calculus case, the framework is insensitive. For non-linear drivers, such as the projection type central to our main result, we show that the full solution triplet $(Y,Z,\Theta)$ provides a richer, geometrically sensitive description of the system, thereby resolving the identifiability challenge that limits classical theories.

\subsection{The Canonical System and the Degeneracy Problem}
To isolate the effects of volatility ambiguity, we consider a driver that is purely quadratic in the volatility process $Z$, which directly corresponds to the generator of a diffusion process.

\begin{definition}[Canonical Driver and System]
\label{def:canonical_driver_revised}
Let the uncertainty set $\Ucal \subset \Sdp$ be a compact set of strictly positive definite matrices. The \emph{canonical driver} for the $\theta$-system is the function $F: \R^{1\times d} \times \Ucal \to \R$ given by:
\begin{equation}
    F(z, \theta) \coloneqq \frac{1}{2} \Tr(z \theta z^T).
\end{equation}
The \emph{canonical $\theta$-system} is the unique solution triplet $(W, Z, \Theta)$ to the $\theta$-BSDE system (Definition \ref{def:theta_bsde_system}) with this driver and the terminal condition $\xi = 0$. The component process $W$ would be the candidate for the \emph{$\theta$-Brownian motion value process}.
\end{definition}

Our well-posedness theorems apply to this driver, as it exhibits quadratic growth in $z$ and is smooth. However, the solution in this specific case is trivial.

\begin{proposition}[Degeneracy of the Canonical System]
\label{prop:degeneracy_of_canonical_system}
Let the canonical driver and uncertainty set $\Ucal$ satisfy the conditions for well-posedness (either Theorem 3.2 or 4.2). The unique solution to the canonical $\theta$-BSDE with $\xi=0$ is the trivial solution $(W, Z) = (0, 0)$.
\end{proposition}

\begin{proof}
We first verify that $(W_t, Z_t) = (0,0)$ for all $t \in [0,T]$ is a valid solution pair.
The terminal condition is met: $W_T = 0 = \xi$.
The backward evolution equation becomes:
\[ 0 = 0 + \int_t^T F(s, 0, \Theta_s) \dd s - \int_t^T 0 \cdot \dd B_s. \]
The driver is $F(s, Z_s, \Theta_s) = F(s, 0, \Theta_s) = \frac{1}{2}\Tr(0 \cdot \Theta_s \cdot 0^T) = 0$. The equation simplifies to $0=0$, which is satisfied. Thus, $(0,0)$ is a solution.
By the uniqueness results established in Theorem 3.2 and Theorem 4.2, which guarantee a unique solution pair $(Y,Z)$ for the effective BSDE, this must be the only solution.
\end{proof}

This result leads to a fundamental impasse for the maximizer process $\Theta_t$.

\begin{remark}[The Maximizer Uniqueness Impasse]
\label{rem:maximizer_impasse}
Proposition \ref{prop:degeneracy_of_canonical_system} demonstrates that the only possible volatility process is $Z_t \equiv 0$. The pointwise optimality condition for $\Theta_t$ is therefore:
\[ \Theta_t(\omega) \in \underset{\theta \in \Ucal}{\argmax} \; F(Z_t(\omega), \theta) = \underset{\theta \in \Ucal}{\argmax} \; F(0, \theta) = \underset{\theta \in \Ucal}{\argmax} \; 0 = \Ucal. \]
The set of maximizers is the entire uncertainty set $\Ucal$. If $\Ucal$ is not a singleton, the maximizer is not unique. This violates the central requirement of both of our well-posedness theorems (either the existence of a unique maximizer map $\Thetacal^*$ or the pathwise uniqueness of the maximizer).

This reveals that a non-trivial system state ($Z_t \not\equiv 0$) is necessary for the ambiguity-resolving mechanism of our framework to become active. The notion of a "$\theta$-Brownian Motion" as a single, canonical process is therefore ill-posed; the meaningful object is the full $\theta$-system driven by a non-trivial claim $\xi$.
\end{remark}

\subsection{Geometric Sensitivity: When does \texorpdfstring{$\Ucal \neq \overlineconv(\Ucal)$}{Ucal != conv(Ucal)} Matter?}
Having established that a non-trivial state is required, we now investigate the core question of identifiability. We analyze the conditions under which the framework can distinguish between a non-convex set $\Ucal_1$ and its convex hull $\Ucal_2 = \overlineconv(\Ucal_1)$.

\begin{theorem}[Indistinguishability under Linear Drivers]
\label{thm:indistinguishability_linear_revised}
Let the driver $F(t,y,z,\theta)$ be affine in the parameter $\theta$. Let $\Ucal_1$ be a compact uncertainty set and $\Ucal_2 = \overlineconv(\Ucal_1)$. Assume the well-posedness conditions of Theorem \ref{thm:existence_geometric} hold for both $\Ucal_1$ and $\Ucal_2$. Let $(Y^{(i)}, Z^{(i)}, \Theta^{(i)})$ be the unique solution for $\Ucal_i$ ($i=1,2$) with a given terminal condition $\xi$. Then the solutions are identical:
\[ (Y^{(1)}_t, Z^{(1)}_t) = (Y^{(2)}_t, Z^{(2)}_t) \quad \text{and} \quad \Theta^{(1)}_t = \Theta^{(2)}_t, \quad P_0 \otimes dt\text{-a.e.} \]
\end{theorem}
\begin{proof}
The proof proceeds in three steps, leveraging the properties of linear optimization and the uniqueness of the BSDE solution.
\begin{enumerate}
    \item \textbf{Coincidence of Effective Drivers:} The effective driver is $\tilde{F}_i(t,y,z) = \sup_{\theta \in \Ucal_i} F(t,y,z,\theta)$. Since $F$ is affine (and continuous) in $\theta$, a fundamental result of convex analysis \cite{Rockafellar1970} states that the supremum of an affine function over a compact set depends only on the set's closed convex hull. Thus,
    \[ \tilde{F}_1(t,y,z) = \sup_{\theta \in \Ucal_1} F(t,y,z,\theta) = \sup_{\theta \in \Ucal_2} F(t,y,z,\theta) = \tilde{F}_2(t,y,z). \]
    The effective drivers for both problems are identical.

    \item \textbf{Coincidence of $(Y,Z)$:} Since the terminal condition $\xi$ and the effective drivers are identical for both systems, the BSDEs governing the $(Y,Z)$ pairs are identical. By the uniqueness part of our main theorems, the solutions must coincide: $(Y^{(1)}, Z^{(1)}) = (Y^{(2)}, Z^{(2)})$. Let us denote this common solution by $(Y,Z)$.

    \item \textbf{Coincidence of $\Theta$:} By definition, $\Theta^{(i)}_t$ is the almost surely unique maximizer of $F(t,Y_t,Z_t,\theta)$ over $\Ucal_i$. Since the map $\theta \mapsto F(t,Y_t,Z_t,\theta)$ is affine and $\Ucal_2$ is compact and convex, Bauer's Maximum Principle implies that the maximum over $\Ucal_2$ is attained at an extreme point of $\Ucal_2$. By uniqueness, the maximizer $\Theta^{(2)}_t$ must be such an extreme point. Furthermore, by the converse of the Krein-Milman theorem (Milman's theorem), the set of extreme points of $\Ucal_2 = \overlineconv(\Ucal_1)$ is contained in $\Ucal_1$. Therefore, the unique maximizer $\Theta^{(2)}_t$ must belong to $\Ucal_1$.
    Since $\Theta^{(2)}_t \in \Ucal_1$ and it maximizes $F$ over the larger set $\Ucal_2$, it must also maximize $F$ over the smaller set $\Ucal_1$. As the maximizer over $\Ucal_1$ is also assumed to be unique a.e., it must be that $\Theta^{(1)}_t = \Theta^{(2)}_t$.
\end{enumerate}
The entire solution triplet is indistinguishable.
\end{proof}

\begin{remark}[The Limitation of G-Calculus]
The canonical G-calculus driver $F(z,\theta) = \frac{1}{2}\Tr(z\theta z^T)$ is linear in $\theta$. Theorem \ref{thm:indistinguishability_linear_revised} rigorously proves that any framework based on such a driver is structurally insensitive to non-convex geometry. This is the fundamental limitation our theory seeks to overcome.
\end{remark}

\subsection{Distinguishing Models with Non-Linear Drivers}
The true power of the $\theta$-framework emerges when the driver is non-linear in $\theta$. The projection driver from our main theorem is the canonical example.

\begin{proposition}[Sensitivity of the Projection Driver]
\label{prop:sensitivity_projection_revised}
Let the driver be the projection driver $F(x,y,z,\theta) = h(x,y,z) - \frac{1}{2}\norm{\theta - G(x,y,z)}_F^2$. Let $\Ucal_1$ be a compact, non-convex set, and $\Ucal_2 = \overlineconv(\Ucal_1)$. The effective drivers are generally distinct and satisfy $\tilde{F}_1 \le \tilde{F}_2$.
\end{proposition}
\begin{proof}
The effective driver is $\tilde{F}_i(v) = h(v) - \frac{1}{2} \dist(G(v), \Ucal_i)^2$, where $v=(x,y,z)$. Since $\Ucal_1 \subset \Ucal_2$, for any point $P=G(v)$, we have $\dist(P, \Ucal_2) \le \dist(P, \Ucal_1)$, which implies $\tilde{F}_2 \ge \tilde{F}_1$.
To show they are distinct, since $\Ucal_1$ is non-convex, there exists a point $P_0 \in \Ucal_2 \setminus \Ucal_1$. If the state of the system is such that $G(v) = P_0$, then $\dist(G(v), \Ucal_2)=0$ while $\dist(G(v), \Ucal_1) > 0$. At this state, $\tilde{F}_1(v) < \tilde{F}_2(v)$.
\end{proof}

\begin{corollary}[Distinct Valuations]
Under the conditions of Proposition \ref{prop:sensitivity_projection_revised}, assuming the non-degeneracy conditions of Theorem \ref{thm:existence_geometric} hold, the value processes satisfy $Y^{(1)}_t \le Y^{(2)}_t$. The inequality is strict with positive probability, provided the map $G$ is sufficiently rich for the process to explore the regions where the drivers differ.
\end{corollary}
\begin{proof}
This follows directly from the comparison principle for BSDEs with (locally) Lipschitz drivers \cite{Kobylanski2000}. Since $\tilde{F}_1 \le \tilde{F}_2$, we have $Y^{(1)} \le Y^{(2)}$. If the strict inequality $\tilde{F}_1 < \tilde{F}_2$ holds on a set of states visited with positive probability (guaranteed by non-degeneracy), the strict comparison principle ensures $Y^{(1)}_t < Y^{(2)}_t$ with positive probability.
\end{proof}

\subsection{The Endogenous Covariance Process as a Distinguishing Statistic}
Even in cases where the value processes might be close, the full $\theta$-system provides a richer, path-dependent statistic that is highly sensitive to the geometry of $\Ucal$.

\begin{definition}[Endogenous Covariance Process]
Let $(Y,Z,\Theta)$ be the unique solution to a $\theta$-BSDE. The \emph{endogenous covariance process} $\mathcal{C}$ is the process defined by
\begin{equation}
    \mathcal{C}_t \coloneqq \int_0^t \Theta_s ds.
\end{equation}
It represents the cumulative, path-dependent ambiguity realized by the system.
\end{definition}

This statistic allows us to rigorously distinguish between models even if the driver were chosen such that the value processes were identical (e.g., if $h$ were adjusted to compensate for the difference in the distance term).

\begin{proposition}[Distinguishability via Endogenous Covariance]
Let $\Ucal_1$ be a non-convex compact set and $\Ucal_2 = \overlineconv(\Ucal_1)$, and use a non-linear driver like the projection driver. Assume the conditions of Theorem \ref{thm:existence_geometric} hold. Let $\mathcal{C}^{(1)}$ and $\mathcal{C}^{(2)}$ be the endogenous covariance processes associated with the solutions. The laws of the processes $\mathcal{C}^{(1)}$ and $\mathcal{C}^{(2)}$ are distinct.
\end{proposition}
\begin{proof}
Let $(Y^{(i)}, Z^{(i)}, \Theta^{(i)})$ be the respective unique solutions. By definition, for $P_0 \otimes dt$-a.e. $(s,\omega)$, we have $\Theta^{(1)}_s(\omega) \in \Ucal_1$ and $\Theta^{(2)}_s(\omega) \in \Ucal_2$.
Since $\Ucal_1$ is a proper non-convex subset of $\Ucal_2$, there exists a non-empty set of points $O \subset \Ucal_2 \setminus \Ucal_1$. Under the non-degeneracy conditions of Theorem \ref{thm:existence_geometric}, the law of the solution process $(Y_s^{(2)}, Z_s^{(2)})$ ensures that the target process $W_s = G(Y_s^{(2)}, Z_s^{(2)})$ explores its state space. With positive probability, it will visit a region where the unique projection onto $\Ucal_2$ is a point $\theta^* \in O$.
For such a path $\omega$ at that time $s$, we have $\Theta^{(2)}_s(\omega) = \theta^* \in \Ucal_2 \setminus \Ucal_1$. However, $\Theta^{(1)}_s(\omega)$ must, by definition, be some other point in $\Ucal_1$. Thus, $\Theta^{(1)}_s(\omega) \neq \Theta^{(2)}_s(\omega)$.
Since $\frac{d}{dt}\mathcal{C}^{(i)}_t = \Theta^{(i)}_t$, the paths of the covariance processes will differ. The set of attainable paths for $\mathcal{C}^{(1)}$ is determined by integrating functions with values in $\Ucal_1$, while for $\mathcal{C}^{(2)}$ it is determined by integrating functions with values in $\Ucal_2$. These sets of paths are different, and thus their laws must be distinct.
\end{proof}

\begin{remark}[The True Innovation]
This section demonstrates that the core innovation of the $\theta$-expectation framework is not the creation of a single "$\theta$-Brownian Motion" process. Instead, it lies in the well-posedness of the full \emph{system} $(Y,Z,\Theta)$. While the value component $Y$ may sometimes coincide with classical valuations, the unique maximizer process $\Theta_t$ provides a path-dependent indicator of the most adversarial model from the primitive set $\Ucal$. This process, and statistics derived from it like the endogenous covariance $\mathcal{C}_t$, are demonstrably sensitive to the specific non-convex geometry of $\Ucal$, thereby founding a richer, more descriptive, and tractable stochastic calculus.
\end{remark}


\section{The Nonlinear Feynman-Kac Formula}
\label{sec:feynman_kac}

This section establishes a unified nonlinear Feynman-Kac formula, rigorously connecting the conditional $\theta$-expectation defined via the solution of the $\theta$-BSDE to a fully nonlinear parabolic partial differential equation (PDE) in the framework of viscosity solutions. This result solidifies the link between the probabilistic representation of our non-subadditive expectation and a corresponding analytic, PDE-based valuation. It is applicable to both classes of well-posed models identified in this paper (Theorem 3.2 and Theorem 4.2), relying only on the properties of the resulting single-valued effective driver.

\subsection{Markovian Framework and the Associated PDE}
We specialize to a Markovian framework. Let the state process $X_t \in \R^k$ evolve according to the stochastic differential equation:
\begin{equation} \label{eq:sde_forward}
    dX_s = b(s, X_s) ds + \sigma(s, X_s) dB_s, \quad X_t = x, \text{ for } s \in [t,T],
\end{equation}
where $B$ is a $d$-dimensional Brownian motion. We denote the solution started at time $t$ from state $x$ as $X_s^{t,x}$. The terminal condition is a function of the final state, $\xi = \phi(X_T)$.

\begin{assumption}[Markovian Primitives] \label{ass:fk_primitives}
\begin{enumerate}[label=(\roman*)]
    \item The coefficients $b: [0,T]\times\R^k \to \R^k$ and $\sigma: [0,T]\times\R^k \to \R^{k\times d}$ are continuous and globally Lipschitz in $x$, uniformly in $t$.
    \item The terminal condition function $\phi: \R^k \to \R$ is continuous and has at most polynomial growth, i.e., $|\phi(x)| \le C(1+|x|^p)$ for some $C,p > 0$.
    \item The primitive driver $F(t,x,y,z,\theta)$ is continuous in all its arguments.
    \item The conditions for well-posedness of the $\theta$-BSDE (either Hypothesis 3.1 or Assumption 4.1) are satisfied. The resulting single-valued effective driver $\tilde{F}(t,x,y,z) \coloneqq \sup_{\theta \in \Ucal} F(t,x,y,z,\theta)$ is continuous and locally Lipschitz in $(y,z)$, uniformly in $(t,x)$. Moreover, $\tilde{F}$ has at most polynomial growth in $(x,y,z)$.
\end{enumerate}
\end{assumption}

In this setting, the solution to the $\theta$-BSDE is a deterministic function of time and state, $Y_t = u(t, X_t)$. The associated fully nonlinear PDE is:
\begin{equation}
\label{eq:nonlinear_pde_revised}
\begin{cases}
    -\partial_t u - \mathcal{L}u - \tilde{F}\big(t, x, u, (\nabla_x u)\sigma(t,x)\big) = 0, & (t,x) \in [0,T) \times \R^k, \\
    u(T,x) = \phi(x), & x \in \R^k.
\end{cases}
\end{equation}
where $\mathcal{L}$ is the infinitesimal generator of $X$:
\[ \mathcal{L}u(t,x) \coloneqq \scpr{b(t,x)}{\nabla_x u(t,x)} + \frac{1}{2} \Tr\left(\sigma(t,x)\sigma(t,x)^T D^2_x u(t,x)\right). \]

\subsection{The Unified Feynman-Kac Theorem}

\begin{theorem}[Unified Nonlinear Feynman-Kac Formula]
\label{thm:feynman_kac_unified}
Let Assumption \ref{ass:fk_primitives} hold. Define the value function $u: [0,T]\times\R^k \to \R$ by the conditional $\theta$-expectation:
\[ u(t,x) \coloneqq \Efrakgen_t[\phi(X_T^{t,x})] = Y_t^{t,x}, \]
where $(Y^{t,x}, Z^{t,x})$ is the unique solution to the effective BSDE
\begin{equation} \label{eq:effective_bsde_fk}
 Y_s = \phi(X_T^{t,x}) + \int_s^T \tilde{F}(r, X_r^{t,x}, Y_r, Z_r) dr - \int_s^T Z_r dB_r, \quad s \in [t,T].
\end{equation}
Then $u(t,x)$ is a continuous function and is the unique viscosity solution to the PDE \eqref{eq:nonlinear_pde_revised} in the class of functions with polynomial growth.
\end{theorem}

\begin{proof}
The proof is executed in four rigorous steps. First, we establish the continuity of the candidate solution $u(t,x)$. Second and third, we prove that $u$ is a viscosity subsolution and supersolution, respectively. Finally, we establish the uniqueness of this solution.

\vspace{1ex}
\noindent\textbf{Step 1: Continuity of the Value Function $u(t,x)$.}
The continuity of $u(t,x)$ follows from the stability properties of BSDEs with respect to their parameters (initial time and state). Let $(t,x)$ and $(t',x')$ be two points in $[0,T]\times\R^k$. Let $(Y,Z)$ and $(Y',Z')$ be the solutions to the effective BSDE \eqref{eq:effective_bsde_fk} corresponding to initial data $(t,x)$ and $(t',x')$, respectively. Then $u(t,x) = Y_t$ and $u(t',x')=Y'_{t'}$.
Standard a priori estimates for BSDEs (see, e.g., \cite[Theorem 2.6]{ElKaroui1997}) yield:
\begin{align*}
|Y_t - Y'_t|^2 &\le C \mathbb{E}\left[ |\phi(X_T^{t,x}) - \phi(X_T^{t',x'})|^2 \mid \mathcal{F}_t \right].
\end{align*}
(For simplicity, we assume $t=t'$ first and consider variation in $x$ only). The continuity of $\phi$ and the continuous dependence of the SDE solution $X_s^{t,x}$ on its initial state $x$ imply that as $x' \to x$, the right-hand side converges to zero by dominated convergence. A similar, though more technical, argument handles variation in $t$. The joint continuity of $u(t,x)$ is a standard result in the theory of BSDEs under our stated assumptions on the coefficients.

\vspace{1ex}
\noindent\textbf{Step 2: $u$ is a viscosity subsolution.}
Let $(t_0, x_0) \in [0,T) \times \R^k$ and let $\varphi \in C^{1,2}([0,T]\times\R^k)$ be a smooth test function such that $u - \varphi$ has a local maximum at $(t_0, x_0)$, and $u(t_0,x_0) = \varphi(t_0,x_0)$. We must prove that
\begin{equation} \label{eq:subsolution_ineq}
    -\partial_t \varphi(t_0,x_0) - \mathcal{L}\varphi(t_0,x_0) - \tilde{F}\big(t_0, x_0, \varphi(t_0,x_0), (\nabla_x \varphi(t_0,x_0))\sigma(t_0,x_0)\big) \le 0.
\end{equation}
Assume for contradiction that the left-hand side is strictly positive. Then for some $\eta > 0$:
\begin{equation} \label{eq:contra_sub}
    -\partial_t \varphi(t_0,x_0) - \mathcal{L}\varphi(t_0,x_0) - \tilde{F}\big(t_0, x_0, \varphi(t_0,x_0), \nabla_x\varphi\sigma(t_0,x_0)\big) = 2\eta > 0.
\end{equation}
By continuity, there exists $\delta > 0$ such that this inequality (with $2\eta$ replaced by $\eta$) holds for all $(t,x)$ in a neighborhood $B_\delta(t_0,x_0) \coloneqq \{(t,x) : |t-t_0| + |x-x_0| < \delta\}$. Also, by the local maximum property, we can choose $\delta$ small enough such that $u(t,x) \le \varphi(t,x)$ for all $(t,x) \in B_\delta(t_0,x_0)$.

Let $X_s \coloneqq X_s^{t_0,x_0}$. Define a stopping time $\tau \coloneqq \inf\{s \ge t_0 : (s,X_s) \notin B_\delta(t_0,x_0)\}$. Let $h>0$ be small enough such that $t_0+h < T$ and the event $\{\tau > t_0+h\}$ has positive probability.
Let $(Y_s, Z_s)$ be the solution to the effective BSDE starting from $(t_0, x_0)$, so $Y_{t_0} = u(t_0,x_0)$. On the interval $[t_0, \tau]$, we have $Y_s = u(s, X_s) \le \varphi(s, X_s)$.
Now, consider the process $Y'_s \coloneqq \varphi(s,X_s)$. By It\^o's formula for $s \in [t_0, \tau]$:
\[
    Y'_s = \varphi(\tau, X_\tau) - \int_s^\tau \left( \partial_r \varphi + \mathcal{L}\varphi \right)(r,X_r) dr - \int_s^\tau (\nabla_x \varphi)\sigma(r,X_r) dB_r.
\]
Let $\hat{Z}_s \coloneqq (\nabla_x \varphi)\sigma(s,X_s)$. Define the process $\hat{F}_s \coloneqq \tilde{F}(s,X_s,Y'_s,\hat{Z}_s)$.
By our contradiction assumption \eqref{eq:contra_sub}, for $s \in [t_0, \tau]$, we have
\[ -(\partial_s \varphi + \mathcal{L}\varphi)(s,X_s) > \eta + \tilde{F}(s,X_s,\varphi(s,X_s), \hat{Z}_s) = \eta + \hat{F}_s. \]
So, we can write the dynamics of $Y'$ as a BSDE with a perturbed driver:
\[ Y'_s = Y'_\tau + \int_s^\tau \left( \hat{F}_r + \Delta_r \right) dr - \int_s^\tau \hat{Z}_r dB_r, \]
where $\Delta_r = -(\partial_r \varphi + \mathcal{L}\varphi)(r,X_r) - \hat{F}_r > \eta > 0$.

Now we compare the two BSDEs on $[t_0, \tau]$:
\begin{align*}
    Y_s &= Y_\tau + \int_s^\tau \tilde{F}(r,X_r, Y_r, Z_r) dr - \int_s^\tau Z_r dB_r \\
    Y'_s &= Y'_\tau + \int_s^\tau \tilde{F}(r,X_r, Y'_r, \hat{Z}_r) dr + \int_s^\tau \Delta_r dr - \int_s^\tau \hat{Z}_r dB_r.
\end{align*}
Let $\delta Y = Y - Y'$, $\delta Z = Z - \hat{Z}$. Subtracting the equations gives:
\[
    \delta Y_s = \delta Y_\tau + \int_s^\tau \left( \tilde{F}(r,X_r, Y_r, Z_r) - \tilde{F}(r,X_r, Y'_r, \hat{Z}_r) \right) dr - \int_s^\tau \Delta_r dr - \int_s^\tau \delta Z_r dB_r.
\]
The term $\delta Y_\tau = Y_\tau - Y'_\tau = u(\tau, X_\tau) - \varphi(\tau, X_\tau) \le 0$. The term $-\int_s^\tau \Delta_r dr \le -\eta(\tau-s)$.
Using standard comparison arguments for BSDEs (which rely on the local Lipschitz nature of $\tilde{F}$), the positive drift term $-\Delta_r$ forces $\delta Y_s < 0$. Specifically, applying a comparison theorem on $[t_0, \tau]$:
\[ Y_{t_0} - Y'_{t_0} \le \mathbb{E}\left[ (Y_\tau - Y'_\tau) - \int_{t_0}^\tau \Delta_r dr \right] \le \mathbb{E}\left[ -\eta(\tau - t_0) \right]. \]
Since $\mathbb{P}(\tau > t_0+h) > 0$ for small $h$, we have $\mathbb{E}[\tau-t_0] > 0$. Thus, the right-hand side is strictly negative.
This implies $Y_{t_0} < Y'_{t_0}$. However, we started with $Y_{t_0} = u(t_0,x_0) = \varphi(t_0,x_0) = Y'_{t_0}$. This is a contradiction. Thus, the assumption \eqref{eq:contra_sub} must be false, and \eqref{eq:subsolution_ineq} must hold.

\vspace{1ex}
\noindent\textbf{Step 3: $u$ is a viscosity supersolution.}
The argument is symmetric. Let $(t_0, x_0) \in [0,T) \times \R^k$ and $\varphi \in C^{1,2}$ be a test function such that $u - \varphi$ has a local minimum at $(t_0, x_0)$ and $u(t_0,x_0) = \varphi(t_0,x_0)$. We must prove:
\begin{equation} \label{eq:supersolution_ineq}
    -\partial_t \varphi(t_0,x_0) - \mathcal{L}\varphi(t_0,x_0) - \tilde{F}\big(t_0, x_0, \varphi(t_0,x_0), (\nabla_x \varphi(t_0,x_0))\sigma(t_0,x_0)\big) \ge 0.
\end{equation}
Assume for contradiction that the left-hand side is strictly negative ($= -2\eta < 0$). We choose $\delta > 0$ such that for $(t,x) \in B_\delta(t_0,x_0)$, we have $u(t,x) \ge \varphi(t,x)$ and $-(\partial_t\varphi+\mathcal{L}\varphi) - \tilde{F}(\cdot, \varphi, (\nabla_x\varphi)\sigma) \le -\eta$.
Following the same procedure as in Step 2, we compare $Y_s = u(s,X_s)$ with $Y'_s = \varphi(s,X_s)$. The difference in their drivers is now $\Delta_r < -\eta < 0$. The terminal value difference is $\delta Y_\tau = u(\tau, X_\tau) - \varphi(\tau, X_\tau) \ge 0$.
The comparison principle now yields:
\[ Y_{t_0} - Y'_{t_0} \ge \mathbb{E}\left[ (Y_\tau - Y'_\tau) - \int_{t_0}^\tau \Delta_r dr \right] \ge \mathbb{E}\left[ \int_{t_0}^\tau (-\Delta_r) dr \right] \ge \mathbb{E}[\eta(\tau-t_0)] > 0. \]
This implies $Y_{t_0} > Y'_{t_0}$, which contradicts the fact that $Y_{t_0} = u(t_0,x_0) = \varphi(t_0,x_0) = Y'_{t_0}$. The contradiction proves that $u$ is a viscosity supersolution.

From Steps 2 and 3, $u(t,x)$ is a viscosity solution to the PDE \eqref{eq:nonlinear_pde_revised}.

\vspace{1ex}
\noindent\textbf{Step 4: Uniqueness.}
Uniqueness is a consequence of the comparison principle for viscosity solutions of second-order parabolic PDEs. The canonical reference is the user's guide by Crandall, Ishii, and Lions \cite{CrandallIshiiLions1992}.
Let $u_1$ and $u_2$ be two viscosity solutions in the class of functions with polynomial growth. The comparison principle states that if $u_1(T,x) \le u_2(T,x)$ for all $x$, then $u_1(t,x) \le u_2(t,x)$ for all $(t,x) \in [0,T]\times\R^k$.
For this principle to apply to our PDE \eqref{eq:nonlinear_pde_revised}, the Hamiltonian
\[ H(t,x,r,p,M) \coloneqq \scpr{b(t,x)}{p} + \frac{1}{2}\Tr(\sigma\sigma^T(t,x) M) + \tilde{F}(t,x,r,p^T\sigma(t,x)) \]
must satisfy certain structural conditions. Under Assumption \ref{ass:fk_primitives}:
\begin{enumerate}
    \item $H$ is continuous in all its arguments. This follows from the continuity of $b, \sigma, \tilde{F}$.
    \item $H$ is degenerate elliptic, i.e., it is non-decreasing in $M$. This holds because $\sigma\sigma^T$ is positive semi-definite.
    \item The dependence of $H$ on $r$ (which corresponds to $u$) must be non-increasing for the comparison to hold. The term $\tilde{F}(t,x,r,z)$ must be non-decreasing in $r$. This is a standard assumption in the BSDE/viscosity solution literature, often ensured by adding a term like $e^{\lambda t}$ to the solution. Assuming this standard structural condition holds, the comparison principle applies. The local Lipschitz property of $\tilde{F}$ is sufficient.
\end{enumerate}
Since our candidate solution $u(t,x)$ and any other potential solution must match the same terminal data $\phi(x)$, the comparison principle implies that they must coincide everywhere. Thus, $u(t,x)$ is the unique viscosity solution in its class.
\end{proof}

\section{Application: General Relativity}

\label{sec:application_gr}

A profound open problem in modern physics lies at the intersection of general relativity and nuclear physics: the determination of the Equation of State (EoS) for matter at supra-nuclear densities, as found inside neutron stars. This physical situation presents a canonical example of structured, non-convex model uncertainty, for which our framework is uniquely suited. The classical sub-additive framework is fundamentally ill-equipped to handle this problem, as its reliance on convexification discards essential information about the distinct nature of competing physical theories.

\subsection{Conceptual Framework: Non-Convex Scientific Uncertainty}
The set of plausible EoS models is not a single convex set. Rather, it is broadly understood to be a disjoint union of distinct theoretical classes. For instance, a simplified but conceptually accurate picture of the total uncertainty set \(\Ucal_{\text{EoS}}\) is:
\[ \Ucal_{\text{EoS}} = \Ucal_{\text{hadronic}} \cup \Ucal_{\text{quark}} \cup \Ucal_{\text{hybrid}}. \]
Here, \(\Ucal_{\text{hadronic}}\) represents models where matter consists solely of hadrons, while \(\Ucal_{\text{quark}}\) represents models allowing for a deconfined quark-matter core. These categories occupy different regions of the parameter space and are separated by physically meaningful gaps.

A classical approach based on sublinear expectations, being dual to a supremum over a convex set of priors, would depend only on the closed convex hull, \(\overlineconv(\Ucal_{\text{EoS}})\). This mathematical procedure would introduce a continuum of "averaged" EoS models that may have no physical basis, representing a blending of distinct physical laws. This fundamentally distorts the nature of the scientific uncertainty.

Our framework, by eschewing sub-additivity, resolves this impasse. It allows for a valuation and inference calculus that respects the non-convex geometry of the primitive uncertainty set. We formulate the problem of constraining the EoS using observational data, such as gravitational waves (GW) from a binary neutron star merger, within the rigorous structure of our main result, \Cref{thm:existence_geometric}.

\subsection{Mathematical Formulation}
We now construct a Markovian system that satisfies the conditions of the Geometric Regularity framework.

\paragraph{\textbf{Primitives of the Model:}}
\begin{itemize}
    \item \textbf{State Process $X_t \in \R^k$:} Let the incoming GW data stream be processed in real time to extract a vector of salient features or sufficient statistics, represented by the state process $X_t$. These statistics could include measures of tidal deformability, post-merger oscillation frequencies, or other quantities sensitive to the EoS. We model its evolution as an Itô process:
    \[ dX_t = b(t, X_t) dt + \sigma(t, X_t) dB_t. \]
    We assume the drift $b$ and diffusion $\sigma$ are smooth with bounded partial derivatives, reflecting a stable data analysis pipeline.

    \item \textbf{Uncertainty Set $\Ucal$:} An EoS model can be parameterized by a vector $p \in \R^m$. The set of physically plausible EoS models, \(\Ucal_{\text{EoS}}\), is a compact subset of $\R^m$. As argued above, \(\Ucal_{\text{EoS}}\) is a finite union of disjoint, compact, and convex sets, \(\Ucal_{\text{EoS}} = \cup_{i=1}^N \mathcal{P}_i\). Our theory is formulated in the space of symmetric matrices $\Sd$, so we introduce a smooth embedding \(\mathcal{E}: \R^m \to \Sd\) to define the primitive uncertainty set for our BSDE:
    \[ \Ucal \coloneqq \mathcal{E}(\Ucal_{\text{EoS}}) \subset \Sd. \]
    This set \(\Ucal\) is compact and non-convex.

    \item \textbf{Driver $F$:} We employ the projection-type driver from \Cref{sec:main_result_geometric}, setting the running utility term $h$ to zero to focus purely on the model selection ambiguity:
    \[ F(t, x, y, z, \theta) = -\frac{1}{2}\norm{\theta - G(t,x)}_F^2. \]

    \item \textbf{Target Map $G$:} The map $G: [0,T] \times \R^k \to \Sd$ represents the instantaneous, ambiguity-neutral best-fit EoS parameter given the data statistics $x$. A transparent and tractable choice is a linear map:
    \[ G(t,x) = G(x) = Ax, \]
    where $A: \R^k \to \Sd$ is a fixed linear transformation representing a baseline inference model.

    \item \textbf{Terminal Payoff $\xi$:} The terminal condition \(\xi = \phi(X_T)\) represents the final scientific utility or value of the discovery, as a function of the total information gathered. We assume $\phi$ is continuous with polynomial growth.
\end{itemize}

\subsection{Verification of the Geometric Regularity Conditions}
We now prove that this physically-motivated model can be constructed to satisfy the hypotheses of our main existence and uniqueness theorem.

\begin{proof}[Verification of Assumption \ref{ass:pathwise_consistency}]
Let $N = \dim(\Sd) = d(d+1)/2$. We verify the conditions in order.
\begin{enumerate}[label=(\roman*)]
    \item \textbf{Projection Driver Structure:} Satisfied by construction, with $h \equiv 0$ and $G$ being a smooth (linear) function of $x$.

    \item \textbf{Geometric Regularity of $\Ucal$:} The primitive sets $\mathcal{P}_i$ are compact and convex, and thus have positive reach. A finite union of disjoint compact sets, each with positive reach, also has positive reach. Furthermore, since $\mathcal{E}$ is a smooth embedding, the image set \(\Ucal = \mathcal{E}(\Ucal_{\text{EoS}})\) is a finite union of disjoint, smoothly embedded compact sets. Such a set is known to have a cut locus that is a semi-algebraic set of Lebesgue measure zero, admitting a finite stratification into smooth submanifolds \cite{Federer1959}. The condition is satisfied.

    \item \textbf{Uniform Ellipticity:} A natural model for the noise in the data statistics $X_t$ is non-degenerate instrumental noise. We can model this by setting $\sigma(t,x) = I_k$, the $k \times k$ identity matrix (assuming $d=k$ for the noise dimension). This yields $\sigma\sigma^T = I_k$, which is uniformly elliptic with $\lambda=1$. The condition is satisfied.

    \item \textbf{Transversality Condition:} This is the crucial hypothesis. Let $g(v) = G(x) = Ax$. The condition requires that for each stratum $\Sigma_j$ of the cut locus $\mathrm{Cut}(\Ucal)$, the map $g$ is transverse to $\Sigma_j$. That is, for every $x \in \R^k$ such that $g(x) \in \Sigma_j$, we must have:
    \[ \Image(D_x g(x)) + T_{g(x)}\Sigma_j = T_{g(x)}\Sd. \]
    Since $g$ is linear, its derivative is the map $A$ itself, so we need $\Image(A) + T_{Ax}\Sigma_j = \Sd$. This is a statement about the linear map $A$ that defines our inference model.

    This condition is not a severe restriction but rather a generic property. The set of linear maps $A: \R^k \to \Sd$ that fail to be transverse to a given submanifold $\Sigma_j$ forms a "thin" set (a lower-dimensional algebraic variety) in the space of all such linear maps. Thus, for almost any choice of the $k$ data statistics (which implicitly defines the map $A$), the transversality condition will hold, provided the dimension $k$ of the data space is sufficiently large. Specifically, for transversality to be possible, the sum of the dimensions of the image of $A$ and the tangent space of the stratum must be at least the dimension of the ambient space. Let $\dim(\Sigma_j) = d_j < N$. The condition becomes $\mathrm{rank}(A) + d_j \ge N$. Since we must satisfy this for all strata, we require $\mathrm{rank}(A) + \max_j(d_j) \ge N$. This can be satisfied even if $k < N$, as long as the data probes directions not contained in the tangent spaces of the cut locus. Thus, for any generic choice of a sufficiently rich statistical model, this condition is met.
\end{enumerate}
All conditions of Assumption \ref{ass:pathwise_consistency} are satisfied.
\end{proof}

\subsection{Interpretation and Scientific Value}
By \Cref{thm:existence_geometric}, our EoS inference problem admits a unique solution triplet $(Y,Z,\Theta)$. The interpretation of this solution provides powerful new tools for scientific inference under ambiguity:
\begin{itemize}
    \item $Y_t = \Efrakgen_t[\phi(X_T)]$ is the robust, dynamically updated value of the final scientific discovery (e.g., the strength of the EoS constraint), evaluated under the most adversarial possible "true" EoS.

    \item $Z_t$ quantifies the sensitivity of this valuation to the instrumental and statistical noise in the data stream.

    \item $\Theta_t$ is the central object of physical interest. It represents the specific EoS model from the primitive set \(\Ucal\), evolving in time, that presents the greatest challenge to the valuation given the data up to time $t$. The path-realization of $\Theta_t(\omega)$ provides a novel, concrete inferential tool. If, for a given observation \(\omega\), the process $\Theta_t(\omega)$ spends the majority of the observation time in one region of $\Ucal$ (e.g., hadronic models), it constitutes strong evidence against the other disjoint regions (e.g., quark models) within this robust, adversarial framework.
\end{itemize}
Our framework thus provides a mathematically rigorous and physically meaningful calculus to navigate non-convex scientific uncertainty. It delivers not only a robust valuation but also a path-dependent "most-likely-adversary" process, \(\Theta_t\), which offers a far richer and more interpretable output than the simple convexification mandated by classical approaches.

\section{Conclusion}
\label{sec:conclusion}

This paper has developed a constructive theory for stochastic valuation under ambiguity that remains sensitive to the specific geometry of non-convex uncertainty sets. We framed the problem as a Backward Stochastic Differential Inclusion, for which existence is general but uniqueness fails. Our main contribution is the identification of two distinct, non-trivial classes of models for which uniqueness can be restored.

First, as a motivating case, we proved well-posedness under a novel analytic condition, the Global Regularity Hypothesis. Second, and more significantly, we have resolved the central challenge of existence and uniqueness for a class of geometrically regular models where this hypothesis fails. This is the paper's principal result (\Cref{thm:existence_geometric}). The proof relies on a conjunction of strong but explicit non-degeneracy conditions, where uniqueness follows from a detailed Malliavin calculus argument that shows that solution paths almost surely avoid the singular sets where the ambiguity is irresolvable.

Crucially, we have shown that both model classes lead to a unified theoretical structure, where the \(\theta\)-BSDE is equivalent to a standard BSDE with a single-valued, (locally) Lipschitz effective driver. This unified framework provides a rigorous and identifiable theory, which we used to define the associated $\theta$-martingales and establish a nonlinear Feynman-Kac formula. By demonstrating that a consistent and tractable theory can be built without the axiom of sub-additivity, this work opens new avenues for the mathematical modeling of structured uncertainty.

\begin{appendices}


\section{A Genuinely Non-Convex Example Satisfying Global Regularity}
\label{app:nonconvex_example_detailed}

The Global Regularity Hypothesis (\Cref{hyp:global_regularity}) imposes exceptionally strong conditions on the problem's primitives, demanding that the argmax correspondence is a single-valued, globally Lipschitz function. While we have argued that this hypothesis is non-generic, it is crucial to establish that it is not vacuous. A non-trivial model satisfying these conditions demonstrates the existence of a benchmark class of tractable non-convex problems and motivates the search for the more general geometric theory developed in \Cref{sec:main_result_geometric}.

This section provides a complete and self-contained proof that such a model exists. We construct a problem where the uncertainty set \(\Ucal\) is manifestly non-convex, yet the interaction between the driver \(F\) and the geometry of \(\Ucal\) forces the optimizing parameter to be unique and to depend on the state variables in a globally regular manner.

\begin{proposition}[A Globally Regular Non-Convex Model]
\label{prop:non_convex_global_regularity}
There exist primitives \((F, \Ucal)\) that satisfy the Global Regularity Hypothesis (\Cref{hyp:global_regularity}) where the uncertainty set \(\Ucal\) is not convex.
\end{proposition}

\begin{proof}
We construct the primitives and verify the hypothesis in three steps.

\vspace{1ex}
\noindent\textbf{Step 1: Construction of the Primitives.}
Let the dimension of the underlying Brownian motion be \(d=2\). The space of parameters is the space of \(2 \times 2\) symmetric matrices, \(\mathbb{S}_2\), equipped with the Frobenius inner product \(\scpr{A}{B}_F \coloneqq \Tr(A^T B) = \Tr(AB)\) and its induced norm \(\norm{A}_F = \sqrt{\Tr(A^2)}\).

\begin{enumerate}[label=(\alph*)]
    \item \textbf{Uncertainty Set \(\Ucal\):} Let the uncertainty set be the disjoint union of two closed balls in \(\mathbb{S}_2\):
    \[
        \Ucal \coloneqq \Bcal(-2I, 1) \cup \Bcal(2I, 1),
    \]
    where \(I\) is the \(2 \times 2\) identity matrix, and \(\Bcal(C,r) \coloneqq \{A \in \mathbb{S}_2 \mid \norm{A-C}_F \le r\}\). The set \(\Ucal\) is compact, as it is a finite union of compact sets. It is manifestly non-convex, as the midpoint between the centers, \(0 \in \mathbb{S}_2\), is not in \(\Ucal\).

    \item \textbf{Driver Function \(F\):} Let the driver \(F\) be independent of \(y\) and have the form:
    \[
        F(t, \omega, z, \theta) = -\frac{1}{2} \norm{\theta - G(z)}_F^2 - \scpr{C_0}{\theta}_F,
    \]
    where \(C_0 \coloneqq 3I\). The function \(G: \R^{1\times 2} \to \mathbb{S}_2\) is any map that is globally Lipschitz with constant \(L_G > 0\) and satisfies the structural condition \(\Tr(G(z)) = 0\) for all \(z \in \R^{1\times 2}\). For example, \(G(z_1, z_2) = z_1 \begin{psmallmatrix} 1 & 0 \\ 0 & -1 \end{psmallmatrix} + z_2 \begin{psmallmatrix} 0 & 1 \\ 1 & 0 \end{psmallmatrix}\) is a valid choice.
\end{enumerate}
The primitives satisfy the baseline conditions of Assumption \ref{ass:bsdi_primitives}. The driver \(F\) is uniformly Lipschitz in \(z\) and \(\theta\). We must now verify the central condition of Hypothesis \ref{hyp:global_regularity}(iii): the existence of a unique, globally regular maximizer map.

\vspace{1ex}
\noindent\textbf{Step 2: Identification of the Unique Maximizer Map \(\Thetacal^*\).}
The pointwise optimality condition requires finding the maximizer of \(F\). This is equivalent to finding the minimizer of \(-F\):
\begin{equation}
\label{eq:proof_argmin}
    \Thetacal^*(z) \in \underset{\theta \in \Ucal}{\argmax} \; F(z, \theta) = \underset{\theta \in \Ucal}{\argmin} \left( \frac{1}{2} \norm{\theta - G(z)}_F^2 + \scpr{C_0}{\theta}_F \right).
\end{equation}
We re-frame the objective function by completing the square. The expression to be minimized is:
\begin{align*}
    \frac{1}{2}\scpr{\theta - G(z)}{\theta - G(z)}_F + \scpr{C_0}{\theta}_F
    &= \frac{1}{2}\left( \norm{\theta}_F^2 - 2\scpr{\theta}{G(z)}_F + \norm{G(z)}_F^2 \right) + \scpr{C_0}{\theta}_F \\
    &= \frac{1}{2}\left( \norm{\theta}_F^2 - 2\scpr{\theta}{G(z) - C_0}_F \right) + \frac{1}{2}\norm{G(z)}_F^2.
\end{align*}
Minimizing this expression over \(\theta\) is equivalent to minimizing \(\norm{\theta}_F^2 - 2\scpr{\theta}{G(z) - C_0}_F\), which is in turn equivalent to minimizing \(\norm{\theta - (G(z) - C_0)}_F^2\). Let us define the target point \(P(z) \in \mathbb{S}_2\) as:
\[
    P(z) \coloneqq G(z) - C_0 = G(z) - 3I.
\]
The optimization problem \eqref{eq:proof_argmin} is thus equivalent to finding the nearest point in the set \(\Ucal\) to the target point \(P(z)\):
\[
    \Thetacal^*(z) \in \underset{\theta \in \Ucal}{\argmin} \; \norm{\theta - P(z)}_F^2 = \proj_{\Ucal}(P(z)).
\]
The set of optimizers is the nearest-point projection of \(P(z)\) onto \(\Ucal\). The key to ensuring uniqueness is to show that \(P(z)\) is never equidistant from the two disjoint components of \(\Ucal\). The projection onto \(\Ucal\) is given by the projection onto the closer of the two balls \(\Bcal(-2I,1)\) and \(\Bcal(2I,1)\).

We compare the squared distance from \(P(z)\) to the centers of the two balls, \(2I\) and \(-2I\). Consider the difference:
\begin{align*}
    \norm{P(z) - 2I}_F^2 - \norm{P(z) - (-2I)}_F^2
    &= \norm{P(z) - 2I}_F^2 - \norm{P(z) + 2I}_F^2 \\
    &= \scpr{P(z)-2I}{P(z)-2I}_F - \scpr{P(z)+2I}{P(z)+2I}_F \\
    &= \left(\norm{P(z)}_F^2 - 4\scpr{P(z)}{I}_F + 4\norm{I}_F^2\right) - \left(\norm{P(z)}_F^2 + 4\scpr{P(z)}{I}_F + 4\norm{I}_F^2\right) \\
    &= -8\scpr{P(z)}{I}_F = -8\Tr(P(z)).
\end{align*}
Now, we use the definition of \(P(z)\) and the properties of the trace and the map \(G\):
\begin{align*}
    -8\Tr(P(z)) &= -8\Tr(G(z) - 3I) \\
    &= -8(\Tr(G(z)) - \Tr(3I)) \\
    &= -8(0 - 3\Tr(I)).
\end{align*}
Since \(d=2\), \(\Tr(I) = 2\). The difference in squared distances is therefore:
\[
    -8(-3 \times 2) = 48 > 0.
\]
This strict inequality \(\norm{P(z) - 2I}_F^2 > \norm{P(z) + 2I}_F^2\) holds for all \(z \in \R^{1\times 2}\). It demonstrates that the target point \(P(z)\) is always strictly closer to the center of the first ball, \(-2I\), than to the center of the second ball, \(2I\). Consequently, the nearest point in \(\Ucal\) to \(P(z)\) must lie in \(\Bcal(-2I, 1)\). The argmax correspondence is therefore single-valued and given by:
\[
    \Thetacal^*(z) = \proj_{\Bcal(-2I, 1)}(P(z)) = \proj_{\Bcal(-2I, 1)}(G(z) - 3I).
\]

\vspace{1ex}
\noindent\textbf{Step 3: Verification of Global Lipschitz Continuity.}
We have established the existence of a unique, single-valued maximizer map \(\Thetacal^*\). It remains to show that this map is globally Lipschitz.

Let \(\Pi(\cdot) \coloneqq \proj_{\Bcal(-2I, 1)}(\cdot)\) denote the projection onto the closed convex set \(\Bcal(-2I, 1)\). A standard result in convex analysis (see, e.g., \cite[Theorem 12.19]{RockafellarWets1998}) is that the projection operator onto a non-empty, closed, convex set in a Hilbert space is non-expansive, i.e., it is \(1\)-Lipschitz.
Therefore, for any two points \(p_1, p_2 \in \mathbb{S}_2\), we have:
\[
    \norm{\Pi(p_1) - \Pi(p_2)}_F \le \norm{p_1 - p_2}_F.
\]
Now consider two states \(z_1, z_2 \in \R^{1\times 2}\). Let \(p_1 = P(z_1) = G(z_1) - 3I\) and \(p_2 = P(z_2) = G(z_2) - 3I\). We can bound the distance between the corresponding maximizers:
\begin{align*}
    \norm{\Thetacal^*(z_1) - \Thetacal^*(z_2)}_F
    &= \norm{\Pi(p_1) - \Pi(p_2)}_F \\
    &\le \norm{p_1 - p_2}_F &&\text{(by non-expansiveness of \(\Pi\))} \\
    &= \norm{(G(z_1) - 3I) - (G(z_2) - 3I)}_F \\
    &= \norm{G(z_1) - G(z_2)}_F \\
    &\le L_G \norm{z_1 - z_2}. &&\text{(since G is globally Lipschitz)}
\end{align*}
This shows that the maximizer map \(\Thetacal^*\) is globally Lipschitz with a Lipschitz constant no larger than \(L_G\). The progressive measurability of \(\Thetacal^*(t, \omega, z)\) is also satisfied, as it is a composition of a Lipschitz (hence continuous) map \(\Pi\) and a function \(G\) which is assumed to have the necessary measurability properties.

All conditions of Hypothesis \ref{hyp:global_regularity} are satisfied for a genuinely non-convex uncertainty set \(\Ucal\), which completes the proof.
\end{proof}

\begin{remark}
This example illustrates the mechanism by which global regularity can be achieved despite non-convexity. The linear term \(\scpr{C_0}{\theta}_F\) in the driver acts as a tilt, ensuring that the optimization landscape is always biased towards one of the convex components of \(\Ucal\). The problem's solution path never encounters the singular region of the state space where the choice of component would be ambiguous. This construction, while specific, proves that the class of models treatable by our first theorem is non-empty and provides a clear contrast with the geometric regularity conditions required when such a uniform tilt is absent.
\end{remark}


\section{An Example Satisfying the Geometric Regularity Conditions}
\label{app:geometric_example_corrected}

The purpose of this appendix is to construct a concrete, non-trivial example of a \(\theta\)-BSDE that is well-posed under the Geometric Regularity framework of Theorem \ref{thm:existence_geometric}, yet demonstrably fails the stronger conditions of the Global Regularity Hypothesis (\Cref{hyp:global_regularity}). This example serves to highlight the necessity and power of the Malliavin calculus argument at the heart of our main theorem, which relies on the more general Transversality Condition.

Crucially, the model is constructed such that the dimension of the driving state process is strictly less than the dimension of the parameter space, a setting where the map to the target space can still satisfy the transversality condition. This explicitly demonstrates the increased scope of our main result.

\subsection{Model Specification}
We define the primitives of our model as follows:
\begin{itemize}
    \item \textbf{Dimensions:} Let the dimension of the underlying Brownian motion be $d=2$. The parameter space is the space of $2 \times 2$ real symmetric matrices, $\Sd = \mathbb{S}_2$, which is a vector space of dimension $\frac{d(d+1)}{2} = 3$. The dimension of the forward state process is $k=2$.
    
    \item \textbf{Forward Process:} Let the state process $X_t = (X_t^1, X_t^2)^T \in \R^2$ evolve as a standard 2-dimensional Brownian motion, i.e.,
    \[ dX_t = I_2 \, dB_t, \quad X_0 = 0. \]
    Here, $B_t$ is a 2-dimensional standard Brownian motion, and the diffusion coefficient is the $2 \times 2$ identity matrix $\sigma(t,x) = I_2$. The drift is $b(t,x) = 0$.
    
    \item \textbf{Uncertainty Set $\Ucal$:} The uncertainty set $\Ucal \subset \mathbb{S}_2$ is the disjoint union of two closed balls centered at scalar multiples of the identity matrix $I \in \mathbb{S}_2$:
    \[ \Ucal \coloneqq \Bcal(-2I, 1) \cup \Bcal(2I, 1), \]
    where $\Bcal(C, r) = \{\theta \in \mathbb{S}_2 \mid \norm{\theta - C}_F \le r\}$ denotes the closed ball with respect to the Frobenius norm. This set is compact and manifestly non-convex.

    \item \textbf{Driver $F$:} The driver $F$ has the projection-type structure required by Assumption \ref{ass:pathwise_consistency}(i):
    \[ F(t,x,y,z,\theta) = -\frac{1}{2}\norm{\theta - G(x)}_F^2. \]
    Here, the function $h$ from the assumption is identically zero, $h(t,x,y,z) \equiv 0$. The map $G: \R^2 \to \mathbb{S}_2$ is defined below.

    \item \textbf{Target Map $G$:} We define the map $G$ to be a simple linear function of the forward state $x$:
    \[ G(x) \coloneqq G(x_1, x_2) = x_1 \begin{pmatrix} 1 & 0 \\ 0 & 0 \end{pmatrix} + x_2 \begin{pmatrix} 0 & 1 \\ 1 & 0 \end{pmatrix} = \begin{pmatrix} x_1 & x_2 \\ x_2 & 0 \end{pmatrix}. \]
    This map is smooth (in fact, linear) in $x$ and independent of $(t,y,z)$.
\end{itemize}

\subsection{Part 1: Verification of Assumption 4.1 (Geometric Regularity)}
We now rigorously verify that the constructed model satisfies all conditions of our main geometric regularity assumption.

\begin{proof}[Verification of Assumption 4.1]
We check each condition in turn.
\begin{enumerate}[label=(\roman*)]
    \item \textbf{Projection Driver Structure:} Satisfied by construction. The functions $b(x)=0$, $\sigma(x)=I_2$, $h(x,y,z)=0$, and $G(x)$ are all smooth ($C^\infty$) with bounded partial derivatives. Both $h$ and $G$ are globally Lipschitz.

    \item \textbf{Geometric Regularity of $\Ucal$:} The set $\Ucal$ is a finite union of disjoint compact convex sets. By the theory of sets with positive reach \cite[Theorem 4.8]{Federer1959}, each component has positive reach, and this property is preserved under a finite disjoint union. The set is also semi-algebraic. Its cut locus is therefore a semi-algebraic set of Lebesgue measure zero that admits a finite stratification, satisfying the assumption.

    \item \textbf{Uniform Ellipticity of Forward Process:} The diffusion matrix of the forward process $X_t$ is $\sigma\sigma^T = I_2 I_2^T = I_2$. For any vector $v \in \R^2$, we have $\scpr{v}{(\sigma\sigma^T)v} = v^T I_2 v = \norm{v}^2$. The uniform ellipticity condition is satisfied with $\lambda = 1$.
    
    \item \textbf{Transversality Condition:} This is the most critical condition to verify. We must show that the map $G(x)$ is transverse to every stratum of the cut locus $\mathrm{Cut}(\Ucal)$.
    
    \paragraph{Step (a): Identify the Cut Locus and its Tangent Space.}
    The cut locus is the set of points in $\mathbb{S}_2$ that are equidistant from the two components of $\Ucal$. For two balls of the same radius centered at $C_1=-2I$ and $C_2=2I$, this is the perpendicular bisecting hyperplane:
    \begin{align*}
        \norm{P - C_1}_F^2 &= \norm{P - C_2}_F^2 \\
        \norm{P + 2I}_F^2 &= \norm{P - 2I}_F^2 \\
        \norm{P}_F^2 + 4\scpr{P}{I}_F + 4\norm{I}_F^2 &= \norm{P}_F^2 - 4\scpr{P}{I}_F + 4\norm{I}_F^2 \\
        8\scpr{P}{I}_F &= 0 \implies \Tr(P) = 0.
    \end{align*}
    The cut locus is the set of all traceless $2\times2$ symmetric matrices, which we denote by $\Sigma$:
    \[ \mathrm{Cut}(\Ucal) = \Sigma \coloneqq \left\{ \begin{pmatrix} a & b \\ b & -a \end{pmatrix} \mid a,b \in \R \right\}. \]
    Since $\Sigma$ is a linear subspace of $\mathbb{S}_2$, it is a smooth submanifold (a single stratum). Its tangent space at any point $P \in \Sigma$ is the subspace itself: $T_P\Sigma = \Sigma$.
    
    \paragraph{Step (b): Identify the Image of the Derivative of $G$.}
    The map is $G(x_1, x_2) = x_1 E_1 + x_2 E_2$, where $E_1 = \begin{psmallmatrix} 1 & 0 \\ 0 & 0 \end{psmallmatrix}$ and $E_2 = \begin{psmallmatrix} 0 & 1 \\ 1 & 0 \end{psmallmatrix}$. Since $G$ is linear, its derivative at any point $x$ is the map itself, $DG_x(v) = G(v)$. The image of the derivative is the subspace spanned by $E_1$ and $E_2$:
    \[ \mathcal{V} \coloneqq \Image(DG_x) = \mathrm{span}\{E_1, E_2\} = \left\{ \begin{pmatrix} \alpha & \beta \\ \beta & 0 \end{pmatrix} \mid \alpha,\beta \in \R \right\}. \]

    
    \paragraph{Step (c): Verify Transversality.}
    The transversality condition requires that for any $x$ such that $G(x) \in \Sigma$, we have $\Image(DG_x) + T_{G(x)}\Sigma = \mathbb{S}_2$. Since both are constant subspaces, we must verify that $\mathcal{V} + \Sigma = \mathbb{S}_2$.
    The space $\Sigma$ is spanned by the basis $\left\{ B_1, B_2 \right\}$, where $B_1 = \begin{psmallmatrix} 1 & 0 \\ 0 & -1 \end{psmallmatrix}$ and $B_2 = \begin{psmallmatrix} 0 & 1 \\ 1 & 0 \end{psmallmatrix}$.
    The space $\mathcal{V}$ is spanned by the basis $\left\{ E_1, E_2 \right\}$. Note that $E_2 = B_2$.
    The sum of the subspaces, $\mathcal{V} + \Sigma$, is spanned by the union of their bases: $\{E_1, E_2, B_1, B_2\} = \{E_1, B_1, B_2\}$.
    The set of spanning matrices is:
    \[ \left\{ \begin{pmatrix} 1 & 0 \\ 0 & 0 \end{pmatrix}, \quad \begin{pmatrix} 1 & 0 \\ 0 & -1 \end{pmatrix}, \quad \begin{pmatrix} 0 & 1 \\ 1 & 0 \end{pmatrix} \right\}. \]
    These three matrices are linearly independent and live in the 3-dimensional space $\mathbb{S}_2$. They therefore form a basis for $\mathbb{S}_2$. Thus, we have shown that $\mathcal{V} + \Sigma = \mathbb{S}_2$. The map $G$ is transverse to the cut locus $\Sigma$.
\end{enumerate}
All conditions of Assumption 4.1 are met. Therefore, Theorem \ref{thm:existence_geometric} applies and guarantees that this \(\theta\)-BSDE admits a unique solution triplet $(Y,Z,\Theta)$.
\end{proof}

\subsection{Part 2: Failure of the Global Regularity Hypothesis}
We now demonstrate that this same model fails Hypothesis 3.1, proving that our main theorem is non-vacuous and covers cases inaccessible to the simpler theory.

\begin{proof}[Demonstration of Failure of Hypothesis 3.1]
The core of Hypothesis 3.1 is the requirement that the argmax correspondence is a single-valued, globally Lipschitz function. We show that it is not single-valued on a set of states that can be visited by the driving process with positive probability.

\paragraph{Step 1: Relate the Maximizer to a Geometric Projection.}
The pointwise optimality condition is
\[ \Theta_t \in \underset{\theta \in \Ucal}{\argmax} \; F(t,X_t,y,z,\theta) = \underset{\theta \in \Ucal}{\argmax} \; \left(-\frac{1}{2}\norm{\theta - G(X_t)}_F^2\right). \]
Maximizing this expression is equivalent to minimizing the squared Frobenius distance. The maximizer is therefore the nearest-point projection of the target point $W_t \coloneqq G(X_t)$ onto the set $\Ucal$:
\[ \Theta_t \in \proj_{\Ucal}(W_t). \]

\paragraph{Step 2: Identify the Singular Set for the Projection.}
As established in the verification of transversality, the projection $\proj_{\Ucal}(P)$ is non-unique if and only if the point $P$ lies on the cut locus, $\mathrm{Cut}(\Ucal) = \Sigma = \{ P \in \mathbb{S}_2 \mid \Tr(P) = 0 \}$.

\paragraph{Step 3: Show the Singular Set is Visited by the Target Process.}
The maximizer is non-unique whenever the target process $W_t$ lands on this cut locus. The target process is $W_t = G(X_t) = \begin{pmatrix} X_t^1 & X_t^2 \\ X_t^2 & 0 \end{pmatrix}$. The condition for non-uniqueness at time $t$ is:
\[ \Tr(W_t) = 0 \iff \Tr\begin{pmatrix} X_t^1 & X_t^2 \\ X_t^2 & 0 \end{pmatrix} = 0 \iff X_t^1 = 0. \]
The process $X_t^1$ is a standard one-dimensional Brownian motion starting at zero. By the properties of Brownian motion, for any $\epsilon > 0$, the probability that the path of $X^1$ hits zero during the time interval $(0, \epsilon]$ is 1. That is,
\[ P_0(\exists t \in (0, \epsilon] \text{ s.t. } X_t^1(\omega) = 0) = 1. \]
This demonstrates that the set of states $(t,x,y,z)$ for which the argmax correspondence is multi-valued is not only non-empty, but is actively and almost surely visited by the driving processes of the system. This implies that the maximizer map $\Thetacal^*$ cannot be represented as a globally single-valued function, let alone a globally Lipschitz one. Hypothesis 3.1 fundamentally fails for this model.
\end{proof}

\subsection{Conclusion}
This example successfully demonstrates the scope and power of our main result. We have constructed a model with a non-convex uncertainty set where the mapping from the system's state to the optimal parameter choice has unavoidable geometric singularities. Since the dimension of the state space ($k=2$) is less than the dimension of the parameter space ($\dim(\mathbb{S}_2)=3$), the Global Regularity Hypothesis is clearly violated.

Nevertheless, because the map $G$ is transverse to the singular stratum of the uncertainty set, the model satisfies the geometric and non-degeneracy conditions of Assumption 4.1. Theorem 4.2 therefore applies, guaranteeing the existence of a unique solution. The theorem implicitly relies on the Malliavin calculus argument, which proves that any solution path is such that the target process $G(X_s)$ almost surely avoids the null set $\mathrm{Cut}(\Ucal)$, thus ensuring pathwise uniqueness of the maximizer and restoring well-posedness. This example confirms that our main result provides a genuine and significant extension of the theory to a broader, more realistic class of non-convex problems.
\end{appendices}

\section{Operational Significance: Dynamic Hedging and Ambiguity-Aware Risk Management}
\label{sec:operational_significance}

A central result of this paper, established in Theorem \ref{thm:indistinguishability_linear_revised}, is that the value process $Y_t$ derived from our $\theta$-BSDE framework is identical to the value obtained from the classical G-calculus, which depends only on the convex hull of the uncertainty set, $\overlineconv(\Ucal)$. This gives rise to a critical question: if the primary valuation output is unchanged, what is the operational significance of the vastly more complex machinery developed in this paper?

This section provides a definitive answer to that question. We argue that the innovation of our framework lies not in altering the initial price of a claim, but in providing a richer, dynamic information set that allows for superior, ambiguity-aware risk management. We demonstrate that the unique maximizer process $\Theta_t$, which is invisible to the G-calculus, is an indispensable tool for optimal decision-making. We formalize this by constructing a dynamic risk management problem where an agent equipped with our framework (the `$\theta$-Agent') demonstrably outperforms an agent using the classical framework (the `G-Agent').

\subsection{The Setup: Two Agents and a Hedging Problem}
Consider an agent who sells a European contingent claim with payoff $\xi = \phi(X_T)$ at time $t=0$. The pricing and hedging environment is characterized by a non-convex, compact uncertainty set $\Ucal \subset \Sdp$ governing volatility, as in Section \ref{sec:theta_system_and_sensitivity}. The canonical driver is $F(z, \theta) = \frac{1}{2}\Tr(z\theta z^T)$. We define two competing risk-management agents:

\begin{definition}[The G-Agent vs. The $\theta$-Agent]
\begin{enumerate}
    \item \textbf{The G-Agent:} This agent operates within the classical sub-additive framework. Their model of ambiguity is the convexified set $\Ucal_G \coloneqq \overlineconv(\Ucal)$. They compute the price of the claim $\xi$ as $Y_0^{(G)}$, the solution to the G-BSDE, and implement the corresponding delta-hedging strategy using the volatility process $Z_t^{(G)}$.
    
    \item \textbf{The $\theta$-Agent:} This agent uses the framework developed in this paper, with the primitive non-convex set $\Ucal$. They compute the price $Y_0^{(\theta)}$ and the hedging process $Z_t^{(\theta)}$ by solving the $\theta$-BSDE.
\end{enumerate}
\end{definition}

By Theorem \ref{thm:indistinguishability_linear_revised}, we know that the initial prices and hedging strategies are identical: $Y_0^{(G)} = Y_0^{(\theta)}$ and the processes $Z^{(G)}$ and $Z^{(\theta)}$ are indistinguishable. The crucial difference lies in their information sets. The G-Agent only observes the history of the Brownian motion, $\mathcal{F}_t^B$. The $\theta$-Agent, by solving the full $\theta$-BSDE system, also observes the unique maximizer process $\Theta_t \equiv \Theta_t^{(\theta)}$, which takes values in the primitive set $\Ucal$. The information filtration for the $\theta$-Agent is $\mathcal{F}_t^\Theta \coloneqq \mathcal{F}_t^B \vee \sigma(\{\Theta_s : s \le t\})$.

\subsection{A Dynamic Risk Management Problem: Optimal Hedging Termination}
Suppose the agent, having sold the claim and initiated the hedge, has the right to terminate the contract at any time $t \in [0,T)$. If they terminate at time $t$, they pay a fixed administrative and market-impact cost $K > 0$ and unwind their hedge. If they do not terminate, they must carry the hedge to maturity and pay the terminal payoff $\xi$, incurring a terminal hedging error. The agent's goal is to choose a stopping time $\tau$ to minimize a total expected cost.

\begin{definition}[Hedging Error and Cost Functional]
Let $(Y_t, Z_t)$ be the unique solution process common to both agents. The terminal hedging error is $e_T \coloneqq Y_0 + \int_0^T Z_s dB_s - \phi(X_T)$. Note that from the BSDE definition, $e_T = \int_0^T F(Z_s, \Theta_s^{\text{true}}) ds$, where $\Theta_s^{\text{true}}$ is the (unknown) realized path of the true volatility parameter. The agent's objective is to solve the optimal stopping problem:
\begin{equation}
\label{eq:optimal_stopping_problem}
    V_t = \inf_{\tau \in \mathcal{T}_{[t,T]}} \mathbb{E}\left[ K \cdot \mathbf{1}_{\tau < T} + L(e_T) \cdot \mathbf{1}_{\tau=T} \mid \mathcal{I}_t \right],
\end{equation}
where $\mathcal{T}_{[t,T]}$ is the set of stopping times with values in $[t,T]$, $L: \R \to \R^+$ is a convex loss function (e.g., $L(x)=x^2$), and $\mathcal{I}_t$ is the agent's information set at time $t$.
\end{definition}

The key distinction is that for the G-Agent, $\mathcal{I}_t = \mathcal{F}_t^B$, while for the $\theta$-Agent, $\mathcal{I}_t = \mathcal{F}_t^\Theta$.

\subsection{Superiority of the Ambiguity-Aware Strategy}
We now state and prove the main result of this section: the additional information contained in $\Theta_t$ leads to a strictly better optimal strategy.

\begin{proposition}[Strict Superiority of the $\theta$-Agent]
\label{prop:agent_superiority}
Let $\Ucal$ be a non-convex, compact set satisfying the geometric regularity conditions of Theorem \ref{thm:existence_geometric}. Let $\Ucal_G = \overlineconv(\Ucal)$ be its convex hull. Assume there exists a state $z^*$ and a time $t_0$ such that the unique maximizer of $F(z^*, \cdot)$ over $\Ucal_G$ lies in the interior of $\Ucal_G \setminus \Ucal$.

Let $V_t^{(G)}$ and $V_t^{(\theta)}$ be the value functions of the optimal stopping problem \eqref{eq:optimal_stopping_problem} for the G-Agent and the $\theta$-Agent, respectively. Then:
\begin{enumerate}[label=(\roman*)]
    \item For all $t \in [0,T]$, $V_t^{(\theta)} \le V_t^{(G)}$ almost surely.
    \item There exists a set of paths of positive probability on which the optimal stopping strategy for the $\theta$-Agent, $\tau^*_{(\theta)}$, is strictly different from that of the G-Agent, $\tau^*_{(G)}$, and $V_t^{(\theta)} < V_t^{(G)}$.
\end{enumerate}
The information provided by the process $\Theta_t$ is operationally significant.
\end{proposition}

\begin{proof}
\noindent\textbf{Part (i): Dominance of the $\theta$-Strategy.}
The information filtration of the $\theta$-Agent is larger than that of the G-Agent: $\mathcal{F}_t^B \subseteq \mathcal{F}_t^\Theta$. The value function of an optimal stopping problem is always monotonically non-increasing with respect to the enlargement of the underlying filtration. This is a classic result; for any stopping time $\tau$ adapted to the smaller filtration $\mathcal{F}_t^B$, it is also adapted to $\mathcal{F}_t^\Theta$. Thus, the $\theta$-Agent is optimizing over a larger set of strategies. By the tower property of conditional expectation, for any $\mathcal{F}_t^B$-measurable random variable $X$, $\mathbb{E}[X|\mathcal{F}_t^B] = \mathbb{E}[\mathbb{E}[X|\mathcal{F}_t^\Theta]|\mathcal{F}_t^B]$. A standard argument for optimal stopping problems then yields $V_t^{(\theta)} \le V_t^{(G)}$.

\vspace{1ex}
\noindent\textbf{Part (ii): Strict Improvement and Strategy Dependence.}
The core of the proof is to show that the conditional expectation of the terminal loss, which defines the continuation value of the problem, is different for the two agents. Let the continuation value at time $t$ be $C_t \coloneqq \mathbb{E}[L(e_T) \mid \mathcal{I}_t]$.

The terminal error is $e_T = -\int_0^T \tilde{F}(Z_s)ds$, where $\tilde{F}(z) = \frac{1}{2}\sup_{\theta' \in \Ucal}\Tr(z\theta'z^T)$. We can rewrite this as $e_T = -\frac{1}{2}\int_0^T \Tr(Z_s \Theta_s^{\text{true}} Z_s^T)ds$. The key unknown is the path of $\Theta_s^{\text{true}}$ for $s>t$.

Let's analyze the continuation values for the two agents.
\begin{itemize}
    \item \textbf{G-Agent's Continuation Value:} $C_t^{(G)} = \mathbb{E}\left[L\left(-\frac{1}{2}\int_0^T \Tr(Z_s \Theta_s^{\text{true}} Z_s^T)ds\right) \mid \mathcal{F}_t^B\right]$. The G-Agent's model of the world is that $\Theta_s^{\text{true}}$ is selected adversarially from $\Ucal_G = \overlineconv(\Ucal)$. Their conditional expectation is an average over all possible scenarios consistent with this assumption.
    
    \item \textbf{$\theta$-Agent's Continuation Value:} $C_t^{(\theta)} = \mathbb{E}\left[L\left(-\frac{1}{2}\int_0^T \Tr(Z_s \Theta_s^{\text{true}} Z_s^T)ds\right) \mid \mathcal{F}_t^\Theta\right]$. The $\theta$-Agent's model is that $\Theta_s^{\text{true}}$ is selected from $\Ucal$. Crucially, they have observed the path $\{\Theta_r\}_{r\le t}$.
\end{itemize}

By assumption, there exists a state $z^*$ such that $\theta_G^* \coloneqq \argmax_{\theta \in \Ucal_G} F(z^*, \theta)$ lies in $\Ucal_G \setminus \Ucal$. Let $\theta_\theta^* \coloneqq \argmax_{\theta \in \Ucal} F(z^*, \theta)$. By construction, $F(z^*, \theta_G^*) > F(z^*, \theta_\theta^*)$.

Under the non-degeneracy conditions of our main theorem, the process $Z_t$ has a law that is absolutely continuous with respect to the Lebesgue measure. Therefore, with positive probability, the process $Z_t$ will enter a small neighborhood of $z^*$ at some time $t$. Let's consider such a path $\omega$ where $Z_t(\omega) \approx z^*$.
On this path:
\begin{itemize}
    \item The $\theta$-Agent observes $\Theta_t(\omega) = \argmax_{\theta \in \Ucal} F(Z_t, \theta) \approx \theta_\theta^*$.
    \item The G-Agent does not observe $\Theta_t$. Their model is driven by the theoretical maximizer $\Theta_t^{(G)} = \argmax_{\theta \in \Ucal_G} F(Z_t, \theta) \approx \theta_G^*$.
\end{itemize}

The $\theta$-Agent's observation of $\Theta_t \approx \theta_\theta^*$ provides powerful information. They know that the currently active adversarial model is $\theta_\theta^*$, a specific point within the primitive set $\Ucal$. Their prediction for the future evolution of the terminal loss $e_T$ will be conditioned on this fact.

The G-Agent, lacking this observation, must average over all possibilities. Their conditional expectation for the terminal loss is based on the premise that the future adversarial volatility could be $\theta_G^*$, an "averaged" model that is physically impossible under the true model specification.

Let us formalize the difference. The conditional laws of the future path $\{\Theta_s^{\text{true}}\}_{s>t}$ are different:
\[ \mathrm{Law}(\{\Theta_s^{\text{true}}\}_{s>t} \mid \mathcal{F}_t^B) \neq \mathrm{Law}(\{\Theta_s^{\text{true}}\}_{s>t} \mid \mathcal{F}_t^\Theta). \]
The second law is more concentrated because it incorporates the knowledge that $\Theta_t$ is currently in a specific region of $\Ucal$. This implies that the conditional expectations of the terminal loss will be different: $C_t^{(G)} \neq C_t^{(\theta)}$ on this set of paths.

Since the continuation values are different, the optimal decisions will be different. The decision to stop at time $t$ is made if the immediate cost of stopping, $K$, is less than the expected cost of continuing, $C_t$. Since $C_t^{(G)} \neq C_t^{(\theta)}$, there will be paths where, for instance,
\[ C_t^{(\theta)} > K \ge C_t^{(G)}. \]
On such a path, the $\theta$-Agent's optimal decision is to stop immediately ($\tau^*_{(\theta)} = t$), while the G-Agent continues. This happens if the observation of $\Theta_t$ indicates that the system has entered a high-risk regime (a region of $\Ucal$ corresponding to high variance), causing the conditional expected loss to spike above the exit cost $K$. The G-Agent, whose conditional expectation averages over high- and low-risk regimes, would underestimate this risk and fail to make the optimal decision.

Because such paths occur with positive probability, the inequality $V_t^{(\theta)} < V_t^{(G)}$ holds on a set of non-zero measure. This proves that the additional information in $\Theta_t$ has strict, positive economic value and is essential for optimal risk management.
\end{proof}

\begin{remark}[The Operational Value of $\Theta_t$]
This proposition provides a rigorous demonstration of the operational value of our framework. While the G-calculus provides a robust price, it is blind to the dynamic evolution of the adversarial model within the ambiguity set. The $\theta$-calculus, by uniquely identifying the active adversary $\Theta_t$, provides a real-time signal of model risk. This signal is crucial for any dynamic decision problem—such as optimal exit, capital adjustment, or secondary hedging—that depends on a forward-looking assessment of risk. An agent who ignores this signal is sub-optimal and will, on average, incur greater costs than an agent who uses it.
\end{remark}



\begin{thebibliography}{99}

\bibitem{Aliprantis2006}
C. D. Aliprantis and K. C. Border,
\textit{Infinite Dimensional Analysis: A Hitchhiker's Guide}, Third Edition,
Springer, 2006.

\bibitem{Antonelli2002}
F. Antonelli and A. Kohatsu-Higa,
\textit{Densities of solutions to backward SDEs},
Potential Analysis, 16(1), pp. 1--33, 2002.

\bibitem{Artzner1999}
P. Artzner, F. Delbaen, J.-M. Eber, and D. Heath,
\textit{Coherent measures of risk},
Mathematical Finance, 9(3), pp. 203--228, 1999.

\bibitem{CastaingValadier1977}
C. Castaing and M. Valadier,
\textit{Convex Analysis and Measurable Multifunctions},
Lecture Notes in Mathematics, vol. 580, Springer, 1977.

\bibitem{CrandallIshiiLions1992}
M. G. Crandall, H. Ishii, and P.-L. Lions,
\textit{User's guide to viscosity solutions of second order partial differential equations},
Bulletin of the American Mathematical Society, 27(1), pp. 1--67, 1992.

\bibitem{ElKaroui1997}
N. El Karoui, S. Peng, and M. C. Quenez,
\textit{Backward stochastic differential equations in finance},
Mathematical Finance, 7(1), pp. 1--71, 1997.

\bibitem{Federer1959}
H. Federer,
\textit{Curvature measures},
Transactions of the American Mathematical Society, 93(3), pp. 418--491, 1959.

\bibitem{Federer1969}
H. Federer,
\textit{Geometric Measure Theory},
Springer-Verlag, 1969.

\bibitem{GuilleminPollack1974}
V. Guillemin and A. Pollack,
\textit{Differential Topology},
Prentice-Hall, 1974.

\bibitem{Knight1921}
F. H. Knight,
\textit{Risk, Uncertainty and Profit},
Houghton Mifflin, 1921.

\bibitem{Kobylanski2000}
M. Kobylanski,
\textit{Backward stochastic differential equations and partial differential equations with quadratic growth},
The Annals of Probability, 28(2), pp. 558--602, 2000.

\bibitem{MaYong1999}
J. Ma and J. Yong,
\textit{Forward-Backward Stochastic Differential Equations and their Applications},
Springer, 1999.

\bibitem{Nualart2006}
D. Nualart,
\textit{The Malliavin Calculus and Related Topics}, Second Edition,
Springer, 2006.

\bibitem{PardouxPeng1990}
E. Pardoux and S. Peng,
\textit{Adapted solution of a backward stochastic differential equation},
Systems \& Control Letters, 14(1), pp. 55--61, 1990.

\bibitem{PardouxPeng1992}
E. Pardoux and S. Peng,
\textit{Backward stochastic differential equations and quasilinear parabolic partial differential equations},
in Stochastic Partial Differential Equations and Their Applications, B.L. Rozovskii and R.B. Sowers, eds., Springer, pp. 200--217, 1992.

\bibitem{Peng2019}
S. Peng,
\textit{Nonlinear Expectations and Stochastic Calculus under Ambiguity},
Springer, 2019.

\bibitem{RevuzYor1999}
D. Revuz and M. Yor,
\textit{Continuous Martingales and Brownian Motion}, Third Edition,
Grundlehren der mathematischen Wissenschaften, vol. 293, Springer-Verlag, Berlin, 1999.

\bibitem{Rockafellar1970}
R. T. Rockafellar,
\textit{Convex Analysis},
Princeton University Press, 1970.

\bibitem{RockafellarWets1998}
R. T. Rockafellar and R. J-B. Wets,
\textit{Variational Analysis},
Grundlehren der mathematischen Wissenschaften, vol. 317, Springer-Verlag, Berlin, 1998.

\bibitem{Evans2010}
L. C. Evans,
\textit{Partial Differential Equations}, Second Edition,
Graduate Studies in Mathematics, vol. 19, American Mathematical Society, 2010.

\end{thebibliography}



\end{document}
